
\Data\ID{1003123401}
\Updated{2020-07-18 01:07:42}
\Level{A}
\SubArea{Powers and Roots of Complex Numbers}
\LANGen{
\Question{
Given  the complex number \( a =\sqrt3\cdot\left( \cos 225^{\circ} + \mathrm{i}\cdot\sin 225^{\circ}\right) \), find  the polar form of \( a^6 \).}
{\( 27\cdot\left(\cos270^{\circ}+\mathrm{i}\cdot\sin270^{\circ}\right) \)}{\( 9\cdot\left(\cos90^{\circ}+\mathrm{i}\cdot\sin90^{\circ}\right) \)}{\( 27\cdot\left(\cos90^{\circ}+\mathrm{i}\cdot\sin90^{\circ}\right) \)}{\( 9\cdot\left(\cos270^{\circ}+\mathrm{i}\cdot\sin270^{\circ}\right) \)}{}{}}
\LANGpl{
\Question{
Dana jest liczba zespolona \( a =\sqrt3\cdot\left( \cos 225^{\circ} + \mathrm{i}\cdot\sin 225^{\circ}\right) \), wskaż postać trygonometryczną \( a^6 \).}
{\( 27\cdot\left(\cos270^{\circ}+\mathrm{i}\cdot\sin270^{\circ}\right) \)}{\( 9\cdot\left(\cos90^{\circ}+\mathrm{i}\cdot\sin90^{\circ}\right) \)}{\( 27\cdot\left(\cos90^{\circ}+\mathrm{i}\cdot\sin90^{\circ}\right) \)}{\( 9\cdot\left(\cos270^{\circ}+\mathrm{i}\cdot\sin270^{\circ}\right) \)}{}{}}
\LANGcs{
\Question{
Je dané komplexní číslo  \( a =\sqrt3\cdot\left( \cos 225^{\circ} + \mathrm{i}\cdot\sin 225^{\circ}\right) \). Určete goniometrický tvar komplexního čísla \( a^6 \).}
{\( 27\cdot\left(\cos270^{\circ}+\mathrm{i}\cdot\sin270^{\circ}\right) \)}{\( 9\cdot\left(\cos90^{\circ}+\mathrm{i}\cdot\sin90^{\circ}\right) \)}{\( 27\cdot\left(\cos90^{\circ}+\mathrm{i}\cdot\sin90^{\circ}\right) \)}{\( 9\cdot\left(\cos270^{\circ}+\mathrm{i}\cdot\sin270^{\circ}\right) \)}{}{}}
\LANGsk{
\Question{
Dané je komplexné číslo  \( a =\sqrt3\cdot\left( \cos 225^{\circ} + \mathrm{i}\cdot\sin 225^{\circ}\right) \). Určte goniometrický tvar komplexného čísla \( a^6 \).}
{\( 27\cdot\left(\cos270^{\circ}+\mathrm{i}\cdot\sin270^{\circ}\right) \)}{\( 9\cdot\left(\cos90^{\circ}+\mathrm{i}\cdot\sin90^{\circ}\right) \)}{\( 27\cdot\left(\cos90^{\circ}+\mathrm{i}\cdot\sin90^{\circ}\right) \)}{\( 9\cdot\left(\cos270^{\circ}+\mathrm{i}\cdot\sin270^{\circ}\right) \)}{}{}}
\LANGes{
\Question{
Dado el número complejo \( a =\sqrt3\cdot\left( \cos 225^{\circ} + \mathrm{i}\cdot\sin 225^{\circ}\right) \), determina la forma polar de \( a^6 \).}
{\( 27\cdot\left(\cos270^{\circ}+\mathrm{i}\cdot\sin270^{\circ}\right) \)}{\( 9\cdot\left(\cos90^{\circ}+\mathrm{i}\cdot\sin90^{\circ}\right) \)}{\( 27\cdot\left(\cos90^{\circ}+\mathrm{i}\cdot\sin90^{\circ}\right) \)}{\( 9\cdot\left(\cos270^{\circ}+\mathrm{i}\cdot\sin270^{\circ}\right) \)}{}{}}
\End

\Data\ID{1003123402}
\Updated{2020-07-18 01:07:56}
\Level{A}
\SubArea{Powers and Roots of Complex Numbers}
\LANGen{
\Question{
Given the complex number \( b=\sqrt[3]2\cdot\left(\cos\frac56\pi+\mathrm{i}\cdot\sin\frac56\pi\right) \), find the polar form of \( b^9 \).}
{\( 8\cdot\left(\cos\frac32\pi+\mathrm{i}\cdot\sin\frac32\pi\right)  \)}{\( 64\cdot\left(\cos\frac12\pi-\mathrm{i}\cdot\sin\frac12\pi\right)  \)}{\( 8\cdot\left(\cos\frac12\pi-\mathrm{i}\cdot\sin\frac12\pi\right)  \)}{\( 64\cdot\left(\cos\frac32\pi+\mathrm{i}\cdot\sin\frac32\pi\right)  \)}{}{}}
\LANGpl{
\Question{
Dana jest liczba zespolona \( b=\sqrt[3]2\cdot\left(\cos\frac56\pi+\mathrm{i}\cdot\sin\frac56\pi\right) \), wskaż postać trygonometryczną \( b^9 \).}
{\( 8\cdot\left(\cos\frac32\pi+\mathrm{i}\cdot\sin\frac32\pi\right)  \)}{\( 64\cdot\left(\cos\frac12\pi-\mathrm{i}\cdot\sin\frac12\pi\right)  \)}{\( 8\cdot\left(\cos\frac12\pi-\mathrm{i}\cdot\sin\frac12\pi\right)  \)}{\( 64\cdot\left(\cos\frac32\pi+\mathrm{i}\cdot\sin\frac32\pi\right)  \)}{}{}}
\LANGcs{
\Question{
Je dané komplexní číslo \( b=\sqrt[3]2\cdot\left(\cos\frac56\pi+\mathrm{i}\cdot\sin\frac56\pi\right) \). Určete goniometrický tvar komplexního čísla \( b^9 \).}
{\( 8\cdot\left(\cos\frac32\pi+\mathrm{i}\cdot\sin\frac32\pi\right)  \)}{\( 64\cdot\left(\cos\frac12\pi-\mathrm{i}\cdot\sin\frac12\pi\right)  \)}{\( 8\cdot\left(\cos\frac12\pi-\mathrm{i}\cdot\sin\frac12\pi\right)  \)}{\( 64\cdot\left(\cos\frac32\pi+\mathrm{i}\cdot\sin\frac32\pi\right)  \)}{}{}}
\LANGsk{
\Question{
Dané je komplexné číslo \( b=\sqrt[3]2\cdot\left(\cos\frac56\pi+\mathrm{i}\cdot\sin\frac56\pi\right) \). Určte goniometrický tvar komplexného čísla \( b^9 \).}
{\( 8\cdot\left(\cos\frac32\pi+\mathrm{i}\cdot\sin\frac32\pi\right)  \)}{\( 64\cdot\left(\cos\frac12\pi-\mathrm{i}\cdot\sin\frac12\pi\right)  \)}{\( 8\cdot\left(\cos\frac12\pi-\mathrm{i}\cdot\sin\frac12\pi\right)  \)}{\( 64\cdot\left(\cos\frac32\pi+\mathrm{i}\cdot\sin\frac32\pi\right)  \)}{}{}}
\LANGes{
\Question{
Dado el número complejo \( b=\sqrt[3]2\cdot\left(\cos\frac56\pi+\mathrm{i}\cdot\sin\frac56\pi\right) \), determina la forma polar de \( b^9 \).}
{\( 8\cdot\left(\cos\frac32\pi+\mathrm{i}\cdot\sin\frac32\pi\right)  \)}{\( 64\cdot\left(\cos\frac12\pi-\mathrm{i}\cdot\sin\frac12\pi\right)  \)}{\( 8\cdot\left(\cos\frac12\pi-\mathrm{i}\cdot\sin\frac12\pi\right)  \)}{\( 64\cdot\left(\cos\frac32\pi+\mathrm{i}\cdot\sin\frac32\pi\right)  \)}{}{}}
\End

\Data\ID{2000002101}
\Updated{2021-12-01 04:12:32}
\Level{A}
\SubArea{Powers and Roots of Complex Numbers}
\LANGen{
\Question{
Find the algebraic form of the complex number \( \left((\sqrt[3]{4}(\cos{\frac{\pi}{6}}+i\sin{\frac{\pi}{6}})\right)^3 \).}
{\( 4i\)}{\( -4i \)}{\( 4 \)}{\( 4\left(\frac{\sqrt{3}}{2} + \frac{1}{2}i\right)\)}{}{}}
\LANGpl{
\Question{
Znajdź postać algebraiczną liczby zespolonej \( \left((\sqrt[3]{4}(\cos{\frac{\pi}{6}}+i\sin{\frac{\pi}{6}})\right)^3 \).}
{\( 4i\)}{\( -4i \)}{\( 4 \)}{\( 4\left(\frac{\sqrt{3}}{2} + \frac{1}{2}i\right)\)}{}{}}
\LANGcs{
\Question{
Jak zapíšete algebraický tvar komplexního čísla \( \left((\sqrt[3]{4}(\cos{\frac{\pi}{6}}+i\sin{\frac{\pi}{6}})\right)^3 \)?}
{\( 4i\)}{\( -4i \)}{\( 4 \)}{\( 4\left(\frac{\sqrt{3}}{2} + \frac{1}{2}i\right)\)}{}{}}
\LANGsk{
\Question{
Ako zapíšete algebraický tvar komplexného čísla \( \left((\sqrt[3]{4}(\cos{\frac{\pi}{6}}+i\sin{\frac{\pi}{6}})\right)^3 \).}
{\( 4i\)}{\( -4i \)}{\( 4 \)}{\( 4\left(\frac{\sqrt{3}}{2} + \frac{1}{2}i\right)\)}{}{}}
\LANGes{
\Question{
Determina la forma algebraica del número complejo \( \left((\sqrt[3]{4}(\cos{\frac{\pi}{6}}+i\sin{\frac{\pi}{6}})\right)^3 \).}
{\( 4i\)}{\( -4i \)}{\( 4 \)}{\( 4\left(\frac{\sqrt{3}}{2} + \frac{1}{2}i\right)\)}{}{}}
\End

\Data\ID{2000002102}
\Updated{2021-12-01 04:12:38}
\Level{A}
\SubArea{Powers and Roots of Complex Numbers}
\LANGen{
\Question{
Consider \( z= \cos{\frac{\pi}{4}} + i\sin{\frac{\pi}{4}} \) and find \(z^9\).}
{\( \cos{\frac{\pi}{4}} + i\sin{\frac{\pi}{4}} \)}{\(  9 \left(\cos{\frac{\pi}{4}} + i\sin{\frac{\pi}{4}}\right) \)}{\( \cos{\frac{9\pi}{4}} - i\sin{\frac{9\pi}{4}} \)}{\( 9\left(\frac{\sqrt{2}}{2} + i\frac{\sqrt{2}}{2} \right) \)}{}{}}
\LANGpl{
\Question{
Rozważ  \( z= \cos{\frac{\pi}{4}} + i\sin{\frac{\pi}{4}} \) i znajdź \(z^9\).}
{\( \cos{\frac{\pi}{4}} + i\sin{\frac{\pi}{4}} \)}{\(  9 \left(\cos{\frac{\pi}{4}} + i\sin{\frac{\pi}{4}}\right) \)}{\( \cos{\frac{9\pi}{4}} - i\sin{\frac{9\pi}{4}} \)}{\( 9\left(\frac{\sqrt{2}}{2} + i\frac{\sqrt{2}}{2} \right) \)}{}{}}
\LANGcs{
\Question{
Určete \(z^9\), pokud platí, že \( z= \cos{\frac{\pi}{4}} + i\sin{\frac{\pi}{4}} \).}
{\( \cos{\frac{\pi}{4}} + i\sin{\frac{\pi}{4}} \)}{\(  9 \left(\cos{\frac{\pi}{4}} + i\sin{\frac{\pi}{4}}\right) \)}{\( \cos{\frac{9\pi}{4}} - i\sin{\frac{9\pi}{4}} \)}{\( 9\left(\frac{\sqrt{2}}{2} + i\frac{\sqrt{2}}{2} \right) \)}{}{}}
\LANGsk{
\Question{
Ak platí, že \( z= \cos{\frac{\pi}{4}} + i\sin{\frac{\pi}{4}} \) určte \(z^9\).}
{\( \cos{\frac{\pi}{4}} + i\sin{\frac{\pi}{4}} \)}{\(  9 \left(\cos{\frac{\pi}{4}} + i\sin{\frac{\pi}{4}}\right) \)}{\( \cos{\frac{9\pi}{4}} - i\sin{\frac{9\pi}{4}} \)}{\( 9\left(\frac{\sqrt{2}}{2} + i\frac{\sqrt{2}}{2} \right) \)}{}{}}
\LANGes{
\Question{
Considera \( z= \cos{\frac{\pi}{4}} + i\sin{\frac{\pi}{4}} \) y calcula \(z^9\).}
{\( \cos{\frac{\pi}{4}} + i\sin{\frac{\pi}{4}} \)}{\(  9 \left(\cos{\frac{\pi}{4}} + i\sin{\frac{\pi}{4}}\right) \)}{\( \cos{\frac{9\pi}{4}} - i\sin{\frac{9\pi}{4}} \)}{\( 9\left(\frac{\sqrt{2}}{2} + i\frac{\sqrt{2}}{2} \right) \)}{}{}}
\End

\Data\ID{2000002103}
\Updated{2021-12-01 04:12:43}
\Level{A}
\SubArea{Powers and Roots of Complex Numbers}
\LANGen{
\Question{
Find the algebraic form of the complex number \( \left(\sqrt{2}\left(\cos{\frac{\pi}{2}}+i\sin{\frac{\pi}{2}}\right)\right)^{18} \).}
{\( -512 \)}{\( 512i \)}{\( 512 \)}{\( -512i \)}{}{}}
\LANGpl{
\Question{
Znajdź postać algebraiczną liczby zespolonej \( \left(\sqrt{2}\left(\cos{\frac{\pi}{2}}+i\sin{\frac{\pi}{2}}\right)\right)^{18} \).}
{\( -512 \)}{\( 512i \)}{\( 512 \)}{\( -512i \)}{}{}}
\LANGcs{
\Question{
Jak zapíšete algebraický tvar komplexního čísla \( \left(\sqrt{2}\left(\cos{\frac{\pi}{2}}+i\sin{\frac{\pi}{2}}\right)\right)^{18} \)?}
{\( -512 \)}{\( 512i \)}{\( 512 \)}{\( -512i \)}{}{}}
\LANGsk{
\Question{
Ako zapíšete algebraický tvar komplexného čísla \( \left(\sqrt{2}\left(\cos{\frac{\pi}{2}}+i\sin{\frac{\pi}{2}}\right)\right)^{18} \)?}
{\( -512 \)}{\( 512i \)}{\( 512 \)}{\( -512i \)}{}{}}
\LANGes{
\Question{
Determina la forma algebraica del número complejo \( \left(\sqrt{2}\left(\cos{\frac{\pi}{2}}+i\sin{\frac{\pi}{2}}\right)\right)^{18} \).}
{\( -512 \)}{\( 512i \)}{\( 512 \)}{\( -512i \)}{}{}}
\End

\Data\ID{2000002104}
\Updated{2021-12-01 04:12:48}
\Level{A}
\SubArea{Powers and Roots of Complex Numbers}
\LANGen{
\Question{
Find the algebraic form of the complex number \(\left(\cos{\left(-\frac{\pi}{3}\right)} + i\sin{\left(-\frac{\pi}{3}\right)} \right)^{99} \).}
{\( -1 \)}{\( -i\)}{\(1\)}{\(i\)}{}{}}
\LANGpl{
\Question{
Znajdź postać algebraiczną liczby zespolonej \(\left(\cos{\left(-\frac{\pi}{3}\right)} + i\sin{\left(-\frac{\pi}{3}\right)} \right)^{99} \).}
{\( -1 \)}{\( -i\)}{\(1\)}{\(i\)}{}{}}
\LANGcs{
\Question{
Jak zapíšete algebraický tvar komplexního čísla \(\left(\cos{\left(-\frac{\pi}{3}\right)} + i\sin{\left(-\frac{\pi}{3}\right)} \right)^{99} \)?}
{\( -1 \)}{\( -i\)}{\(1\)}{\(i\)}{}{}}
\LANGsk{
\Question{
Ako zapíšete algebraický tvar komplexného čísla  \(\left(\cos{\left(-\frac{\pi}{3}\right)} + i\sin{\left(-\frac{\pi}{3}\right)} \right)^{99} \)?}
{\( -1 \)}{\( -i\)}{\(1\)}{\(i\)}{}{}}
\LANGes{
\Question{
Determina la forma algebraica del número complejo \(\left(\cos{\left(-\frac{\pi}{3}\right)} + i\sin{\left(-\frac{\pi}{3}\right)} \right)^{99} \).}
{\( -1 \)}{\( -i\)}{\(1\)}{\(i\)}{}{}}
\End

\Data\ID{2000002105}
\Updated{2023-02-21 01:02:33}
\Level{A}
\SubArea{Powers and Roots of Complex Numbers}
\LANGen{
\Question{
Find the value of the principal argument \( \varphi\) of \( \left(\cos{\frac{2\pi}{3} }+ i\sin{\frac{2\pi}{3} }\right)^{18} \). The principal value is the corresponding angle \(\varphi \in (-\pi; \pi] \).}
{\( 0 \)}{\( \pi \)}{\( -\pi \)}{\( \frac{\pi}{3} \)}{}{}}
\LANGpl{
\Question{
Niech  \( \varphi\) będzie argumentem głównym, \( (-\pi < \varphi\leq \pi) \), liczby zespolonej.  Znajdż argument główny \( \varphi\) liczby \( \left(\cos{\frac{2\pi}{3} }+ i\sin{\frac{2\pi}{3} }\right)^{18} \).}
{\( 0 \)}{\( \pi \)}{\( -\pi \)}{\( \frac{\pi}{3} \)}{}{}}
\LANGcs{
\Question{
Určete hlavní hodnotu \( \varphi\) argumentu komplexního čísla \( \left(\cos{\frac{2\pi}{3} }+ i\sin{\frac{2\pi}{3} }\right)^{18} \). Hlavní hodnotou přitom rozumíme příslušný úhel \(\varphi \in (-\pi; \pi\rangle \).}
{\( 0 \)}{\( \pi \)}{\( -\pi \)}{\( \frac{\pi}{3} \)}{}{}}
\LANGsk{
\Question{
Určte hlavnú hodnotu \( \varphi\) argumentu komplexného čísla \( \left(\cos{\frac{2\pi}{3} }+ i\sin{\frac{2\pi}{3} }\right)^{18} \). Hlavnou hodnotou je príslušný uhol \(\varphi \in (-\pi; \pi\rangle \).}
{\( 0 \)}{\( \pi \)}{\( -\pi \)}{\( \frac{\pi}{3} \)}{}{}}
\LANGes{
\Question{
Sea \( \varphi\) el argumento principal de un número complejo \( (-\pi < \varphi \leq \pi) \). Determina el valor del argumento principal \( \varphi\) de \( \left(\cos{\frac{2\pi}{3} }+ i\sin{\frac{2\pi}{3} }\right)^{18} \).}
{\( 0 \)}{\( \pi \)}{\( -\pi \)}{\( \frac{\pi}{3} \)}{}{}}
\End

\Data\ID{2000002106}
\Updated{2021-12-26 08:12:25}
\Level{A}
\SubArea{Powers and Roots of Complex Numbers}
\LANGen{
\Question{
Which of the following numbers  is not the value of the argument of \( \left(\cos{\frac{\pi}{4}} + i\sin{\frac{\pi}{4}}\right)^5 \)?}
{\( \frac{\pi}{4} \)}{\( \frac{5\pi}{4} \)}{\( -\frac{3\pi}{4} \)}{\( \frac{13\pi}{4} \)}{}{}}
\LANGpl{
\Question{
Która z poniższych liczb nie jest wartością argumentu liczby \( \left(\cos{\frac{\pi}{4}} + i\sin{\frac{\pi}{4}}\right)^5 \)?}
{\( \frac{\pi}{4} \)}{\( \frac{5\pi}{4} \)}{\( -\frac{3\pi}{4} \)}{\( \frac{13\pi}{4} \)}{}{}}
\LANGcs{
\Question{
Které z uvedených čísel není hodnotou argumentu \( \left(\cos{\frac{\pi}{4}} + i\sin{\frac{\pi}{4}}\right)^5 \)?}
{\( \frac{\pi}{4} \)}{\( \frac{5\pi}{4} \)}{\( -\frac{3\pi}{4} \)}{\( \frac{13\pi}{4} \)}{}{}}
\LANGsk{
\Question{
Ktoré z uvedených čísel  nie je  hodnotou argumentu \( \left(\cos{\frac{\pi}{4}} + i\sin{\frac{\pi}{4}}\right)^5 \)?}
{\( \frac{\pi}{4} \)}{\( \frac{5\pi}{4} \)}{\( -\frac{3\pi}{4} \)}{\( \frac{13\pi}{4} \)}{}{}}
\LANGes{
\Question{
¿Cuál de los siguientes números no es el argumento de \( \left(\cos{\frac{\pi}{4}} + i\sin{\frac{\pi}{4}}\right)^5 \)?}
{\( \frac{\pi}{4} \)}{\( \frac{5\pi}{4} \)}{\( -\frac{3\pi}{4} \)}{\( \frac{13\pi}{4} \)}{}{}}
\End

\Data\ID{2000002107}
\Updated{2021-12-01 04:12:03}
\Level{A}
\SubArea{Powers and Roots of Complex Numbers}
\LANGen{
\Question{
Find the algebraic form of \( (\cos{45^{\circ}} +i\sin{45^{\circ}})^9 \).}
{\( \frac{\sqrt{2}}{2} + i\frac{\sqrt{2}}{2} \)}{\( -\frac{\sqrt{2}}{2} - i\frac{\sqrt{2}}{2} \)}{\( \frac{\sqrt{3}}{2} - \frac{i}{2} \)}{\( i +1 \)}{}{}}
\LANGpl{
\Question{
Znajdź postać algebraiczną liczby  \( (\cos{45^{\circ}} +i\sin{45^{\circ}})^9 \).}
{\( \frac{\sqrt{2}}{2} + i\frac{\sqrt{2}}{2} \)}{\( -\frac{\sqrt{2}}{2} - i\frac{\sqrt{2}}{2} \)}{\( \frac{\sqrt{3}}{2} - \frac{i}{2} \)}{\( i +1 \)}{}{}}
\LANGcs{
\Question{
Jak napíšete v algebraickém tvaru  \( (\cos{45^{\circ}} +i\sin{45^{\circ}})^9 \)?}
{\( \frac{\sqrt{2}}{2} + i\frac{\sqrt{2}}{2} \)}{\( -\frac{\sqrt{2}}{2} - i\frac{\sqrt{2}}{2} \)}{\( \frac{\sqrt{3}}{2} - \frac{i}{2} \)}{\( i +1 \)}{}{}}
\LANGsk{
\Question{
Nájdi algebraický tvar \( (\cos{45^{\circ}} +i\sin{45^{\circ}})^9 \).}
{\( \frac{\sqrt{2}}{2} + i\frac{\sqrt{2}}{2} \)}{\( -\frac{\sqrt{2}}{2} - i\frac{\sqrt{2}}{2} \)}{\( \frac{\sqrt{3}}{2} - \frac{i}{2} \)}{\( i +1 \)}{}{}}
\LANGes{
\Question{
Determina la forma algebraica de \( (\cos{45^{\circ}} +i\sin{45^{\circ}})^9 \).}
{\( \frac{\sqrt{2}}{2} + i\frac{\sqrt{2}}{2} \)}{\( -\frac{\sqrt{2}}{2} - i\frac{\sqrt{2}}{2} \)}{\( \frac{\sqrt{3}}{2} - \frac{i}{2} \)}{\( i +1 \)}{}{}}
\End

\Data\ID{2000002108}
\Updated{2023-02-21 01:02:25}
\Level{A}
\SubArea{Powers and Roots of Complex Numbers}
\LANGen{
\Question{
Find the value of the principal argument \( \varphi\) of \( \left(3\left(\cos{\frac{3\pi}{2} }+ i\sin{\frac{3\pi}{2} }\right)\right)^{13} \). The principal value is the corresponding angle \(\varphi \in (-\pi; \pi] \).}
{\( -\frac{\pi}{2} \)}{\( \frac{\pi}{2} \)}{\( 0 \)}{\( \frac{3}{26}\pi \)}{}{}}
\LANGpl{
\Question{
Niech \( \varphi\) będzie argumentem głównym, \( (-\pi < \varphi\leq \pi) \), liczby zespolonej. Znajdź argument główny \( \varphi\) liczby \( \left(3\left(\cos{\frac{3\pi}{2} }+ i\sin{\frac{3\pi}{2} }\right)\right)^{13} \).}
{\( -\frac{\pi}{2} \)}{\( \frac{\pi}{2} \)}{\( 0 \)}{\( \frac{3}{26}\pi \)}{}{}}
\LANGcs{
\Question{
Určete hlavní hodnotu \( \varphi\) argumentu komplexního čísla \( \left(3\left(\cos{\frac{3\pi}{2} }+ i\sin{\frac{3\pi}{2} }\right)\right)^{13} \). Hlavní hodnotou přitom rozumíme příslušný úhel \(\varphi \in (-\pi; \pi\rangle \).}
{\( -\frac{\pi}{2} \)}{\( \frac{\pi}{2} \)}{\( 0 \)}{\( \frac{3}{26}\pi \)}{}{}}
\LANGsk{
\Question{
Určte hlavnú hodnotu \( \varphi\) argumentu komplexného čísla \( \left(3\left(\cos{\frac{3\pi}{2} }+ i\sin{\frac{3\pi}{2} }\right)\right)^{13} \). Hlavnou hodnotou je príslušný uhol \(\varphi \in (-\pi; \pi\rangle \).}
{\( -\frac{\pi}{2} \)}{\( \frac{\pi}{2} \)}{\( 0 \)}{\( \frac{3}{26}\pi \)}{}{}}
\LANGes{
\Question{
Sea \( \varphi\) el argumento principal de un número complejo \( (-\pi < \varphi\leq \pi) \). Determina el valor del argumento principal \( \varphi\) de \( \left(3\left(\cos{\frac{3\pi}{2} }+ i\sin{\frac{3\pi}{2} }\right)\right)^{13} \).}
{\( -\frac{\pi}{2} \)}{\( \frac{\pi}{2} \)}{\( 0 \)}{\( \frac{3}{26}\pi \)}{}{}}
\End

\Data\ID{2000002109}
\Updated{2021-12-01 04:12:41}
\Level{A}
\SubArea{Powers and Roots of Complex Numbers}
\LANGen{
\Question{
Given \( z= \sqrt[3]{3} \left(\cos{\frac{3\pi}{4}}+i\sin{\frac{3\pi}{4}}\right) \), determine which of the numbers below does not represent  \( z^6\).}
{\( 9 \)}{\( 9i \)}{\( 9\left(\cos{\frac{9\pi}{2}}+i\sin{\frac{9\pi}{2}}\right) \)}{\( 9\left(\cos{\frac{\pi}{2}}+i\sin{\frac{\pi}{2}}\right) \)}{}{}}
\LANGpl{
\Question{
Dane \( z= \sqrt[3]{3} \left(\cos{\frac{3\pi}{4}}+i\sin{\frac{3\pi}{4}}\right) \), określ, która z poniższych liczb  nie reprezentuje \( z^6\).}
{\( 9 \)}{\( 9i \)}{\( 9\left(\cos{\frac{9\pi}{2}}+i\sin{\frac{9\pi}{2}}\right) \)}{\( 9\left(\cos{\frac{\pi}{2}}+i\sin{\frac{\pi}{2}}\right) \)}{}{}}
\LANGcs{
\Question{
Je dáno komplexní číslo \( z= \sqrt[3]{3} \left(\cos{\frac{3\pi}{4}}+i\sin{\frac{3\pi}{4}}\right) \). Určete, které z níže uvedených čísel nevyjadřuje \( z^6\).}
{\( 9 \)}{\( 9i \)}{\( 9\left(\cos{\frac{9\pi}{2}}+i\sin{\frac{9\pi}{2}}\right) \)}{\( 9\left(\cos{\frac{\pi}{2}}+i\sin{\frac{\pi}{2}}\right) \)}{}{}}
\LANGsk{
\Question{
Je dané komplexné číslo \( z= \sqrt[3]{3} \left(\cos{\frac{3\pi}{4}}+i\sin{\frac{3\pi}{4}}\right) \), určte, ktoré z nasledujúcich čísel nevyjadruje   \( z^6\).}
{\( 9 \)}{\( 9i \)}{\( 9\left(\cos{\frac{9\pi}{2}}+i\sin{\frac{9\pi}{2}}\right) \)}{\( 9\left(\cos{\frac{\pi}{2}}+i\sin{\frac{\pi}{2}}\right) \)}{}{}}
\LANGes{
\Question{
Dado \( z= \sqrt[3]{3} \left(\cos{\frac{3\pi}{4}}+i\sin{\frac{3\pi}{4}}\right) \), determina cuál de los siguientes números  no representa   \( z^6\).}
{\( 9 \)}{\( 9i \)}{\( 9\left(\cos{\frac{9\pi}{2}}+i\sin{\frac{9\pi}{2}}\right) \)}{\( 9\left(\cos{\frac{\pi}{2}}+i\sin{\frac{\pi}{2}}\right) \)}{}{}}
\End

\Data\ID{2000002110}
\Updated{2021-12-01 04:12:46}
\Level{A}
\SubArea{Powers and Roots of Complex Numbers}
\LANGen{
\Question{
Find the algebraic form of the complex number \( \left(\cos{ \frac{\pi}{6}} + i\sin{ \frac{\pi}{6}}\right)^{-6} \).}
{\( -1\)}{\( i\)}{\( -i \)}{\( 1\)}{}{}}
\LANGpl{
\Question{
Znajdź postać algebraiczną liczby zespolonej \( \left(\cos{ \frac{\pi}{6}} + i\sin{ \frac{\pi}{6}}\right)^{-6} \).}
{\( -1\)}{\( i\)}{\( -i \)}{\( 1\)}{}{}}
\LANGcs{
\Question{
Jaký je algebraický tvar komplexního čísla \( \left(\cos{ \frac{\pi}{6}} + i\sin{ \frac{\pi}{6}}\right)^{-6} \)?}
{\( -1\)}{\( i\)}{\( -i \)}{\( 1\)}{}{}}
\LANGsk{
\Question{
Aký je algebraický tvar komplexného čísla \( \left(\cos{ \frac{\pi}{6}} + i\sin{ \frac{\pi}{6}}\right)^{-6} \)?}
{\( -1\)}{\( i\)}{\( -i \)}{\( 1\)}{}{}}
\LANGes{
\Question{
Determina la forma algebraica del número complejo \( \left(\cos{ \frac{\pi}{6}} + i\sin{ \frac{\pi}{6}}\right)^{-6} \).}
{\( -1\)}{\( i\)}{\( -i \)}{\( 1\)}{}{}}
\End

\Data\ID{2010004601}
\Updated{2022-08-25 10:08:26}
\Level{A}
\SubArea{Powers and Roots of Complex Numbers}
\LANGen{
\Question{
Evaluate \(\left (\cos \frac{\pi }{4} + \mathrm{i}\sin \frac{\pi }{4}\right )^{40}\) and find the algebraic form of the result.}
{\(1\)}{\(-1\)}{\( -\mathrm{i}\)}{\( \mathrm{i} \)}{}{}}
\LANGpl{
\Question{
Oblicz \(\left (\cos \frac{\pi }{4} + \mathrm{i}\sin \frac{\pi }{4}\right )^{40}\) i wynik zapisz w postaci algebraicznej.}
{\(1\)}{\(-1\)}{\( -\mathrm{i}\)}{\( \mathrm{i} \)}{}{}}
\LANGcs{
\Question{
Zjednodušte \(\left (\cos \frac{\pi }{4} + \mathrm{i}\sin \frac{\pi }{4}\right )^{40}\) a zapište výraz v algebraickém tvaru.}
{\(1\)}{\(-1\)}{\( -\mathrm{i}\)}{\( \mathrm{i} \)}{}{}}
\LANGsk{
\Question{
Zjednodušte \(\left (\cos \frac{\pi }{4} + \mathrm{i}\sin \frac{\pi }{4}\right )^{40}\) a zapíšte výraz v algebrickom tvare.}
{\(1\)}{\(-1\)}{\( -\mathrm{i}\)}{\( \mathrm{i} \)}{}{}}
\LANGes{
\Question{
Simplifica \(\left (\cos \frac{\pi }{4} + \mathrm{i}\sin \frac{\pi }{4}\right )^{40}\) y determina la forma algebraica del resultado.}
{\(1\)}{\(-1\)}{\( -\mathrm{i}\)}{\( \mathrm{i} \)}{}{}}
\End

\Data\ID{2010004602}
\Updated{2022-08-25 10:08:26}
\Level{A}
\SubArea{Powers and Roots of Complex Numbers}
\LANGen{
\Question{
Identify a complex number that is not equal to \(\left (\cos \frac{\pi }{3} + \mathrm{i}\sin \frac{\pi }{3}\right )^{19}\).}
{\( \frac{1}{2} - \frac{\sqrt{3}}{2}\mathrm{i}\)}{\( \cos \frac{\pi }{3} + \mathrm{i}\sin \frac{\pi }{3} \)}{\( \frac{1}{2} + \frac{\sqrt{3}}{2}\mathrm{i}\)}{\( \cos \left(-\frac{5\pi }{3}\right) + \mathrm{i}\sin \left(-\frac{5\pi }{3}\right) \)}{}{}}
\LANGpl{
\Question{
Wskaż liczbę zespoloną, która nie jest równa liczbie  \(\left (\cos \frac{\pi }{3} + \mathrm{i}\sin \frac{\pi }{3}\right )^{19}\).}
{\( \frac{1}{2} - \frac{\sqrt{3}}{2}\mathrm{i}\)}{\( \cos \frac{\pi }{3} + \mathrm{i}\sin \frac{\pi }{3} \)}{\( \frac{1}{2} + \frac{\sqrt{3}}{2}\mathrm{i}\)}{\( \cos \left(-\frac{5\pi }{3}\right) + \mathrm{i}\sin \left(-\frac{5\pi }{3}\right) \)}{}{}}
\LANGcs{
\Question{
Najděte komplexní číslo, které není rovno \(\left (\cos \frac{\pi }{3} + \mathrm{i}\sin \frac{\pi }{3}\right )^{19}\).}
{\( \frac{1}{2} - \frac{\sqrt{3}}{2}\mathrm{i}\)}{\( \cos \frac{\pi }{3} + \mathrm{i}\sin \frac{\pi }{3} \)}{\( \frac{1}{2} + \frac{\sqrt{3}}{2}\mathrm{i}\)}{\( \cos \left(-\frac{5\pi }{3}\right) + \mathrm{i}\sin \left(-\frac{5\pi }{3}\right) \)}{}{}}
\LANGsk{
\Question{
Nájdite komplexné číslo, ktoré nie je rovné \(\left (\cos \frac{\pi }{3} + \mathrm{i}\sin \frac{\pi }{3}\right )^{19}\).}
{\( \frac{1}{2} - \frac{\sqrt{3}}{2}\mathrm{i}\)}{\( \cos \frac{\pi }{3} + \mathrm{i}\sin \frac{\pi }{3} \)}{\( \frac{1}{2} + \frac{\sqrt{3}}{2}\mathrm{i}\)}{\( \cos \left(-\frac{5\pi }{3}\right) + \mathrm{i}\sin \left(-\frac{5\pi }{3}\right) \)}{}{}}
\LANGes{
\Question{
Identifica un número complejo que no es igual a \(\left (\cos \frac{\pi }{3} + \mathrm{i}\sin \frac{\pi }{3}\right )^{19}\).}
{\( \frac{1}{2} - \frac{\sqrt{3}}{2}\mathrm{i}\)}{\( \cos \frac{\pi }{3} + \mathrm{i}\sin \frac{\pi }{3} \)}{\( \frac{1}{2} + \frac{\sqrt{3}}{2}\mathrm{i}\)}{\( \cos \left(-\frac{5\pi }{3}\right) + \mathrm{i}\sin \left(-\frac{5\pi }{3}\right) \)}{}{}}
\End

\Data\ID{2010004603}
\Updated{2022-08-25 10:08:26}
\Level{A}
\SubArea{Powers and Roots of Complex Numbers}
\LANGen{
\Question{
Identify a complex number that is not equal to \(\left (\cos \frac{\pi }{6} + \mathrm{i}\sin \frac{\pi }{6}\right )^{13}\).}
{\( \frac{\sqrt{3}}{2} - \frac{1}{2}\mathrm{i}\)}{\( \frac{\sqrt{3}}{2} + \frac{1}{2}\mathrm{i}\)}{\( \cos \frac{\pi }{6} + \mathrm{i}\sin \frac{\pi }{6} \)}{\( \cos \frac{25\pi }{6} + \mathrm{i}\sin \frac{25\pi }{6} \)}{}{}}
\LANGpl{
\Question{
Wskaż liczbę zespoloną, która nie jest równa liczbie  \(\left (\cos \frac{\pi }{6} + \mathrm{i}\sin \frac{\pi }{6}\right )^{13}\).}
{\( \frac{\sqrt{3}}{2} - \frac{1}{2}\mathrm{i}\)}{\( \frac{\sqrt{3}}{2} + \frac{1}{2}\mathrm{i}\)}{\( \cos \frac{\pi }{6} + \mathrm{i}\sin \frac{\pi }{6} \)}{\( \cos \frac{25\pi }{6} + \mathrm{i}\sin \frac{25\pi }{6} \)}{}{}}
\LANGcs{
\Question{
Najděte komplexní číslo, které není rovno \(\left (\cos \frac{\pi }{6} + \mathrm{i}\sin \frac{\pi }{6}\right )^{13}\).}
{\( \frac{\sqrt{3}}{2} - \frac{1}{2}\mathrm{i}\)}{\( \frac{\sqrt{3}}{2} + \frac{1}{2}\mathrm{i}\)}{\( \cos \frac{\pi }{6} + \mathrm{i}\sin \frac{\pi }{6} \)}{\( \cos \frac{25\pi }{6} + \mathrm{i}\sin \frac{25\pi }{6} \)}{}{}}
\LANGsk{
\Question{
Nájdite komplexné číslo, ktoré nie je rovné \(\left (\cos \frac{\pi }{6} + \mathrm{i}\sin \frac{\pi }{6}\right )^{13}\).}
{\( \frac{\sqrt{3}}{2} - \frac{1}{2}\mathrm{i}\)}{\( \frac{\sqrt{3}}{2} + \frac{1}{2}\mathrm{i}\)}{\( \cos \frac{\pi }{6} + \mathrm{i}\sin \frac{\pi }{6} \)}{\( \cos \frac{25\pi }{6} + \mathrm{i}\sin \frac{25\pi }{6} \)}{}{}}
\LANGes{
\Question{
Identifica un número complejo que no es igual a \(\left (\cos \frac{\pi }{6} + \mathrm{i}\sin \frac{\pi }{6}\right )^{13}\).}
{\( \frac{\sqrt{3}}{2} - \frac{1}{2}\mathrm{i}\)}{\( \frac{\sqrt{3}}{2} + \frac{1}{2}\mathrm{i}\)}{\( \cos \frac{\pi }{6} + \mathrm{i}\sin \frac{\pi }{6} \)}{\( \cos \frac{25\pi }{6} + \mathrm{i}\sin \frac{25\pi }{6} \)}{}{}}
\End

\Data\ID{2010004604}
\Updated{2022-08-25 10:08:26}
\Level{A}
\SubArea{Powers and Roots of Complex Numbers}
\LANGen{
\Question{
Find the algebraic form of the complex number:
\[ 
 \left (\cos \frac{\pi }
{3} + \mathrm{i}\sin \frac{\pi }
{3}\right )^{5}
\]}
{\(\frac{1} 
 {2} - \mathrm{i}\frac{\sqrt{3}} 
 {2} \)}{\(\frac{\sqrt{3}} 
 {2} +\mathrm{i} \frac{1} 
 {2} \)}{\(\frac{\sqrt{3}} 
 {2} -\mathrm{i} \frac{1} 
 {2} \)}{\(\frac{1} 
 {2} + \mathrm{i}\frac{\sqrt{3}} 
 {2} \)}{}{}}
\LANGpl{
\Question{
Zapisz w postaci algebraicznej liczbę:
\[ 
 \left (\cos \frac{\pi }
{3} + \mathrm{i}\sin \frac{\pi }
{3}\right )^{5}
\]}
{\(\frac{1} 
 {2} - \mathrm{i}\frac{\sqrt{3}} 
 {2} \)}{\(\frac{\sqrt{3}} 
 {2} +\mathrm{i} \frac{1} 
 {2} \)}{\(\frac{\sqrt{3}} 
 {2} -\mathrm{i} \frac{1} 
 {2} \)}{\(\frac{1} 
 {2} + \mathrm{i}\frac{\sqrt{3}} 
 {2} \)}{}{}}
\LANGcs{
\Question{
Zjednodušte a zapište v algebraickém tvaru komplexní číslo:
\[ 
 \left (\cos \frac{\pi }
{3} + \mathrm{i}\sin \frac{\pi }
{3}\right )^{5}
\]}
{\(\frac{1} 
 {2} - \mathrm{i}\frac{\sqrt{3}} 
 {2} \)}{\(\frac{\sqrt{3}} 
 {2} +\mathrm{i} \frac{1} 
 {2} \)}{\(\frac{\sqrt{3}} 
 {2} -\mathrm{i} \frac{1} 
 {2} \)}{\(\frac{1} 
 {2} + \mathrm{i}\frac{\sqrt{3}} 
 {2} \)}{}{}}
\LANGsk{
\Question{
Zjednodušte a zapíšte v algebrickom tvare komplexné číslo:
\[ 
 \left (\cos \frac{\pi }
{3} + \mathrm{i}\sin \frac{\pi }
{3}\right )^{5}
\]}
{\(\frac{1} 
 {2} - \mathrm{i}\frac{\sqrt{3}} 
 {2} \)}{\(\frac{\sqrt{3}} 
 {2} +\mathrm{i} \frac{1} 
 {2} \)}{\(\frac{\sqrt{3}} 
 {2} -\mathrm{i} \frac{1} 
 {2} \)}{\(\frac{1} 
 {2} + \mathrm{i}\frac{\sqrt{3}} 
 {2} \)}{}{}}
\LANGes{
\Question{
Calcula el siguiente número complejo.
\[ 
 \left (\cos \frac{\pi }
{3} + \mathrm{i}\sin \frac{\pi }
{3}\right )^{5}
\]}
{\(\frac{1} 
 {2} - \mathrm{i}\frac{\sqrt{3}} 
 {2} \)}{\(\frac{\sqrt{3}} 
 {2} +\mathrm{i} \frac{1} 
 {2} \)}{\(\frac{\sqrt{3}} 
 {2} -\mathrm{i} \frac{1} 
 {2} \)}{\(\frac{1} 
 {2} + \mathrm{i}\frac{\sqrt{3}} 
 {2} \)}{}{}}
\End

\Data\ID{2010004605}
\Updated{2022-08-25 10:08:26}
\Level{A}
\SubArea{Powers and Roots of Complex Numbers}
\LANGen{
\Question{
Find the algebraic form of the complex number:
\[ 
 \left (\frac{1}
{2} +\sin \frac{\pi } 
{6} + \mathrm{i}\sin \frac{\pi}{2} \right )^{3}
\]}
{\(- 2 +2\mathrm{i}\)}{\(2 -2\mathrm{i}\)}{\(-2 - 2\mathrm{i}\)}{\(2 + 2\mathrm{i}\)}{}{}}
\LANGpl{
\Question{
Zapisz w postaci algebraicznej liczbę:
\[ 
 \left (\frac{1}
{2} +\sin \frac{\pi } 
{6} + \mathrm{i}\sin \frac{\pi}{2} \right )^{3}
\]}
{\(- 2 +2\mathrm{i}\)}{\(2 -2\mathrm{i}\)}{\(-2 - 2\mathrm{i}\)}{\(2 + 2\mathrm{i}\)}{}{}}
\LANGcs{
\Question{
Zjednodušte a zapište v algebraickém tvaru komplexní číslo:
\[ 
 \left (\frac{1}
{2} +\sin \frac{\pi } 
{6} + \mathrm{i}\sin \frac{\pi}{2} \right )^{3}
\]}
{\(- 2 +2\mathrm{i}\)}{\(2 -2\mathrm{i}\)}{\(-2 - 2\mathrm{i}\)}{\(2 + 2\mathrm{i}\)}{}{}}
\LANGsk{
\Question{
Zjednodušte a zapíšte v algebrickom tvare komplexné číslo:
\[ 
 \left (\frac{1}
{2} +\sin \frac{\pi } 
{6} + \mathrm{i}\sin \frac{\pi}{2} \right )^{3}
\]}
{\(- 2 +2\mathrm{i}\)}{\(2 -2\mathrm{i}\)}{\(-2 - 2\mathrm{i}\)}{\(2 + 2\mathrm{i}\)}{}{}}
\LANGes{
\Question{
Calcula el siguiente número complejo.
\[ 
 \left (\frac{1}
{2} +\sin \frac{\pi } 
{6} + \mathrm{i}\sin \frac{\pi}{2} \right )^{3}
\]}
{\(- 2 +2\mathrm{i}\)}{\(2 -2\mathrm{i}\)}{\(-2 - 2\mathrm{i}\)}{\(2 + 2\mathrm{i}\)}{}{}}
\End

\Data\ID{2010004606}
\Updated{2023-02-21 10:02:12}
\Level{A}
\SubArea{Powers and Roots of Complex Numbers}
\LANGen{
\Question{
Evaluate \( (\cos 33^{\circ} + \mathrm{i}\sin 33^{\circ})^{10} \) and find the algebraic form of the result.}
{\( \frac{\sqrt{3}}{2} - \frac{1}{2}\mathrm{i}\)}{\( \frac{\sqrt{3}}{2} + \frac{1}{2}\mathrm{i}\)}{\( \frac{1}{2} + \frac{\sqrt{3}}{2}\mathrm{i}\)}{\( \frac{1}{2} - \frac{\sqrt{3}}{2}\mathrm{i}\)}{}{}}
\LANGpl{
\Question{
Oblicz \( (\cos 33^{\circ} + \mathrm{i}\sin 33^{\circ})^{10} \) i wynik zapisz w postaci algebraicznej.}
{\( \frac{\sqrt{3}}{2} - \frac{1}{2}\mathrm{i}\)}{\( \frac{\sqrt{3}}{2} + \frac{1}{2}\mathrm{i}\)}{\( \frac{1}{2} + \frac{\sqrt{3}}{2}\mathrm{i}\)}{\( \frac{1}{2} - \frac{\sqrt{3}}{2}\mathrm{i}\)}{}{}}
\LANGcs{
\Question{
Zjednodušte \( (\cos 33^{\circ} + \mathrm{i}\sin 33^{\circ})^{10} \) a zapište v algebraickém tvaru.}
{\( \frac{\sqrt{3}}{2} - \frac{1}{2}\mathrm{i}\)}{\( \frac{\sqrt{3}}{2} + \frac{1}{2}\mathrm{i}\)}{\( \frac{1}{2} + \frac{\sqrt{3}}{2}\mathrm{i}\)}{\( \frac{1}{2} - \frac{\sqrt{3}}{2}\mathrm{i}\)}{}{}}
\LANGsk{
\Question{
Zjednodušte \( (\cos 33^{\circ} + \mathrm{i}\sin 33^{\circ})^{10} \) a zapíšte v algebrickom tvare.}
{\( \frac{\sqrt{3}}{2} - \frac{1}{2}\mathrm{i}\)}{\( \frac{\sqrt{3}}{2} + \frac{1}{2}\mathrm{i}\)}{\( \frac{1}{2} + \frac{\sqrt{3}}{2}\mathrm{i}\)}{\( \frac{1}{2} - \frac{\sqrt{3}}{2}\mathrm{i}\)}{}{}}
\LANGes{
\Question{
Simplifica \( (\cos 33^{\circ} + \mathrm{i}\sin 33^{\circ})^{10} \) y determina la forma algebraica del resultado.}
{\( \frac{\sqrt{3}}{2} - \frac{1}{2}\mathrm{i}\)}{\( \frac{\sqrt{3}}{2} + \frac{1}{2}\mathrm{i}\)}{\( \frac{1}{2} + \frac{\sqrt{3}}{2}\mathrm{i}\)}{\( \frac{1}{2} - \frac{\sqrt{3}}{2}\mathrm{i}\)}{}{}}
\End

\Data\ID{2010004607}
\Updated{2023-02-21 10:02:43}
\Level{A}
\SubArea{Powers and Roots of Complex Numbers}
\LANGen{
\Question{
Evaluate \( (\cos 27^{\circ} + \mathrm{i}\sin 27^{\circ})^{5} \) and find the algebraic form of the result.}
{\( -\frac{\sqrt{2}}{2} + \frac{\sqrt{2}}{2}\mathrm{i}\)}{\( \frac{\sqrt{2}}{2} + \frac{\sqrt{2}}{2}\mathrm{i}\)}{\( -\frac{\sqrt{2}}{2} - \frac{\sqrt{2}}{2}\mathrm{i}\)}{\( \frac{\sqrt{2}}{2} - \frac{\sqrt{2}}{2}\mathrm{i}\)}{}{}}
\LANGpl{
\Question{
Oblicz \( (\cos 27^{\circ} + \mathrm{i}\sin 27^{\circ})^{5} \) i wynik zapisz w postaci algebraicznej.}
{\( -\frac{\sqrt{2}}{2} + \frac{\sqrt{2}}{2}\mathrm{i}\)}{\( \frac{\sqrt{2}}{2} + \frac{\sqrt{2}}{2}\mathrm{i}\)}{\( -\frac{\sqrt{2}}{2} - \frac{\sqrt{2}}{2}\mathrm{i}\)}{\( \frac{\sqrt{2}}{2} - \frac{\sqrt{2}}{2}\mathrm{i}\)}{}{}}
\LANGcs{
\Question{
Zjednodušte \( (\cos 27^{\circ} + \mathrm{i}\sin 27^{\circ})^{5} \) a zapište v algebraickém tvaru.}
{\( -\frac{\sqrt{2}}{2} + \frac{\sqrt{2}}{2}\mathrm{i}\)}{\( \frac{\sqrt{2}}{2} + \frac{\sqrt{2}}{2}\mathrm{i}\)}{\( -\frac{\sqrt{2}}{2} - \frac{\sqrt{2}}{2}\mathrm{i}\)}{\( \frac{\sqrt{2}}{2} - \frac{\sqrt{2}}{2}\mathrm{i}\)}{}{}}
\LANGsk{
\Question{
Zjednodušte \( (\cos 27^{\circ} + \mathrm{i}\sin 27^{\circ})^{5} \) a zapíšte v algebrickom tvare.}
{\( -\frac{\sqrt{2}}{2} + \frac{\sqrt{2}}{2}\mathrm{i}\)}{\( \frac{\sqrt{2}}{2} + \frac{\sqrt{2}}{2}\mathrm{i}\)}{\( -\frac{\sqrt{2}}{2} - \frac{\sqrt{2}}{2}\mathrm{i}\)}{\( \frac{\sqrt{2}}{2} - \frac{\sqrt{2}}{2}\mathrm{i}\)}{}{}}
\LANGes{
\Question{
Simplifica \( (\cos 27^{\circ} + \mathrm{i}\sin 27^{\circ})^{5} \) y determina la forma algebraica del resultado.}
{\( -\frac{\sqrt{2}}{2} + \frac{\sqrt{2}}{2}\mathrm{i}\)}{\( \frac{\sqrt{2}}{2} + \frac{\sqrt{2}}{2}\mathrm{i}\)}{\( -\frac{\sqrt{2}}{2} - \frac{\sqrt{2}}{2}\mathrm{i}\)}{\( \frac{\sqrt{2}}{2} - \frac{\sqrt{2}}{2}\mathrm{i}\)}{}{}}
\End

\Data\ID{2010004608}
\Updated{2022-08-25 10:08:30}
\Level{A}
\SubArea{Powers and Roots of Complex Numbers}
\LANGen{
\Question{
Evaluate \( (1+\mathrm{i})^{22} \) and write the result in algebraic form.}
{\( -2^{11}\mathrm{i} \)}{\( 2^{11}\mathrm{i} \)}{\( 2^{22}\mathrm{i} \)}{\( -2^{22}\mathrm{i} \)}{}{}}
\LANGpl{
\Question{
Oblicz \( (1+\mathrm{i})^{22} \) i wynik zapisz w postaci algebraicznej.}
{\( -2^{11}\mathrm{i} \)}{\( 2^{11}\mathrm{i} \)}{\( 2^{22}\mathrm{i} \)}{\( -2^{22}\mathrm{i} \)}{}{}}
\LANGcs{
\Question{
Zjednodušte \( (1+\mathrm{i})^{22} \) a zapište výsledek v algebraickém tvaru.}
{\( -2^{11}\mathrm{i} \)}{\( 2^{11}\mathrm{i} \)}{\( 2^{22}\mathrm{i} \)}{\( -2^{22}\mathrm{i} \)}{}{}}
\LANGsk{
\Question{
Zjednodušte \( (1+\mathrm{i})^{22} \) a zapíšte výsledok v algebrickom tvare.}
{\( -2^{11}\mathrm{i} \)}{\( 2^{11}\mathrm{i} \)}{\( 2^{22}\mathrm{i} \)}{\( -2^{22}\mathrm{i} \)}{}{}}
\LANGes{
\Question{
Simplifica \( (1+\mathrm{i})^{22} \) y escribe el resultado en forma algebraica.}
{\( -2^{11}\mathrm{i} \)}{\( 2^{11}\mathrm{i} \)}{\( 2^{22}\mathrm{i} \)}{\( -2^{22}\mathrm{i} \)}{}{}}
\End

\Data\ID{2010004609}
\Updated{2022-08-25 10:08:30}
\Level{A}
\SubArea{Powers and Roots of Complex Numbers}
\LANGen{
\Question{
Evaluate \( (2+2\mathrm{i})^{100} \) and write the result in algebraic form.}
{\( -2^{150} \)}{\( -2^{150}\mathrm{i} \)}{\( 2^{150}\mathrm{i} \)}{\( 2^{150}\)}{}{}}
\LANGpl{
\Question{
Oblicz \( (2+2\mathrm{i})^{100} \) i wynik zapisz w postaci algebraicznej.}
{\( -2^{150} \)}{\( -2^{150}\mathrm{i} \)}{\( 2^{150}\mathrm{i} \)}{\( 2^{150}\)}{}{}}
\LANGcs{
\Question{
Zjednodušte \( (2+2\mathrm{i})^{100} \) a zapište výsledek v algebraickém tvaru.}
{\( -2^{150} \)}{\( -2^{150}\mathrm{i} \)}{\( 2^{150}\mathrm{i} \)}{\( 2^{150}\)}{}{}}
\LANGsk{
\Question{
Zjednodušte \( (2+2\mathrm{i})^{100} \) a zapíšte výsledok v algebrickom tvare.}
{\( -2^{150} \)}{\( -2^{150}\mathrm{i} \)}{\( 2^{150}\mathrm{i} \)}{\( 2^{150}\)}{}{}}
\LANGes{
\Question{
Simplifica \( (2+2\mathrm{i})^{100} \)  y escribe el resultado en forma algebraica.}
{\( -2^{150} \)}{\( -2^{150}\mathrm{i} \)}{\( 2^{150}\mathrm{i} \)}{\( 2^{150}\)}{}{}}
\End

\Data\ID{2010004610}
\Updated{2022-08-25 10:08:30}
\Level{A}
\SubArea{Powers and Roots of Complex Numbers}
\LANGen{
\Question{
Identify the complex number that equals to  \((\sqrt{3} - \mathrm{i})^{37} \).}
{\(2^{36}(\sqrt{3}-\mathrm{i})\)}{\(2^{36}(-1+\mathrm{i}\sqrt{3})\)}{\(2^{36}(\sqrt{3}+\mathrm{i})\)}{\(2^{36}(-1-\mathrm{i}\sqrt{3})\)}{}{}}
\LANGpl{
\Question{
Znajdź liczbę zespoloną równą liczbie  \((\sqrt{3} - \mathrm{i})^{37} \).}
{\(2^{36}(\sqrt{3}-\mathrm{i})\)}{\(2^{36}(-1+\mathrm{i}\sqrt{3})\)}{\(2^{36}(\sqrt{3}+\mathrm{i})\)}{\(2^{36}(-1-\mathrm{i}\sqrt{3})\)}{}{}}
\LANGcs{
\Question{
Najděte komplexní číslo, které je rovno \((\sqrt{3} - \mathrm{i})^{37} \).}
{\(2^{36}(\sqrt{3}-\mathrm{i})\)}{\(2^{36}(-1+\mathrm{i}\sqrt{3})\)}{\(2^{36}(\sqrt{3}+\mathrm{i})\)}{\(2^{36}(-1-\mathrm{i}\sqrt{3})\)}{}{}}
\LANGsk{
\Question{
Nájdite komplexné číslo, ktoré je rovné \((\sqrt{3} - \mathrm{i})^{37} \).}
{\(2^{36}(\sqrt{3}-\mathrm{i})\)}{\(2^{36}(-1+\mathrm{i}\sqrt{3})\)}{\(2^{36}(\sqrt{3}+\mathrm{i})\)}{\(2^{36}(-1-\mathrm{i}\sqrt{3})\)}{}{}}
\LANGes{
\Question{
Identifica el número complejo que equivale a \((\sqrt{3} - \mathrm{i})^{37} \).}
{\(2^{36}(\sqrt{3}-\mathrm{i})\)}{\(2^{36}(-1+\mathrm{i}\sqrt{3})\)}{\(2^{36}(\sqrt{3}+\mathrm{i})\)}{\(2^{36}(-1-\mathrm{i}\sqrt{3})\)}{}{}}
\End

\Data\ID{2010004611}
\Updated{2022-08-25 10:08:30}
\Level{A}
\SubArea{Powers and Roots of Complex Numbers}
\LANGen{
\Question{
Identify the complex number that equals to  \(( -1+ \mathrm{i}\sqrt{3})^{67} \).}
{\(2^{66}(-1+\mathrm{i}\sqrt{3})\)}{\(2^{66}(1-\mathrm{i}\sqrt{3})\)}{\(2^{66}(\sqrt{3}+\mathrm{i})\)}{\(2^{66}(\sqrt{3}-\mathrm{i})\)}{}{}}
\LANGpl{
\Question{
Znajdź liczbę zespoloną równą liczbie  \(( -1+ \mathrm{i}\sqrt{3})^{67} \).}
{\(2^{66}(-1+\mathrm{i}\sqrt{3})\)}{\(2^{66}(1-\mathrm{i}\sqrt{3})\)}{\(2^{66}(\sqrt{3}+\mathrm{i})\)}{\(2^{66}(\sqrt{3}-\mathrm{i})\)}{}{}}
\LANGcs{
\Question{
Které komplexní číslo je rovno  \(( -1+ \mathrm{i}\sqrt{3})^{67} \)?}
{\(2^{66}(-1+\mathrm{i}\sqrt{3})\)}{\(2^{66}(1-\mathrm{i}\sqrt{3})\)}{\(2^{66}(\sqrt{3}+\mathrm{i})\)}{\(2^{66}(\sqrt{3}-\mathrm{i})\)}{}{}}
\LANGsk{
\Question{
Ktoré komplexné číslo je rovné  \(( -1+ \mathrm{i}\sqrt{3})^{67} \)?}
{\(2^{66}(-1+\mathrm{i}\sqrt{3})\)}{\(2^{66}(1-\mathrm{i}\sqrt{3})\)}{\(2^{66}(\sqrt{3}+\mathrm{i})\)}{\(2^{66}(\sqrt{3}-\mathrm{i})\)}{}{}}
\LANGes{
\Question{
Identifica el número complejo que equivale a \(( -1+ \mathrm{i}\sqrt{3})^{67} \).}
{\(2^{66}(-1+\mathrm{i}\sqrt{3})\)}{\(2^{66}(1-\mathrm{i}\sqrt{3})\)}{\(2^{66}(\sqrt{3}+\mathrm{i})\)}{\(2^{66}(\sqrt{3}-\mathrm{i})\)}{}{}}
\End

\Data\ID{2010004612}
\Updated{2022-08-25 10:08:30}
\Level{A}
\SubArea{Powers and Roots of Complex Numbers}
\LANGen{
\Question{
Find the absolute value and the value of the argument of \( z=(1-\mathrm{i}\sqrt{3})^4\).}
{\( |z|=16\); \(\varphi = \frac{2}{3}\pi\)}{\( |z|=8\); \(\varphi = \frac{\pi}{3}\)}{\( |z|=256\); \(\varphi = \frac{2}{3}\pi\)}{\( |z|=2\); \(\varphi = \frac{2}{3}\pi\)}{}{}}
\LANGpl{
\Question{
Znajdź wartość bezwzględną i wartość argumentu liczby  \( z=(1-\mathrm{i}\sqrt{3})^4\).}
{\( |z|=16\); \(\varphi = \frac{2}{3}\pi\)}{\( |z|=8\); \(\varphi = \frac{\pi}{3}\)}{\( |z|=256\); \(\varphi = \frac{2}{3}\pi\)}{\( |z|=2\); \(\varphi = \frac{2}{3}\pi\)}{}{}}
\LANGcs{
\Question{
Vypočtěte absolutní hodnotu a hodnotu argumentu komplexního čísla \( z=(1-\mathrm{i}\sqrt{3})^4\).}
{\( |z|=16\); \(\varphi = \frac{2}{3}\pi\)}{\( |z|=8\); \(\varphi = \frac{\pi}{3}\)}{\( |z|=256\); \(\varphi = \frac{2}{3}\pi\)}{\( |z|=2\); \(\varphi = \frac{2}{3}\pi\)}{}{}}
\LANGsk{
\Question{
Vypočítajte absolútnu hodnotu a hodnotu argumentu komplexného čísla \( z=(1-\mathrm{i}\sqrt{3})^4\).}
{\( |z|=16\); \(\varphi = \frac{2}{3}\pi\)}{\( |z|=8\); \(\varphi = \frac{\pi}{3}\)}{\( |z|=256\); \(\varphi = \frac{2}{3}\pi\)}{\( |z|=2\); \(\varphi = \frac{2}{3}\pi\)}{}{}}
\LANGes{
\Question{
Determina el valor absoluto y el valor del argumento de \( z=(1-\mathrm{i}\sqrt{3})^4\).}
{\( |z|=16\); \(\varphi = \frac{2}{3}\pi\)}{\( |z|=8\); \(\varphi = \frac{\pi}{3}\)}{\( |z|=256\); \(\varphi = \frac{2}{3}\pi\)}{\( |z|=2\); \(\varphi = \frac{2}{3}\pi\)}{}{}}
\End

\Data\ID{2010004613}
\Updated{2022-08-25 10:08:30}
\Level{A}
\SubArea{Powers and Roots of Complex Numbers}
\LANGen{
\Question{
Find the absolute value and the value of the argument of \( z=(2+\mathrm{i}\sqrt{12})^5\).}
{\( |z|=1024\); \(\varphi = \frac{5}{3}\pi\)}{\( |z|=512\); \(\varphi = \frac{5}{3}\pi\)}{\( |z|=1024\); \(\varphi = \frac{\pi}{3}\)}{\( |z|=4\); \(\varphi = \frac{5}{3}\pi\)}{}{}}
\LANGpl{
\Question{
Znajdź wartość bezwzględną i wartość argumentu liczby  \( z=(2+\mathrm{i}\sqrt{12})^5\).}
{\( |z|=1024\); \(\varphi = \frac{5}{3}\pi\)}{\( |z|=512\); \(\varphi = \frac{5}{3}\pi\)}{\( |z|=1024\); \(\varphi = \frac{\pi}{3}\)}{\( |z|=4\); \(\varphi = \frac{5}{3}\pi\)}{}{}}
\LANGcs{
\Question{
Vypočtěte absolutní hodnotu a hodnotu argumentu komplexního čísla \( z=(2+\mathrm{i}\sqrt{12})^5\).}
{\( |z|=1024\); \(\varphi = \frac{5}{3}\pi\)}{\( |z|=512\); \(\varphi = \frac{5}{3}\pi\)}{\( |z|=1024\); \(\varphi = \frac{\pi}{3}\)}{\( |z|=4\); \(\varphi = \frac{5}{3}\pi\)}{}{}}
\LANGsk{
\Question{
Vypočítajte absolútnu hodnotu a hodnotu argumentu komplexného čísla \( z=(2+\mathrm{i}\sqrt{12})^5\).}
{\( |z|=1024\); \(\varphi = \frac{5}{3}\pi\)}{\( |z|=512\); \(\varphi = \frac{5}{3}\pi\)}{\( |z|=1024\); \(\varphi = \frac{\pi}{3}\)}{\( |z|=4\); \(\varphi = \frac{5}{3}\pi\)}{}{}}
\LANGes{
\Question{
Determina el valor absoluto y el valor del argumento de \( z=(2+\mathrm{i}\sqrt{12})^5\).}
{\( |z|=1024\); \(\varphi = \frac{5}{3}\pi\)}{\( |z|=512\); \(\varphi = \frac{5}{3}\pi\)}{\( |z|=1024\); \(\varphi = \frac{\pi}{3}\)}{\( |z|=4\); \(\varphi = \frac{5}{3}\pi\)}{}{}}
\End

\Data\ID{2010004614}
\Updated{2022-08-25 10:08:30}
\Level{A}
\SubArea{Powers and Roots of Complex Numbers}
\LANGen{
\Question{
Evaluate \( (1-\mathrm{i})^{100} \) and write the result in algebraic form.}
{\( -2^{50}\)}{\( 2^{50}\)}{\( -2^{50}\mathrm{i} \)}{\( 2^{50}\mathrm{i} \)}{}{}}
\LANGpl{
\Question{
Oblicz \( (1-\mathrm{i})^{100} \) i wynik zapisz w postaci algebraicznej.}
{\( -2^{50}\)}{\( 2^{50}\)}{\( -2^{50}\mathrm{i} \)}{\( 2^{50}\mathrm{i} \)}{}{}}
\LANGcs{
\Question{
Zjednodušte \( (1-\mathrm{i})^{100} \) a zapište výsledek v algebraickém tvaru.}
{\( -2^{50}\)}{\( 2^{50}\)}{\( -2^{50}\mathrm{i} \)}{\( 2^{50}\mathrm{i} \)}{}{}}
\LANGsk{
\Question{
Zjednodušte \( (1-\mathrm{i})^{100} \) a zapíšte výsledok v algebrickom tvare.}
{\( -2^{50}\)}{\( 2^{50}\)}{\( -2^{50}\mathrm{i} \)}{\( 2^{50}\mathrm{i} \)}{}{}}
\LANGes{
\Question{
Simplifica  \( (1-\mathrm{i})^{100} \) y escribe el resultado en forma algebraica.}
{\( -2^{50}\)}{\( 2^{50}\)}{\( -2^{50}\mathrm{i} \)}{\( 2^{50}\mathrm{i} \)}{}{}}
\End

\Data\ID{2010004615}
\Updated{2022-08-25 10:08:30}
\Level{A}
\SubArea{Powers and Roots of Complex Numbers}
\LANGen{
\Question{
Evaluate \( (-2+2\mathrm{i})^{8} \) and write the result in algebraic form.}
{\( 2^{12}\)}{\( 2^{12}\mathrm{i}\)}{\( -2^{12}\mathrm{i} \)}{\( -2^{12} \)}{}{}}
\LANGpl{
\Question{
Oblicz  \( (-2+2\mathrm{i})^{8} \) i wynik zapisz w postaci algebraicznej.}
{\( 2^{12}\)}{\( 2^{12}\mathrm{i}\)}{\( -2^{12}\mathrm{i} \)}{\( -2^{12} \)}{}{}}
\LANGcs{
\Question{
Zjednodušte \( (-2+2\mathrm{i})^{8} \) a zapište výsledek v algebraickém tvaru.}
{\( 2^{12}\)}{\( 2^{12}\mathrm{i}\)}{\( -2^{12}\mathrm{i} \)}{\( -2^{12} \)}{}{}}
\LANGsk{
\Question{
Zjednodušte \( (-2+2\mathrm{i})^{8} \) a zapíšte výsledok v algebrickom tvare.}
{\( 2^{12}\)}{\( 2^{12}\mathrm{i}\)}{\( -2^{12}\mathrm{i} \)}{\( -2^{12} \)}{}{}}
\LANGes{
\Question{
Simplifica \( (-2+2\mathrm{i})^{8} \) y escribe el resultado en forma algebraica.}
{\( 2^{12}\)}{\( 2^{12}\mathrm{i}\)}{\( -2^{12}\mathrm{i} \)}{\( -2^{12} \)}{}{}}
\End

\Data\ID{2010004616}
\Updated{2023-02-21 10:02:18}
\Level{A}
\SubArea{Powers and Roots of Complex Numbers}
\LANGen{
\Question{
Let \( z \in \mathbb{C}\). The value of the argument of \(z^6\) is \(270^{\circ}\) and \(|z|^6=27\). Find \(z\).}
{\( z=\frac{\sqrt{6}}{2}(1+\mathrm{i})\)}{\( z=\frac{\sqrt{6}}{2}(1-\mathrm{i})\)}{\( z=\sqrt{3}\mathrm{i}\)}{\( z=3(\cos 45^{\circ} + \mathrm{i} \sin 45^{\circ})\)}{}{}}
\LANGpl{
\Question{
Niech  \( z \in \mathbb{C}\). Wartość argumentu liczby  \(z^6\) wynosi \(270^{\circ}\)  i  \(|z|^6=27\). Oblicz  \(z\).}
{\( z=\frac{\sqrt{6}}{2}(1+\mathrm{i})\)}{\( z=\frac{\sqrt{6}}{2}(1-\mathrm{i})\)}{\( z=\sqrt{3}\mathrm{i}\)}{\( z=3(\cos 45^{\circ} + \mathrm{i} \sin 45^{\circ})\)}{}{}}
\LANGcs{
\Question{
Vypočtěte komplexní číslo \(z\), je-li hodnota argumentu \(z^6\) rovna \(270^{\circ}\) a \(|z|^6=27\).}
{\( z=\frac{\sqrt{6}}{2}(1+\mathrm{i})\)}{\( z=\frac{\sqrt{6}}{2}(1-\mathrm{i})\)}{\( z=\sqrt{3}\mathrm{i}\)}{\( z=3(\cos 45^{\circ} + \mathrm{i} \sin 45^{\circ})\)}{}{}}
\LANGsk{
\Question{
Vypočítajte komplexné číslo \(z\), ak hodnota argumentu \(z^6\)  je \(270^{\circ}\) a \(|z|^6=27\).}
{\( z=\frac{\sqrt{6}}{2}(1+\mathrm{i})\)}{\( z=\frac{\sqrt{6}}{2}(1-\mathrm{i})\)}{\( z=\sqrt{3}\mathrm{i}\)}{\( z=3(\cos 45^{\circ} + \mathrm{i} \sin 45^{\circ})\)}{}{}}
\LANGes{
\Question{
Sea \( z \in \mathbb{C}\). El valor del argumento de \(z^6\) es \(270^{\circ}\) y  \(|z|^6=27\). Determina \(z\).}
{\( z=\frac{\sqrt{6}}{2}(1+\mathrm{i})\)}{\( z=\frac{\sqrt{6}}{2}(1-\mathrm{i})\)}{\( z=\sqrt{3}\mathrm{i}\)}{\( z=3(\cos 45^{\circ} + \mathrm{i} \sin 45^{\circ})\)}{}{}}
\End

\Data\ID{2010004617}
\Updated{2023-02-21 10:02:04}
\Level{A}
\SubArea{Powers and Roots of Complex Numbers}
\LANGen{
\Question{
Let \( z \in \mathbb{C}\). The value of the argument of \(z^5\) is \(300^{\circ}\) and \(|z|^5=\frac1{32}\). Find \(z\).}
{\( z=\frac{1}{4}(1+\mathrm{i}\sqrt{3})\)}{\( z=\frac{1}{4}(1-\mathrm{i}\sqrt{3})\)}{\( z=-\frac{1}{2}\mathrm{i}\)}{\( z=\frac{1}{2}(\cos 60^{\circ} - \mathrm{i} \sin 60^{\circ})\)}{}{}}
\LANGpl{
\Question{
Niech \( z \in \mathbb{C}\). Wartość argumentu liczby \(z^5\) wynosi \(300^{\circ}\) i \(|z|^5=\frac1{32}\). Oblicz \(z\).}
{\( z=\frac{1}{4}(1+\mathrm{i}\sqrt{3})\)}{\( z=\frac{1}{4}(1-\mathrm{i}\sqrt{3})\)}{\( z=-\frac{1}{2}\mathrm{i}\)}{\( z=\frac{1}{2}(\cos 60^{\circ} - \mathrm{i} \sin 60^{\circ})\)}{}{}}
\LANGcs{
\Question{
Vypočtěte komplexní číslo \(z\), je-li hodnota agrumentu \(z^5\) rovna \(300^{\circ}\) a \(|z|^5=\frac1{32}\).}
{\( z=\frac{1}{4}(1+\mathrm{i}\sqrt{3})\)}{\( z=\frac{1}{4}(1-\mathrm{i}\sqrt{3})\)}{\( z=-\frac{1}{2}\mathrm{i}\)}{\( z=\frac{1}{2}(\cos 60^{\circ} - \mathrm{i} \sin 60^{\circ})\)}{}{}}
\LANGsk{
\Question{
Vypočítajte komplexné číslo \(z\), ak hodnota argumentu \(z^5\) je \(300^{\circ}\) a \(|z|^5=\frac1{32}\).}
{\( z=\frac{1}{4}(1+\mathrm{i}\sqrt{3})\)}{\( z=\frac{1}{4}(1-\mathrm{i}\sqrt{3})\)}{\( z=-\frac{1}{2}\mathrm{i}\)}{\( z=\frac{1}{2}(\cos 60^{\circ} - \mathrm{i} \sin 60^{\circ})\)}{}{}}
\LANGes{
\Question{
Sea \( z \in \mathbb{C}\). El valor del argumento de \(z^5\) es \(300^{\circ}\) y \(|z|^5=\frac1{32}\). Determina \(z\).}
{\( z=\frac{1}{4}(1+\mathrm{i}\sqrt{3})\)}{\( z=\frac{1}{4}(1-\mathrm{i}\sqrt{3})\)}{\( z=-\frac{1}{2}\mathrm{i}\)}{\( z=\frac{1}{2}(\cos 60^{\circ} - \mathrm{i} \sin 60^{\circ})\)}{}{}}
\End

\Data\ID{9000035709}
\Updated{2020-07-18 01:07:41}
\Level{A}
\SubArea{Powers and Roots of Complex Numbers}
\LANGen{
\Question{
Simplify \((1 -\mathrm{i})^{-3}\).}
{\(-\frac{1} 
{4} + \frac{1} 
{4}\mathrm{i}\)}{\(1 + 3\mathrm{i}\)}{\(- 2 - 2\mathrm{i}\)}{\(\frac{1}
{2} + \frac{1} 
{2}\mathrm{i}\)}{}{}}
\LANGpl{
\Question{
Uprość  \((1 -\mathrm{i})^{-3}\).}
{\(-\frac{1} 
{4} + \frac{1} 
{4}\mathrm{i}\)}{\(1 + 3\mathrm{i}\)}{\(- 2 - 2\mathrm{i}\)}{\(\frac{1}
{2} + \frac{1} 
{2}\mathrm{i}\)}{}{}}
\LANGcs{
\Question{
Komplexní číslo \(z=(1 -\mathrm{i})^{-3}\) 
má tvar:}
{\(-\frac{1} 
{4} + \frac{1} 
{4}\mathrm{i}\)}{\(1 + 3\mathrm{i}\)}{\(- 2 - 2\mathrm{i}\)}{\(\frac{1}
{2} + \frac{1} 
{2}\mathrm{i}\)}{}{}}
\LANGsk{
\Question{
Zjednodušte \((1 -\mathrm{i})^{-3}\).}
{\(-\frac{1} 
{4} + \frac{1} 
{4}\mathrm{i}\)}{\(1 + 3\mathrm{i}\)}{\(- 2 - 2\mathrm{i}\)}{\(\frac{1}
{2} + \frac{1} 
{2}\mathrm{i}\)}{}{}}
\LANGes{
\Question{
Simplifica \((1 -\mathrm{i})^{-3}\).}
{\(-\frac{1} 
{4} + \frac{1} 
{4}\mathrm{i}\)}{\(1 + 3\mathrm{i}\)}{\(- 2 - 2\mathrm{i}\)}{\(\frac{1}
{2} + \frac{1} 
{2}\mathrm{i}\)}{}{}}
\End

\Data\ID{9000035808}
\Updated{2021-04-12 08:04:59}
\Level{A}
\SubArea{Powers and Roots of Complex Numbers}
\LANGen{
\Question{
Evaluate \((1 -\mathrm{i})^{10}\).}
{\(- 32\mathrm{i}\)}{\(32\)}{\(32\mathrm{i}\)}{\(- 32\)}{}{}}
\LANGpl{
\Question{
Oblicz \((1 -\mathrm{i})^{10}\).}
{\(- 32\mathrm{i}\)}{\(32\)}{\(32\mathrm{i}\)}{\(- 32\)}{}{}}
\LANGcs{
\Question{
Komplexní číslo \(z=(1 -\mathrm{i})^{10}\) 
se rovná:}
{\(- 32\mathrm{i}\)}{\(32\)}{\(32\mathrm{i}\)}{\(- 32\)}{}{}}
\LANGsk{
\Question{
Určte algebrický tvar komplexného čísla \(z=(1 -\mathrm{i})^{10}\).}
{\(- 32\mathrm{i}\)}{\(32\)}{\(32\mathrm{i}\)}{\(- 32\)}{}{}}
\LANGes{
\Question{
Calcula \((1 -\mathrm{i})^{10}\).}
{\(- 32\mathrm{i}\)}{\(32\)}{\(32\mathrm{i}\)}{\(- 32\)}{}{}}
\End

\Data\ID{9000035809}
\Updated{2020-07-18 01:07:15}
\Level{A}
\SubArea{Powers and Roots of Complex Numbers}
\LANGen{
\Question{
Given the complex number \(z = -1 + \mathrm{i}\),
find the angle in the polar form of the number
\(z^{6}\).}
{\(\frac{\pi }
{2}\)}{\(\frac{3\pi }
{2}\)}{\(\frac{3\pi }
{4}\)}{\(\frac{7\pi }
{4}\)}{}{}}
\LANGpl{
\Question{
Dana jest liczba zespolona  \(z = -1 + \mathrm{i}\),
wyznacz kąt w postaci biegunowej liczby
\(z^{6}\).}
{\(\frac{\pi }
{2}\)}{\(\frac{3\pi }
{2}\)}{\(\frac{3\pi }
{4}\)}{\(\frac{7\pi }
{4}\)}{}{}}
\LANGcs{
\Question{
Je dáno komplexní číslo \(z = -1 + \mathrm{i}\).
Hlavní hodnota argumentu čísla \(z^{6}\) 
je:}
{\(\frac{\pi }
{2}\)}{\(\frac{3\pi }
{2}\)}{\(\frac{3\pi }
{4}\)}{\(\frac{7\pi }
{4}\)}{}{}}
\LANGsk{
\Question{
Je dané komplexné číslo \(z = -1 + \mathrm{i}\).
Hlavná hodnota argumentu čísla 
\(z^{6}\) je:}
{\(\frac{\pi }
{2}\)}{\(\frac{3\pi }
{2}\)}{\(\frac{3\pi }
{4}\)}{\(\frac{7\pi }
{4}\)}{}{}}
\LANGes{
\Question{
Dado el número complejo \(z = -1 + \mathrm{i}\),
determina la forma polar de
\(z^{6}\).}
{\(\frac{\pi }
{2}\)}{\(\frac{3\pi }
{2}\)}{\(\frac{3\pi }
{4}\)}{\(\frac{7\pi }
{4}\)}{}{}}
\End

\Data\ID{9000037403}
\Updated{2022-05-18 05:05:36}
\Level{A}
\SubArea{Powers and Roots of Complex Numbers}
\LANGen{
\Question{
Given \(z = \sqrt{3}\left (\cos \frac{\pi }{3} + \mathrm{i}\sin \frac{\pi }{3}\right )\),
find \(z^{4}\).}
{\(-\frac{9} 
{2} -\frac{9\mathrm{i}\sqrt{3}} 
 {2} \)}{\(\frac{9}
{2} + \frac{9\mathrm{i}\sqrt{3}} 
 {2} \)}{\(\frac{3}
{2}\)}{\(-\frac{3} 
{2}\)}{}{}}
\LANGpl{
\Question{
Podano  \(z = \sqrt{3}\left (\cos \frac{\pi }{3} + \mathrm{i}\sin \frac{\pi }{3}\right )\),
wskaż  \(z^{4}\).}
{\(-\frac{9} 
{2} -\frac{9\mathrm{i}\sqrt{3}} 
 {2} \)}{\(\frac{9}
{2} + \frac{9\mathrm{i}\sqrt{3}} 
 {2} \)}{\(\frac{3}
{2}\)}{\(-\frac{3} 
{2}\)}{}{}}
\LANGcs{
\Question{
Určete \(z^{4}\),
když \(z = \sqrt{3}\left (\cos \frac{\pi }{3} + \mathrm{i}\sin \frac{\pi }{3}\right )\).}
{\(-\frac{9} 
{2} -\frac{9\mathrm{i}\sqrt{3}} 
 {2} \)}{\(\frac{9}
{2} + \frac{9\mathrm{i}\sqrt{3}} 
 {2} \)}{\(\frac{3}
{2}\)}{\(-\frac{3} 
{2}\)}{}{}}
\LANGsk{
\Question{
Je dané \(z = \sqrt{3}\left (\cos \frac{\pi }{3} + \mathrm{i}\sin \frac{\pi }{3}\right )\),
určte \(z^{4}\).}
{\(-\frac{9} 
{2} -\frac{9\mathrm{i}\sqrt{3}} 
 {2} \)}{\(\frac{9}
{2} + \frac{9\mathrm{i}\sqrt{3}} 
 {2} \)}{\(\frac{3}
{2}\)}{\(-\frac{3} 
{2}\)}{}{}}
\LANGes{
\Question{
Dado \(z = \sqrt{3}\left (\cos \frac{\pi }{3} + \mathrm{i}\sin \frac{\pi }{3}\right )\),
calcula \(z^{4}\).}
{\(-\frac{9} 
{2} -\frac{9\mathrm{i}\sqrt{3}} 
 {2} \)}{\(\frac{9}
{2} + \frac{9\mathrm{i}\sqrt{3}} 
 {2} \)}{\(\frac{3}
{2}\)}{\(-\frac{3} 
{2}\)}{}{}}
\End

\Data\ID{9000037404}
\Updated{2021-08-15 08:08:31}
\Level{A}
\SubArea{Powers and Roots of Complex Numbers}
\LANGen{
\Question{
Given \(z = \sqrt{2}\left (\cos \frac{\pi }{3} -\mathrm{i}\sin \frac{\pi }{3}\right )\),
find \(z^{2}\).}
{\(- 1 -\mathrm{i}\sqrt{3}\)}{\(1 + \mathrm{i}\sqrt{3}\)}{\(- 2 -\mathrm{i}\sqrt{2}\)}{\(2 + \mathrm{i}\sqrt{2}\)}{}{}}
\LANGpl{
\Question{
Podano  \(z = \sqrt{2}\left (\cos \frac{\pi }{3} -\mathrm{i}\sin \frac{\pi }{3}\right )\),
wskaż \(z^{2}\).}
{\(- 1 -\mathrm{i}\sqrt{3}\)}{\(1 + \mathrm{i}\sqrt{3}\)}{\(- 2 -\mathrm{i}\sqrt{2}\)}{\(2 + \mathrm{i}\sqrt{2}\)}{}{}}
\LANGcs{
\Question{
Určete \(z^{2}\), víte-li, že \(z = \sqrt{2}\left (\cos \frac{\pi }{3} -\mathrm{i}\sin \frac{\pi }{3}\right )\).}
{\(- 1 -\mathrm{i}\sqrt{3}\)}{\(1 + \mathrm{i}\sqrt{3}\)}{\(- 2 -\mathrm{i}\sqrt{2}\)}{\(2 + \mathrm{i}\sqrt{2}\)}{}{}}
\LANGsk{
\Question{
Je dané \(z = \sqrt{2}\left (\cos \frac{\pi }{3} -\mathrm{i}\sin \frac{\pi }{3}\right )\),
určte \(z^{2}\).}
{\(- 1 -\mathrm{i}\sqrt{3}\)}{\(1 + \mathrm{i}\sqrt{3}\)}{\(- 2 -\mathrm{i}\sqrt{2}\)}{\(2 + \mathrm{i}\sqrt{2}\)}{}{}}
\LANGes{
\Question{
Dado \(z = \sqrt{2}\left (\cos \frac{\pi }{3} -\mathrm{i}\sin \frac{\pi }{3}\right )\),
determina \(z^{2}\).}
{\(- 1 -\mathrm{i}\sqrt{3}\)}{\(1 + \mathrm{i}\sqrt{3}\)}{\(- 2 -\mathrm{i}\sqrt{2}\)}{\(2 + \mathrm{i}\sqrt{2}\)}{}{}}
\End

\Data\ID{9000037405}
\Updated{2023-02-21 10:02:00}
\Level{A}
\SubArea{Powers and Roots of Complex Numbers}
\LANGen{
\Question{
Evaluate \(\left (1 + \mathrm{i}\right )^{7}\).}
{\(8 - 8\mathrm{i}\)}{\(7 - 7\mathrm{i}\)}{\(1\)}{\(-\mathrm{i}\)}{}{}}
\LANGpl{
\Question{
Oblicz \(\left (1 + \mathrm{i}\right )^{7}\).}
{\(8 - 8\mathrm{i}\)}{\(7 - 7\mathrm{i}\)}{\(1\)}{\(-\mathrm{i}\)}{}{}}
\LANGcs{
\Question{
Zjednodušte \(\left (1 + \mathrm{i}\right )^{7}\).}
{\(8 - 8\mathrm{i}\)}{\(7 - 7\mathrm{i}\)}{\(1\)}{\(-\mathrm{i}\)}{}{}}
\LANGsk{
\Question{
Určte hodnotu \(\left (1 + \mathrm{i}\right )^{7}\).}
{\(8 - 8\mathrm{i}\)}{\(7 - 7\mathrm{i}\)}{\(1\)}{\(-\mathrm{i}\)}{}{}}
\LANGes{
\Question{
Calcula \(\left (1 + \mathrm{i}\right )^{7}\).}
{\(8 - 8\mathrm{i}\)}{\(7 - 7\mathrm{i}\)}{\(1\)}{\(-\mathrm{i}\)}{}{}}
\End

\Data\ID{9000037406}
\Updated{2023-02-21 10:02:14}
\Level{A}
\SubArea{Powers and Roots of Complex Numbers}
\LANGen{
\Question{
Evaluate \(\left (\cos \frac{\pi }{4} + \mathrm{i}\sin \frac{\pi }{4}\right )^{40}\).}
{\(1\)}{\(1 + \mathrm{i}\)}{\(\mathrm{i}\)}{\(1 -\mathrm{i}\)}{}{}}
\LANGpl{
\Question{
Oblicz  \(\left (\cos \frac{\pi }{4} + \mathrm{i}\sin \frac{\pi }{4}\right )^{40}\).}
{\(1\)}{\(1 + \mathrm{i}\)}{\(\mathrm{i}\)}{\(1 -\mathrm{i}\)}{}{}}
\LANGcs{
\Question{
Určete algebraický tvar komplexního čísla \(\left (\cos \frac{\pi }{4} + \mathrm{i}\sin \frac{\pi }{4}\right )^{40}\).}
{\(1\)}{\(1 + \mathrm{i}\)}{\(\mathrm{i}\)}{\(1 -\mathrm{i}\)}{}{}}
\LANGsk{
\Question{
Určte algebraický tvar komplexného čísla \(\left (\cos \frac{\pi }{4} + \mathrm{i}\sin \frac{\pi }{4}\right )^{40}\).}
{\(1\)}{\(1 + \mathrm{i}\)}{\(\mathrm{i}\)}{\(1 -\mathrm{i}\)}{}{}}
\LANGes{
\Question{
Calcula \(\left (\cos \frac{\pi }{4} + \mathrm{i}\sin \frac{\pi }{4}\right )^{40}\).}
{\(1\)}{\(1 + \mathrm{i}\)}{\(\mathrm{i}\)}{\(1 -\mathrm{i}\)}{}{}}
\End

\Data\ID{9000037407}
\Updated{2023-02-21 10:02:28}
\Level{A}
\SubArea{Powers and Roots of Complex Numbers}
\LANGen{
\Question{
Evaluate \(\left (\cos \frac{\pi }{2} + \mathrm{i}\sin \frac{\pi }{2}\right )^{13}\) and find the algebraic form of the result.}
{\(\mathrm{i}\)}{\(1 + 2\mathrm{i}\)}{\(1 -\mathrm{i}\)}{\(1\)}{}{}}
\LANGpl{
\Question{
Oblicz \(\left (\cos \frac{\pi }{2} + \mathrm{i}\sin \frac{\pi }{2}\right )^{13}\).}
{\(\mathrm{i}\)}{\(1 + 2\mathrm{i}\)}{\(1 -\mathrm{i}\)}{\(1\)}{}{}}
\LANGcs{
\Question{
Určete algebraický tvar komplexního čísla \(\left (\cos \frac{\pi }{2} + \mathrm{i}\sin \frac{\pi }{2}\right )^{13}\).}
{\(\mathrm{i}\)}{\(1 + 2\mathrm{i}\)}{\(1 -\mathrm{i}\)}{\(1\)}{}{}}
\LANGsk{
\Question{
Určte algebraický tvar komplexného čísla \(\left (\cos \frac{\pi }{2} + \mathrm{i}\sin \frac{\pi }{2}\right )^{13}\).}
{\(\mathrm{i}\)}{\(1 + 2\mathrm{i}\)}{\(1 -\mathrm{i}\)}{\(1\)}{}{}}
\LANGes{
\Question{
Calcula \(\left (\cos \frac{\pi }{2} + \mathrm{i}\sin \frac{\pi }{2}\right )^{13}\).}
{\(\mathrm{i}\)}{\(1 + 2\mathrm{i}\)}{\(1 -\mathrm{i}\)}{\(1\)}{}{}}
\End

\Data\ID{9000037410}
\Updated{2023-02-21 10:02:29}
\Level{A}
\SubArea{Powers and Roots of Complex Numbers}
\LANGen{
\Question{
Evaluate \(\left (1 -\mathrm{i}\right )^{3}\).}
{\(- 2 - 2\mathrm{i}\)}{\(2 + 2\mathrm{i}\)}{\(1 + \mathrm{i}\)}{\(\mathrm{i}\)}{}{}}
\LANGpl{
\Question{
Oblicz \(\left (1 -\mathrm{i}\right )^{3}\).}
{\(- 2 - 2\mathrm{i}\)}{\(2 + 2\mathrm{i}\)}{\(1 + \mathrm{i}\)}{\(\mathrm{i}\)}{}{}}
\LANGcs{
\Question{
Zjednodušte \(\left (1 -\mathrm{i}\right )^{3}\).}
{\(- 2 - 2\mathrm{i}\)}{\(2 + 2\mathrm{i}\)}{\(1 + \mathrm{i}\)}{\(\mathrm{i}\)}{}{}}
\LANGsk{
\Question{
Zjednodušte \(\left (1 -\mathrm{i}\right )^{3}\).}
{\(- 2 - 2\mathrm{i}\)}{\(2 + 2\mathrm{i}\)}{\(1 + \mathrm{i}\)}{\(\mathrm{i}\)}{}{}}
\LANGes{
\Question{
Calcula \(\left (1 -\mathrm{i}\right )^{3}\).}
{\(- 2 - 2\mathrm{i}\)}{\(2 + 2\mathrm{i}\)}{\(1 + \mathrm{i}\)}{\(\mathrm{i}\)}{}{}}
\End

\Data\ID{9000070101}
\Updated{2022-04-23 05:04:14}
\Level{A}
\SubArea{Powers and Roots of Complex Numbers}
\LANGen{
\Question{
Find the algebraic form of the complex number:
\[ 
 \left (\cos \frac{\pi }
{4} + \mathrm{i}\sin \frac{\pi }
{4}\right )^{3}
\]}
{\(-\frac{\sqrt{2}} 
 {2} + \mathrm{i}\frac{\sqrt{2}} 
 {2} \)}{\(-\frac{\sqrt{2}} 
 {2} -\mathrm{i}\frac{\sqrt{2}} 
 {2} \)}{\(\frac{\sqrt{2}}
 {2} -\mathrm{i}\frac{\sqrt{2}} 
 {2} \)}{\(\frac{\sqrt{2}}
 {2} + \mathrm{i}\frac{\sqrt{2}} 
 {2} \)}{}{}}
\LANGpl{
\Question{
Wyrażenie
\[ 
 \left (\cos \frac{\pi }
{4} + \mathrm{i}\sin \frac{\pi }
{4}\right )^{3}
\] jest równe:}
{\(-\frac{\sqrt{2}} 
 {2} + \mathrm{i}\frac{\sqrt{2}} 
 {2} \)}{\(-\frac{\sqrt{2}} 
 {2} -\mathrm{i}\frac{\sqrt{2}} 
 {2} \)}{\(\frac{\sqrt{2}}
 {2} -\mathrm{i}\frac{\sqrt{2}} 
 {2} \)}{\(\frac{\sqrt{2}}
 {2} + \mathrm{i}\frac{\sqrt{2}} 
 {2} \)}{}{}}
\LANGcs{
\Question{
Určete algebraický tvar daného komplexního čísla.
\[\left (\cos \frac{\pi }{4} + \mathrm{i}\sin \frac{\pi }{4}\right )^{3}\]}
{\(-\frac{\sqrt{2}} 
 {2} + \mathrm{i}\frac{\sqrt{2}} 
 {2} \)}{\(-\frac{\sqrt{2}} 
 {2} -\mathrm{i}\frac{\sqrt{2}} 
 {2} \)}{\(\frac{\sqrt{2}}
 {2} -\mathrm{i}\frac{\sqrt{2}} 
 {2} \)}{\(\frac{\sqrt{2}}
 {2} + \mathrm{i}\frac{\sqrt{2}} 
 {2} \)}{}{}}
\LANGsk{
\Question{
Určte algebraický tvar daného komplexného čísla. 
\[ 
 \left (\cos \frac{\pi }
{4} + \mathrm{i}\sin \frac{\pi }
{4}\right )^{3}
\]}
{\(-\frac{\sqrt{2}} 
 {2} + \mathrm{i}\frac{\sqrt{2}} 
 {2} \)}{\(-\frac{\sqrt{2}} 
 {2} -\mathrm{i}\frac{\sqrt{2}} 
 {2} \)}{\(\frac{\sqrt{2}}
 {2} -\mathrm{i}\frac{\sqrt{2}} 
 {2} \)}{\(\frac{\sqrt{2}}
 {2} + \mathrm{i}\frac{\sqrt{2}} 
 {2} \)}{}{}}
\LANGes{
\Question{
Calcula el siguiente número complejo.
\[ 
 \left (\cos \frac{\pi }
{4} + \mathrm{i}\sin \frac{\pi }
{4}\right )^{3}
\]}
{\(-\frac{\sqrt{2}} 
 {2} + \mathrm{i}\frac{\sqrt{2}} 
 {2} \)}{\(-\frac{\sqrt{2}} 
 {2} -\mathrm{i}\frac{\sqrt{2}} 
 {2} \)}{\(\frac{\sqrt{2}}
 {2} -\mathrm{i}\frac{\sqrt{2}} 
 {2} \)}{\(\frac{\sqrt{2}}
 {2} + \mathrm{i}\frac{\sqrt{2}} 
 {2} \)}{}{}}
\End

\Data\ID{9000070102}
\Updated{2020-07-18 01:07:06}
\Level{A}
\SubArea{Powers and Roots of Complex Numbers}
\LANGen{
\Question{
Evaluate the following complex number. 
\[ 
 \left (\cos \frac{\pi }
{3} + \mathrm{i}\sin \frac{\pi }
{3}\right )^{10}
\]}
{\(-\frac{1} 
{2} -\mathrm{i}\frac{\sqrt{3}} 
 {2} \)}{\(-\frac{\sqrt{3}} 
 {2} -\frac{1} 
{2}\mathrm{i}\)}{\(-\frac{\sqrt{3}} 
 {2} + \frac{1} 
{2}\mathrm{i}\)}{\(-\frac{1} 
{2} + \mathrm{i}\frac{\sqrt{3}} 
 {2} \)}{}{}}
\LANGpl{
\Question{
Wyrażenie 
\[ 
 \left (\cos \frac{\pi }
{3} + \mathrm{i}\sin \frac{\pi }
{3}\right )^{10}
\] jest równe:}
{\(-\frac{1} 
{2} -\mathrm{i}\frac{\sqrt{3}} 
 {2} \)}{\(-\frac{\sqrt{3}} 
 {2} -\frac{1} 
{2}\mathrm{i}\)}{\(-\frac{\sqrt{3}} 
 {2} + \frac{1} 
{2}\mathrm{i}\)}{\(-\frac{1} 
{2} + \mathrm{i}\frac{\sqrt{3}} 
 {2} \)}{}{}}
\LANGcs{
\Question{
Algebraický tvar komplexního čísla \(\left (\cos \frac{\pi }{3} + \mathrm{i}\sin \frac{\pi }{3}\right )^{10}\) 
je roven:}
{\(-\frac{1} 
{2} -\mathrm{i}\frac{\sqrt{3}} 
 {2} \)}{\(-\frac{\sqrt{3}} 
 {2} -\frac{1} 
{2}\mathrm{i}\)}{\(-\frac{\sqrt{3}} 
 {2} + \frac{1} 
{2}\mathrm{i}\)}{\(-\frac{1} 
{2} + \mathrm{i}\frac{\sqrt{3}} 
 {2} \)}{}{}}
\LANGsk{
\Question{
Určte algebraický tvar daného komplexného čísla. 
\[ 
 \left (\cos \frac{\pi }
{3} + \mathrm{i}\sin \frac{\pi }
{3}\right )^{10}
\]}
{\(-\frac{1} 
{2} -\mathrm{i}\frac{\sqrt{3}} 
 {2} \)}{\(-\frac{\sqrt{3}} 
 {2} -\frac{1} 
{2}\mathrm{i}\)}{\(-\frac{\sqrt{3}} 
 {2} + \frac{1} 
{2}\mathrm{i}\)}{\(-\frac{1} 
{2} + \mathrm{i}\frac{\sqrt{3}} 
 {2} \)}{}{}}
\LANGes{
\Question{
Calcula el siguiente número complejo.
\[ 
 \left (\cos \frac{\pi }
{3} + \mathrm{i}\sin \frac{\pi }
{3}\right )^{10}
\]}
{\(-\frac{1} 
{2} -\mathrm{i}\frac{\sqrt{3}} 
 {2} \)}{\(-\frac{\sqrt{3}} 
 {2} -\frac{1} 
{2}\mathrm{i}\)}{\(-\frac{\sqrt{3}} 
 {2} + \frac{1} 
{2}\mathrm{i}\)}{\(-\frac{1} 
{2} + \mathrm{i}\frac{\sqrt{3}} 
 {2} \)}{}{}}
\End

\Data\ID{9000070103}
\Updated{2020-07-18 01:07:16}
\Level{A}
\SubArea{Powers and Roots of Complex Numbers}
\LANGen{
\Question{
Evaluate the following complex number. 
\[ 
 \left (\cos \pi + \mathrm{i}\sin \pi \right )^{9}
\]}
{\(- 1\)}{\(1\)}{\(\mathrm{i}\)}{\(-\mathrm{i}\)}{}{}}
\LANGpl{
\Question{
Wyrażenie  
\[ 
 \left (\cos \pi + \mathrm{i}\sin \pi \right )^{9}
\] jest równe:}
{\(- 1\)}{\(1\)}{\(\mathrm{i}\)}{\(-\mathrm{i}\)}{}{}}
\LANGcs{
\Question{
Určete algebraický tvar daného komplexního čísla.
\[\left (\cos \pi + \mathrm{i}\sin \pi \right )^{9}\]}
{\(- 1\)}{\(1\)}{\(\mathrm{i}\)}{\(-\mathrm{i}\)}{}{}}
\LANGsk{
\Question{
Určte algebraický tvar daného komplexného čísla. 
\[ 
 \left (\cos \pi + \mathrm{i}\sin \pi \right )^{9}
\]}
{\(- 1\)}{\(1\)}{\(\mathrm{i}\)}{\(-\mathrm{i}\)}{}{}}
\LANGes{
\Question{
Calcula el siguiente número complejo.
\[ 
 \left (\cos \pi + \mathrm{i}\sin \pi \right )^{9}
\]}
{\(- 1\)}{\(1\)}{\(\mathrm{i}\)}{\(-\mathrm{i}\)}{}{}}
\End

\Data\ID{9000070104}
\Updated{2020-07-18 01:07:28}
\Level{A}
\SubArea{Powers and Roots of Complex Numbers}
\LANGen{
\Question{
Evaluate the following complex number. 
\[ 
 \left (\sin 2\pi + \mathrm{i}\cos 2\pi \right )^{11}
\]}
{\(-\mathrm{i}\)}{\(- 1\)}{\(1\)}{\(\mathrm{i}\)}{}{}}
\LANGpl{
\Question{
Wyrażenie 
\[ 
 \left (\sin 2\pi + \mathrm{i}\cos 2\pi \right )^{11}
\] jest równe:}
{\(-\mathrm{i}\)}{\(- 1\)}{\(1\)}{\(\mathrm{i}\)}{}{}}
\LANGcs{
\Question{
Určete algebraický tvar daného komplexního čísla.
 \[\left (\sin 2\pi + \mathrm{i}\cos 2\pi \right )^{11}\]}
{\(-\mathrm{i}\)}{\(- 1\)}{\(1\)}{\(\mathrm{i}\)}{}{}}
\LANGsk{
\Question{
Určte algebraický tvar daného komplexného čísla. 
\[ 
 \left (\sin 2\pi + \mathrm{i}\cos 2\pi \right )^{11}
\]}
{\(-\mathrm{i}\)}{\(- 1\)}{\(1\)}{\(\mathrm{i}\)}{}{}}
\LANGes{
\Question{
Calcula el siguiente número complejo.
\[ 
 \left (\sin 2\pi + \mathrm{i}\cos 2\pi \right )^{11}
\]}
{\(-\mathrm{i}\)}{\(- 1\)}{\(1\)}{\(\mathrm{i}\)}{}{}}
\End

\Data\ID{9000070105}
\Updated{2020-07-18 01:07:44}
\Level{A}
\SubArea{Powers and Roots of Complex Numbers}
\LANGen{
\Question{
Evaluate the following complex number. 
\[ 
 \mathrm{i}^{13}
\]}
{\(\cos \frac{\pi }
{2} + \mathrm{i}\sin \frac{\pi }
{2}\)}{\(\cos \frac{\pi }
{2} -\mathrm{i}\sin \frac{\pi }
{2}\)}{\(\sin \frac{\pi }
{2} + \mathrm{i}\cos \frac{\pi }
{2}\)}{\(\cos \frac{3\pi }
{2} + \mathrm{i}\sin \frac{3\pi }
{2}\)}{}{}}
\LANGpl{
\Question{
Wyrażenie 
\[ 
 \mathrm{i}^{13}
\] jest równe:}
{\(\cos \frac{\pi }
{2} + \mathrm{i}\sin \frac{\pi }
{2}\)}{\(\cos \frac{\pi }
{2} -\mathrm{i}\sin \frac{\pi }
{2}\)}{\(\sin \frac{\pi }
{2} + \mathrm{i}\cos \frac{\pi }
{2}\)}{\(\cos \frac{3\pi }
{2} + \mathrm{i}\sin \frac{3\pi }
{2}\)}{}{}}
\LANGcs{
\Question{
Goniometrický tvar komplexního čísla \(\mathrm{i}^{13}\) 
je roven:}
{\(\cos \frac{\pi }{2} + \mathrm{i}\sin \frac{\pi }{2}\)}{\(\cos \frac{\pi }
{2} -\mathrm{i}\sin \frac{\pi }
{2}\)}{\(\sin \frac{\pi }
{2} + \mathrm{i}\cos \frac{\pi }
{2}\)}{\(\cos \frac{3\pi }
{2} + \mathrm{i}\sin \frac{3\pi }
{2}\)}{}{}}
\LANGsk{
\Question{
Určte goniometrický tvar daného komplexného čísla. 
\[ 
 \mathrm{i}^{13}
\]}
{\(\cos \frac{\pi }
{2} + \mathrm{i}\sin \frac{\pi }
{2}\)}{\(\cos \frac{\pi }
{2} -\mathrm{i}\sin \frac{\pi }
{2}\)}{\(\sin \frac{\pi }
{2} + \mathrm{i}\cos \frac{\pi }
{2}\)}{\(\cos \frac{3\pi }
{2} + \mathrm{i}\sin \frac{3\pi }
{2}\)}{}{}}
\LANGes{
\Question{
Calcula el siguiente número complejo.
\[ 
 \mathrm{i}^{13}
\]}
{\(\cos \frac{\pi }
{2} + \mathrm{i}\sin \frac{\pi }
{2}\)}{\(\cos \frac{\pi }
{2} -\mathrm{i}\sin \frac{\pi }
{2}\)}{\(\sin \frac{\pi }
{2} + \mathrm{i}\cos \frac{\pi }
{2}\)}{\(\cos \frac{3\pi }
{2} + \mathrm{i}\sin \frac{3\pi }
{2}\)}{}{}}
\End

\Data\ID{9000070106}
\Updated{2023-02-21 10:02:04}
\Level{A}
\SubArea{Powers and Roots of Complex Numbers}
\LANGen{
\Question{
Evaluate the following complex number. 
\[ 
 (1 -\mathrm{i})^{8}
\]}
{\(16\)}{\(- 16\mathrm{i}\)}{\(16\mathrm{i}\)}{\(- 16\)}{}{}}
\LANGpl{
\Question{
Wyrażenie 
\[ 
 (1 -\mathrm{i})^{8}
\] jest równe:}
{\(16\)}{\(- 16\mathrm{i}\)}{\(16\mathrm{i}\)}{\(- 16\)}{}{}}
\LANGcs{
\Question{
Určete algebraický tvar daného komplexního čísla. 
\[ 
 (1 -\mathrm{i})^{8}
\]}
{\(16\)}{\(- 16\mathrm{i}\)}{\(16\mathrm{i}\)}{\(- 16\)}{}{}}
\LANGsk{
\Question{
Určte algebraický tvar daného komplexného čísla. 
\[ 
 (1 -\mathrm{i})^{8}
\]}
{\(16\)}{\(- 16\mathrm{i}\)}{\(16\mathrm{i}\)}{\(- 16\)}{}{}}
\LANGes{
\Question{
Calcula el siguiente número complejo.
\[ 
 (1 -\mathrm{i})^{8}
\]}
{\(16\)}{\(- 16\mathrm{i}\)}{\(16\mathrm{i}\)}{\(- 16\)}{}{}}
\End

\Data\ID{9000070107}
\Updated{2021-06-30 11:06:28}
\Level{A}
\SubArea{Powers and Roots of Complex Numbers}
\LANGen{
\Question{
Find the algebraic form of the complex number:
\[ 
 \left (\frac{1}
{2} +\cos \frac{\pi } 
{3} + \mathrm{i}\cos 2\pi \right )^{5}
\]}
{\(- 4 - 4\mathrm{i}\)}{\(- 4 + 4\mathrm{i}\)}{\(4 - 4\mathrm{i}\)}{\(4 + 4\mathrm{i}\)}{}{}}
\LANGpl{
\Question{
Wyrażenie 
\[ 
 \left (\frac{1}
{2} +\cos \frac{\pi } 
{3} + \mathrm{i}\cos 2\pi \right )^{5}
\] jest równe:}
{\(- 4 - 4\mathrm{i}\)}{\(- 4 + 4\mathrm{i}\)}{\(4 - 4\mathrm{i}\)}{\(4 + 4\mathrm{i}\)}{}{}}
\LANGcs{
\Question{
Určete algebraický tvar daného komplexního čísla. 
\[\left (\frac{1}
{2} +\cos \frac{\pi } 
{3} + \mathrm{i}\cos 2\pi \right )^{5}\]}
{\(- 4 - 4\mathrm{i}\)}{\(- 4 + 4\mathrm{i}\)}{\(4 - 4\mathrm{i}\)}{\(4 + 4\mathrm{i}\)}{}{}}
\LANGsk{
\Question{
Určte algebraický tvar daného komplexného čísla. 
\[ 
 \left (\frac{1}
{2} +\cos \frac{\pi } 
{3} + \mathrm{i}\cos 2\pi \right )^{5}
\]}
{\(- 4 - 4\mathrm{i}\)}{\(- 4 + 4\mathrm{i}\)}{\(4 - 4\mathrm{i}\)}{\(4 + 4\mathrm{i}\)}{}{}}
\LANGes{
\Question{
Calcula el siguiente número complejo.
\[ 
 \left (\frac{1}
{2} +\cos \frac{\pi } 
{3} + \mathrm{i}\cos 2\pi \right )^{5}
\]}
{\(- 4 - 4\mathrm{i}\)}{\(- 4 + 4\mathrm{i}\)}{\(4 - 4\mathrm{i}\)}{\(4 + 4\mathrm{i}\)}{}{}}
\End

\Data\ID{9000070108}
\Updated{2023-02-21 10:02:54}
\Level{A}
\SubArea{Powers and Roots of Complex Numbers}
\LANGen{
\Question{
Evaluate the following complex number. 
\[ 
 \left (\frac{1}
{2} + \frac{\sqrt{3}} 
 {2} \mathrm{i}\right )^{6}
\]}
{\(1\)}{\(- 1\)}{\(\mathrm{i}\)}{\(-\mathrm{i}\)}{}{}}
\LANGpl{
\Question{
Wyrażenie 
\[ 
 \left (\frac{1}
{2} + \frac{\sqrt{3}} 
 {2} \mathrm{i}\right )^{6}
\] jest równe:}
{\(1\)}{\(- 1\)}{\(\mathrm{i}\)}{\(-\mathrm{i}\)}{}{}}
\LANGcs{
\Question{
Určete algebraický tvar daného komplexního čísla. 
\[ 
 \left (\frac{1}
{2} + \frac{\sqrt{3}} 
 {2} \mathrm{i}\right )^{6}
\]}
{\(1\)}{\(- 1\)}{\(\mathrm{i}\)}{\(-\mathrm{i}\)}{}{}}
\LANGsk{
\Question{
Určte algebraický tvar daného komplexného čísla. 
\[ 
 \left (\frac{1}
{2} + \frac{\sqrt{3}} 
 {2} \mathrm{i}\right )^{6}
\]}
{\(1\)}{\(- 1\)}{\(\mathrm{i}\)}{\(-\mathrm{i}\)}{}{}}
\LANGes{
\Question{
Calcula el siguiente número complejo.
\[ 
 \left (\frac{1}
{2} + \frac{\sqrt{3}} 
 {2} \mathrm{i}\right )^{6}
\]}
{\(1\)}{\(- 1\)}{\(\mathrm{i}\)}{\(-\mathrm{i}\)}{}{}}
\End

\Data\ID{9000070109}
\Updated{2023-02-21 10:02:44}
\Level{A}
\SubArea{Powers and Roots of Complex Numbers}
\LANGen{
\Question{
Evaluate the following complex number. 
\[ 
 \left (\sqrt{3} -\mathrm{i}\right )^{3}
\]}
{\(- 8\mathrm{i}\)}{\(8\)}{\(- 8\)}{\(8\mathrm{i}\)}{}{}}
\LANGpl{
\Question{
Wyrażenie 
\[ 
 \left (\sqrt{3} -\mathrm{i}\right )^{3}
\] jest równe:}
{\(- 8\mathrm{i}\)}{\(8\)}{\(- 8\)}{\(8\mathrm{i}\)}{}{}}
\LANGcs{
\Question{
Určete algebrický tvar daného komplexního čísla. 
\[ 
 \left (\sqrt{3} -\mathrm{i}\right )^{3}
\]}
{\(- 8\mathrm{i}\)}{\(8\)}{\(- 8\)}{\(8\mathrm{i}\)}{}{}}
\LANGsk{
\Question{
Určte algebrický tvar daného komplexného čísla. 
\[ 
 \left (\sqrt{3} -\mathrm{i}\right )^{3}
\]}
{\(- 8\mathrm{i}\)}{\(8\)}{\(- 8\)}{\(8\mathrm{i}\)}{}{}}
\LANGes{
\Question{
Calcula el siguiente número complejo.
\[ 
 \left (\sqrt{3} -\mathrm{i}\right )^{3}
\]}
{\(- 8\mathrm{i}\)}{\(8\)}{\(- 8\)}{\(8\mathrm{i}\)}{}{}}
\End

\Data\ID{1003109301}
\Updated{2023-02-21 11:02:34}
\Level{B}
\SubArea{Powers and Roots of Complex Numbers}
\LANGen{
\Question{
Consider the equation \( x^4 + 81 = 0 \). Find the sum of all its roots in the set of complex numbers.}
{\( 0 \)}{\( 3 \)}{\( -3 \)}{\( 81 \)}{\( - 81 \)}{}}
\LANGpl{
\Question{
Dane jest równanie \( x^4 + 81 = 0 \). Oblicz sumę wszystkich jego pierwiastków w zbiorze liczb zespolonych.}
{\( 0 \)}{\( 3 \)}{\( -3 \)}{\( 81 \)}{\( - 81 \)}{}}
\LANGcs{
\Question{
Součet kořenů binomické rovnice \( x^4 + 81 = 0 \) je:}
{\( 0 \)}{\( 3 \)}{\( -3 \)}{\( 81 \)}{\( - 81 \)}{}}
\LANGsk{
\Question{
Súčet koreňov binomickej rovnice \( x^4 + 81 = 0 \) je:}
{\( 0 \)}{\( 3 \)}{\( -3 \)}{\( 81 \)}{\( - 81 \)}{}}
\LANGes{
\Question{
Considera la ecuación \( x^4 + 81 = 0 \). Determina la suma de todas sus raíces en el conjunto de los números complejos.}
{\( 0 \)}{\( 3 \)}{\( -3 \)}{\( 81 \)}{\( - 81 \)}{}}
\End

\Data\ID{1003109302}
\Updated{2023-02-21 11:02:54}
\Level{B}
\SubArea{Powers and Roots of Complex Numbers}
\LANGen{
\Question{
Consider the equation \( x^4 - 81 = 0 \). Find the product of all its roots in the set of complex numbers.}
{\( -81 \)}{\( 81 \)}{\( -3 \)}{\( 3 \)}{\( 0 \)}{}}
\LANGpl{
\Question{
Dane jest równanie  \( x^4 - 81 = 0 \). Oblicz iloczyn wszystkich jego pierwiastków w zbiorze liczb zespolonych.}
{\( -81 \)}{\( 81 \)}{\( -3 \)}{\( 3 \)}{\( 0 \)}{}}
\LANGcs{
\Question{
Součin kořenů binomické rovnice \( x^4 - 81 = 0 \) je:}
{\( -81 \)}{\( 81 \)}{\( -3 \)}{\( 3 \)}{\( 0 \)}{}}
\LANGsk{
\Question{
Súčin koreňov binomickej rovnice \( x^4 - 81 = 0 \) je:}
{\( -81 \)}{\( 81 \)}{\( -3 \)}{\( 3 \)}{\( 0 \)}{}}
\LANGes{
\Question{
Considera la ecuación \( x^4 - 81 = 0 \). Determina el producto de todas sus raíces en el conjunto de los números complejos.}
{\( -81 \)}{\( 81 \)}{\( -3 \)}{\( 3 \)}{\( 0 \)}{}}
\End

\Data\ID{1003109305}
\Updated{2023-02-21 11:02:11}
\Level{B}
\SubArea{Powers and Roots of Complex Numbers}
\LANGen{
\Question{
Solve the following equation in the set of complex numbers. (Solve the equation by substitution.)
\[ (2x + 3)^4 - 256 = 0 \]}
{\( \left\{-\frac72;\frac12;-\frac32\pm2\mathrm{i} \right\} \)}{\( \left\{-\frac72;\frac12;\frac32\pm2\mathrm{i} \right\} \)}{\( \left\{\frac72;-\frac12;\frac32\pm2\mathrm{i} \right\} \)}{\( \left\{\frac72;-\frac12;-\frac32\pm2\mathrm{i} \right\} \)}{}{}}
\LANGpl{
\Question{
Rozwiąż  równanie w zbiorze liczb zespolonych. (Rozwiąż równanie przez podstawienie.)
\[ (2x + 3)^4 - 256 = 0 \]}
{\( \left\{-\frac72;\frac12;-\frac32\pm2\mathrm{i} \right\} \)}{\( \left\{-\frac72;\frac12;\frac32\pm2\mathrm{i} \right\} \)}{\( \left\{\frac72;-\frac12;\frac32\pm2\mathrm{i} \right\} \)}{\( \left\{\frac72;-\frac12;-\frac32\pm2\mathrm{i} \right\} \)}{}{}}
\LANGcs{
\Question{
Najděte množinu komplexních kořenů binomické rovnice \[ (2x + 3)^4 - 256 = 0. \]  (Řešte pomocí substituce.)}
{\( \left\{-\frac72;\frac12;-\frac32\pm2\mathrm{i} \right\} \)}{\( \left\{-\frac72;\frac12;\frac32\pm2\mathrm{i} \right\} \)}{\( \left\{\frac72;-\frac12;\frac32\pm2\mathrm{i} \right\} \)}{\( \left\{\frac72;-\frac12;-\frac32\pm2\mathrm{i} \right\} \)}{}{}}
\LANGsk{
\Question{
Nájdite množinu komplexných koreňov binomickej rovnice \[ (2x + 3)^4 - 256 = 0. \]  (Riešte pomocou substitúcie.)}
{\( \left\{-\frac72;\frac12;-\frac32\pm2\mathrm{i} \right\} \)}{\( \left\{-\frac72;\frac12;\frac32\pm2\mathrm{i} \right\} \)}{\( \left\{\frac72;-\frac12;\frac32\pm2\mathrm{i} \right\} \)}{\( \left\{\frac72;-\frac12;-\frac32\pm2\mathrm{i} \right\} \)}{}{}}
\LANGes{
\Question{
Resuelve la siguiente ecuación en el conjunto de los números complejos. (Resuelve la ecuación por el método de sustitución).
\[ (2x + 3)^4 - 256 = 0 \]}
{\( \left\{-\frac72;\frac12;-\frac32\pm2\mathrm{i} \right\} \)}{\( \left\{-\frac72;\frac12;\frac32\pm2\mathrm{i} \right\} \)}{\( \left\{\frac72;-\frac12;\frac32\pm2\mathrm{i} \right\} \)}{\( \left\{\frac72;-\frac12;-\frac32\pm2\mathrm{i} \right\} \)}{}{}}
\End

\Data\ID{1103109303}
\Updated{2023-02-21 11:02:52}
\Level{B}
\SubArea{Powers and Roots of Complex Numbers}
\IMG{1103109303.pdf}
\LANGen{
\Question{
Consider an equation \( x^n+b=0 \), where \( n \) is a positive integer and \( b \) is a real number. The points that correspond to the roots of the equation are marked in the figure as black points. Find the equation.}
{\( x^8 - 256 = 0 \)}{\( x^8 + 256 = 0 \)}{\( x^4 + 16 = 0 \)}{\( x^4 - 16 = 0 \)}{\( x^6 - 64 = 0 \)}{\( x^6 + 64 = 0 \)}}
\LANGpl{
\Question{
Dane jest równanie \( x^n+b=0 \), gdzie \( n \) jest  dodatnią liczbą całkowitą, a  \( b \) jest liczbą rzeczywistą.  Punkty odpowiadające pierwiastkom równania są oznaczone na rysunku kolorem czarnym. Wskaż  równanie.}
{\( x^8 - 256 = 0 \)}{\( x^8 + 256 = 0 \)}{\( x^4 + 16 = 0 \)}{\( x^4 - 16 = 0 \)}{\( x^6 - 64 = 0 \)}{\( x^6 + 64 = 0 \)}}
\LANGcs{
\Question{
Na obrázku jsou černými body zobrazeny kořeny binomické rovnice \( x^n+b=0 \), kde \( n \) je přirozené číslo a \( b \) je reálné číslo. Určete tuto rovnici.}
{\( x^8 - 256 = 0 \)}{\( x^8 + 256 = 0 \)}{\( x^4 + 16 = 0 \)}{\( x^4 - 16 = 0 \)}{\( x^6 - 64 = 0 \)}{\( x^6 + 64 = 0 \)}}
\LANGsk{
\Question{
Na obrázku sú čiernymi bodmi zobrazené korene binomickej rovnice \( x^n+b=0 \), kde \( n \) je prirodzené číslo a \( b \) je reálne číslo. Určte túto rovnicu.}
{\( x^8 - 256 = 0 \)}{\( x^8 + 256 = 0 \)}{\( x^4 + 16 = 0 \)}{\( x^4 - 16 = 0 \)}{\( x^6 - 64 = 0 \)}{\( x^6 + 64 = 0 \)}}
\LANGes{
\Question{
Considera la ecuación \( x^n+b=0 \), donde \( n \) es un número natural y \( b \) es un número real. Los puntos que corresponden a las raíces de la ecuación están marcados en la figura como puntos negros. Determina la ecuación.}
{\( x^8 - 256 = 0 \)}{\( x^8 + 256 = 0 \)}{\( x^4 + 16 = 0 \)}{\( x^4 - 16 = 0 \)}{\( x^6 - 64 = 0 \)}{\( x^6 + 64 = 0 \)}}
\End

\Data\ID{1103109304}
\Updated{2023-02-21 11:02:29}
\Level{B}
\SubArea{Powers and Roots of Complex Numbers}
\LANGen{
\Question{
Choose the figure in which the points marked in black correspond to the roots of the equation  \( x^5 + 1 = 0 \).}
{}{}{}{}{}{}}
\LANGpl{
\Question{
Wskaż figurę, na której punkty zaznaczone na czarno odpowiadają pierwiastkom  równania \( x^5 + 1 = 0 \).}
{}{}{}{}{}{}}
\LANGcs{
\Question{
Na kterém obrázku jsou černými body zobrazeny kořeny rovnice  \( x^5 + 1 = 0 \)?}
{}{}{}{}{}{}}
\LANGsk{
\Question{
Na ktorom obrázku sú čiernymi bodmi zobrazené korene binomickej rovnice  \( x^5 + 1 = 0 \)?}
{}{}{}{}{}{}}
\LANGes{
\Question{
Elige la figura en la que los puntos marcados en negro corresponden a las raíces de la ecuación \( x^5 + 1 = 0 \).}
{}{}{}{}{}{}}

\IMGanswer{1103109304a.pdf}
\IMGanswer{1103109304b.pdf}
\IMGanswer{1103109304c.pdf}
\IMGanswer{1103109304d.pdf}\End

\Data\ID{2000002601}
\Updated{2023-02-21 11:02:29}
\Level{B}
\SubArea{Powers and Roots of Complex Numbers}
\LANGen{
\Question{
Consider the equation \(x^4+16=0\). Which of the given numbers is the solution of this equation?}
{\( 2(\cos{\frac{\pi}{4}}+ i \sin{\frac{\pi}{4}}) \)}{\( 2i\)}{\( -2i\)}{\( 2(\cos{{\pi}}+ i \sin{{\pi}}) \)}{}{}}
\LANGpl{
\Question{
Wskaż rozwiązanie równania \(x^4+16=0\).}
{\( 2(\cos{\frac{\pi}{4}}+ i \sin{\frac{\pi}{4}}) \)}{\( 2i\)}{\( -2i\)}{\( 2(\cos{{\pi}}+ i \sin{{\pi}}) \)}{}{}}
\LANGcs{
\Question{
Mějme rovnici \(x^4+16=0\). Které z daných čísel je řešením této rovnice?}
{\( 2(\cos{\frac{\pi}{4}}+ i \sin{\frac{\pi}{4}}) \)}{\( 2i\)}{\( -2i\)}{\( 2(\cos{{\pi}}+ i \sin{{\pi}}) \)}{}{}}
\LANGsk{
\Question{
Ktoré z daných čísel je riešením rovnice \(x^4+16=0\)?}
{\( 2(\cos{\frac{\pi}{4}}+ i \sin{\frac{\pi}{4}}) \)}{\( 2i\)}{\( -2i\)}{\( 2(\cos{{\pi}}+ i \sin{{\pi}}) \)}{}{}}
\LANGes{
\Question{
Considera la ecuación \(x^4+16=0\). ¿Cuál de los números dados es la solución de esta ecuación?}
{\( 2(\cos{\frac{\pi}{4}}+ i \sin{\frac{\pi}{4}}) \)}{\( 2i\)}{\( -2i\)}{\( 2(\cos{{\pi}}+ i \sin{{\pi}}) \)}{}{}}
\End

\Data\ID{2000002602}
\Updated{2023-02-21 11:02:41}
\Level{B}
\SubArea{Powers and Roots of Complex Numbers}
\LANGen{
\Question{
Consider the equation  \(x^4 =1\), where \(x\) is a complex variable. Which of the following statements is true?}
{The equation has four different complex roots.}{The equation has no real root.}{The equation has two double roots:  \(x_{1,2}=1\) and \(x_{3,4}=-1\).}{The equation has the root \(x=1+i\).}{}{}}
\LANGpl{
\Question{
Dane jest równanie   \(x^4 =1\),  gdzie \(x\) jest zmienną zespoloną. Wybierz  poprawne stwierdzenie.}
{Równanie ma cztery różne rozwiązania.}{Równanie nie ma rozwiązania.}{Równanie ma dwa rozwiązania:  \(x_{1,2}=1\)  i  \(x_{3,4}=-1\).}{Rozwiązaniem równania jest \(x=1+i\).}{}{}}
\LANGcs{
\Question{
Mějme rovnici \(x^4 =1\), kde \(x\) je komplexní proměnná. Které z následujících tvrzení je pravdivé?}
{Rovnice má čtyři různé komplexní kořeny.}{Rovnice nemá reálný kořen.}{Rovnice má dva dvojnásobné kořeny:  \(x_{1,2}=1\) a \(x_{3,4}=-1\).}{Rovnice má kořen \(x=1+i\).}{}{}}
\LANGsk{
\Question{
Daná je rovnica \(x^4 =1\), kde \(x\)  je komplexná premenná. Ktoré z nasledujúcich tvrdení o riešení rovnice je pravdivé?}
{Rovnica má štyri rôzne komplexné korene.}{Rovnica nemá reálny koreň.}{Rovnica má dva dvojnásobné korene: \(x_{1,2}=1\) a \(x_{3,4}=-1\).}{Rovnica má koreň \(x=1+i\).}{}{}}
\LANGes{
\Question{
Considera la ecuación  \(x^4 =1\), donde \(x\) es una variable compleja. ¿Cuál de las siguientes proposiciones es verdadera?}
{La ecuación tiene cuatro raíces complejas diferentes.}{La ecuación no tiene ninguna raíz real.}{La ecuación tiene dos raíces dobles:  \(x_{1,2}=1\) y \(x_{3,4}=-1\).}{La ecuación tiene como raíz \(x=1+i\).}{}{}}
\End

\Data\ID{2000002603}
\Updated{2023-02-21 11:02:55}
\Level{B}
\SubArea{Powers and Roots of Complex Numbers}
\LANGen{
\Question{
One of the roots of the equation \(x^3-8=0\) is \(x_1 = -1-i\sqrt{3}\). Find the sum of all its roots.}
{\( 0\)}{\( -8\)}{\( -2i\sqrt{3} \)}{\(-4\)}{}{}}
\LANGpl{
\Question{
Jednym z pierwiastków równania \(x^3-8=0\)  jest \(x_1 = -1-i\sqrt{3}\). Znajdź sumę wszystkich jego pierwiastków.}
{\( 0\)}{\( -8\)}{\( -2i\sqrt{3} \)}{\(-4\)}{}{}}
\LANGcs{
\Question{
Jedním z kořenů rovnice \(x^3-8=0\) je \(x_1 = -1-i\sqrt{3}\). Určete součet všech kořenů rovnice.}
{\( 0\)}{\( -8\)}{\( -2i\sqrt{3} \)}{\(-4\)}{}{}}
\LANGsk{
\Question{
Jeden z koreňov rovnice \(x^3-8=0\) je \(x_1 = -1-i\sqrt{3}\). 
Určte súčet všetkých koreňov tejto rovnice.}
{\( 0\)}{\( -8\)}{\( -2i\sqrt{3} \)}{\(-4\)}{}{}}
\LANGes{
\Question{
Una de las raíces de la ecuación \(x^3-8=0\) es \(x_1 = -1-i\sqrt{3}\). Calcula la suma de todas sus raíces.}
{\( 0\)}{\( -8\)}{\( -2i\sqrt{3} \)}{\(-4\)}{}{}}
\End

\Data\ID{2000002604}
\Updated{2023-02-21 11:02:09}
\Level{B}
\SubArea{Powers and Roots of Complex Numbers}
\LANGen{
\Question{
Find the solution set of the equation \(x^4+81=0\) if you know that one of its roots is \(\frac{3}{\sqrt{2}}(1+i)\).}
{\( \left\{ \frac{3}{\sqrt{2}}(1+i); -\frac{3}{\sqrt{2}}(1+i); \frac{3}{\sqrt{2}}(1-i);-\frac{3}{\sqrt{2}}(1-i) \right\} \)}{\( \left\{ \frac{3}{\sqrt{2}}(1+i); -\frac{3}{\sqrt{2}}(1+i);3;-3 \right\} \)}{\( \left\{ \frac{3}{\sqrt{2}}(1+i); \frac{3}{\sqrt{2}}(1-i);3i;-3i \right\} \)}{\( \left\{\frac{3}{\sqrt{2}}(1+i);\frac{3}{\sqrt{2}}(1-i) \right\}\)}{}{}}
\LANGpl{
\Question{
Wskaż zbiór rozwiązań równania \(x^4+81=0\) wiedząc, że jednym z jego pierwiastków  jest \(\frac{3}{\sqrt{2}}(1+i)\).}
{\( \left\{ \frac{3}{\sqrt{2}}(1+i); -\frac{3}{\sqrt{2}}(1+i); \frac{3}{\sqrt{2}}(1-i);-\frac{3}{\sqrt{2}}(1-i) \right\} \)}{\( \left\{ \frac{3}{\sqrt{2}}(1+i); -\frac{3}{\sqrt{2}}(1+i);3;-3 \right\} \)}{\( \left\{ \frac{3}{\sqrt{2}}(1+i); \frac{3}{\sqrt{2}}(1-i);3i;-3i \right\} \)}{\( \left\{\frac{3}{\sqrt{2}}(1+i);\frac{3}{\sqrt{2}}(1-i) \right\}\)}{}{}}
\LANGcs{
\Question{
Určete množinu řešení rovnice \(x^4+81=0\), když víte, že jedním z kořenů je \(\frac{3}{\sqrt{2}}(1+i)\).}
{\( \left\{ \frac{3}{\sqrt{2}}(1+i); -\frac{3}{\sqrt{2}}(1+i); \frac{3}{\sqrt{2}}(1-i);-\frac{3}{\sqrt{2}}(1-i) \right\} \)}{\( \left\{ \frac{3}{\sqrt{2}}(1+i); -\frac{3}{\sqrt{2}}(1+i);3;-3 \right\} \)}{\( \left\{ \frac{3}{\sqrt{2}}(1+i); \frac{3}{\sqrt{2}}(1-i);3i;-3i \right\} \)}{\( \left\{\frac{3}{\sqrt{2}}(1+i);\frac{3}{\sqrt{2}}(1-i) \right\}\)}{}{}}
\LANGsk{
\Question{
Určte množinu všetkých riešení rovnice \(x^4+81=0\), ak poznáte jeden z koreňov \(\frac{3}{\sqrt{2}}(1+i)\).}
{\( \left\{ \frac{3}{\sqrt{2}}(1+i); -\frac{3}{\sqrt{2}}(1+i); \frac{3}{\sqrt{2}}(1-i);-\frac{3}{\sqrt{2}}(1-i) \right\} \)}{\( \left\{ \frac{3}{\sqrt{2}}(1+i); -\frac{3}{\sqrt{2}}(1+i);3;-3 \right\} \)}{\( \left\{ \frac{3}{\sqrt{2}}(1+i); \frac{3}{\sqrt{2}}(1-i);3i;-3i \right\} \)}{\( \left\{\frac{3}{\sqrt{2}}(1+i);\frac{3}{\sqrt{2}}(1-i) \right\}\)}{}{}}
\LANGes{
\Question{
Determina el conjunto  solución de la ecuación \(x^4+81=0\) si sabes que una de sus raíces es\(\frac{3}{\sqrt{2}}(1+i)\).}
{\( \left\{ \frac{3}{\sqrt{2}}(1+i); -\frac{3}{\sqrt{2}}(1+i); \frac{3}{\sqrt{2}}(1-i);-\frac{3}{\sqrt{2}}(1-i) \right\} \)}{\( \left\{ \frac{3}{\sqrt{2}}(1+i); -\frac{3}{\sqrt{2}}(1+i);3;-3 \right\} \)}{\( \left\{ \frac{3}{\sqrt{2}}(1+i); \frac{3}{\sqrt{2}}(1-i);3i;-3i \right\} \)}{\( \left\{\frac{3}{\sqrt{2}}(1+i);\frac{3}{\sqrt{2}}(1-i) \right\}\)}{}{}}
\End

\Data\ID{2000002605}
\Updated{2023-02-21 11:02:25}
\Level{B}
\SubArea{Powers and Roots of Complex Numbers}
\LANGen{
\Question{
How many solutions does the equation \(2x^4=32\) have in the set of complex numbers?}
{four}{one}{two}{eight}{}{}}
\LANGpl{
\Question{
Ile rozwiązań ma równanie \(2x^4=32\) w zbiorze liczb zespolonych?}
{cztery}{jedno}{dwa}{osiem}{}{}}
\LANGcs{
\Question{
Kolik řešení má rovnice \(2x^4=32\) v množině komplexních čísel?}
{čtyři}{jedno}{dvě}{osm}{}{}}
\LANGsk{
\Question{
Koľko riešení má rovnica \(2x^4=32\) v množine komplexných čísel?}
{štyri}{jedno}{dve}{osem}{}{}}
\LANGes{
\Question{
¿Cuántas soluciones tiene la ecuación \(2x^4=32\) en el conjunto de números complejos?}
{cuatro}{uno}{dos}{ocho}{}{}}
\End

\Data\ID{2000002606}
\Updated{2023-02-21 11:02:38}
\Level{B}
\SubArea{Powers and Roots of Complex Numbers}
\LANGen{
\Question{
Imagine all the solutions of the equation \(x^6 -64 =0\) shown as points in the complex plane. Find the false statement.}
{Two points lie on the imaginary axis.}{The values of the arguments of any two solutions differ by an integer multiple of  \(\frac{\pi}{3}\).}{All solutions of the equation lie on a circle centered at the origin with a radius of \(2\).}{Two points lie on the real axis.}{}{}}
\LANGpl{
\Question{
Rozważ wszystkie rozwiązania równania \(x^6 -64 =0\)  przedstawione jako punkty na płaszczyźnie zespolonej. Wskaż fałszywe stwierdzenie.}
{Na osi urojonej leżą dwa punkty.}{Dwa rozwiązania różnią się całkowitą wielkokrotnością \(\frac{\pi}{3}\).}{Wszystkie rozwiązania równania leżą na okręgu o promieniu \(2\).}{Na osi rzeczywistej leżą dwa punkty.}{}{}}
\LANGcs{
\Question{
Řešení rovnice \(x^6 -64 =0\) jsou zobrazena jako body komplexní roviny. Vyberte nepravdivý výrok.}
{Dva body leží na imaginární ose.}{Hodnoty argumentů každých dvou řešení se liší o celočíselný násobek \(\frac{\pi}{3}\).}{Všechna řešení rovnice leží na kružnici se středem v počátku a poloměrem \(2\).}{Dva body leží na reálné ose.}{}{}}
\LANGsk{
\Question{
Všetky riešenia rovnice \(x^6 -64 =0\) sú zobrazené ako body komplexnej roviny. Vyberte nepravdivý výrok.}
{Dva body ležia na imaginárnej osi.}{Hodnoty argumentov každých dvoch riešení sa líšia o celočíselný násobok \(\frac{\pi}{3}\).}{Všetky riešenia rovnice ležia na kružnici so stredom v počiatku súradného systému s polomerom \(2\).}{Dva body ležia na reálnej osi.}{}{}}
\LANGes{
\Question{
Imagina que todas las soluciones de la ecuación \(x^6 -64 =0\) se muestran como puntos en el plano complejo. Determina la proposición falsa.}
{Dos puntos se encuentran en el eje imaginario.}{Los valores de los argumentos de dos soluciones cualesquiera  difieren en un múltiplo entero \(\frac{\pi}{3}\).}{Todas las soluciones de la ecuación se encuentran en un círculo centrado en el origen con un radio de \(2\).}{Dos puntos se encuentran en el eje real.}{}{}}
\End

\Data\ID{2000002608}
\Updated{2023-02-21 11:02:05}
\Level{B}
\SubArea{Powers and Roots of Complex Numbers}
\LANGen{
\Question{
Find the right formula for solving the equation \(x^5 +32=0\).}
{\( x_k = \sqrt[5]{|-32|}( \cos\frac{\pi +2k\pi}{5}+ i\sin \frac{\pi +2k\pi}{5})\), \(k=0,1,2,3,4\)}{\( x_k = \sqrt[5]{-32}( \cos\frac{\pi +2k\pi}{5}+ i\sin \frac{\pi +2k\pi}{5})\), \(k=0,1,2,3,4\)}{\( x_k = \sqrt[5]{|-32|}( \cos \frac{\pi +k\pi}{5}+ i\sin \frac{\pi +k\pi}{5})\), \(k=0,1,2,3,4\)}{\( x_k = \sqrt[5]{|-32|}( \cos \frac{\pi +2k\pi}{5}+ \sin \frac{\pi +2k\pi}{5})\), \(k=0,1,2,3,4\)}{}{}}
\LANGpl{
\Question{
Wskaż odpowiedni wzór na rozwiązanie równania \(x^5 +32=0\).}
{\( x_k = \sqrt[5]{|-32|}( \cos\frac{\pi +2k\pi}{5}+ i\sin \frac{\pi +2k\pi}{5})\), \(k=0,1,2,3,4\)}{\( x_k = \sqrt[5]{-32}( \cos\frac{\pi +2k\pi}{5}+ i\sin \frac{\pi +2k\pi}{5})\), \(k=0,1,2,3,4\)}{\( x_k = \sqrt[5]{|-32|}( \cos \frac{\pi +k\pi}{5}+ i\sin \frac{\pi +k\pi}{5})\), \(k=0,1,2,3,4\)}{\( x_k = \sqrt[5]{|-32|}( \cos \frac{\pi +2k\pi}{5}+ \sin \frac{\pi +2k\pi}{5})\), \(k=0,1,2,3,4\)}{}{}}
\LANGcs{
\Question{
Vyberte správný vzorec pro řešení rovnice \(x^5 +32=0\).}
{\( x_k = \sqrt[5]{|-32|}( \cos\frac{\pi +2k\pi}{5}+ i\sin \frac{\pi +2k\pi}{5})\), \(k=0,1,2,3,4\)}{\( x_k = \sqrt[5]{-32}( \cos\frac{\pi +2k\pi}{5}+ i\sin \frac{\pi +2k\pi}{5})\), \(k=0,1,2,3,4\)}{\( x_k = \sqrt[5]{|-32|}( \cos \frac{\pi +k\pi}{5}+ i\sin \frac{\pi +k\pi}{5})\), \(k=0,1,2,3,4\)}{\( x_k = \sqrt[5]{|-32|}( \cos \frac{\pi +2k\pi}{5}+ \sin \frac{\pi +2k\pi}{5})\), \(k=0,1,2,3,4\)}{}{}}
\LANGsk{
\Question{
Vyberte správný vzorec pre riešenie rovnice \(x^5 +32=0\).}
{\( x_k = \sqrt[5]{|-32|}( \cos\frac{\pi +2k\pi}{5}+ i\sin \frac{\pi +2k\pi}{5})\), \(k=0,1,2,3,4\)}{\( x_k = \sqrt[5]{-32}( \cos\frac{\pi +2k\pi}{5}+ i\sin \frac{\pi +2k\pi}{5})\), \(k=0,1,2,3,4\)}{\( x_k = \sqrt[5]{|-32|}( \cos \frac{\pi +k\pi}{5}+ i\sin \frac{\pi +k\pi}{5})\), \(k=0,1,2,3,4\)}{\( x_k = \sqrt[5]{|-32|}( \cos \frac{\pi +2k\pi}{5}+ \sin \frac{\pi +2k\pi}{5})\), \(k=0,1,2,3,4\)}{}{}}
\LANGes{
\Question{
Determina la fórmula correcta para resolver la ecuación \(x^5 +32=0\)}
{\( x_k = \sqrt[5]{|-32|}( \cos\frac{\pi +2k\pi}{5}+ i\sin \frac{\pi +2k\pi}{5})\), \(k=0,1,2,3,4\)}{\( x_k = \sqrt[5]{-32}( \cos\frac{\pi +2k\pi}{5}+ i\sin \frac{\pi +2k\pi}{5})\), \(k=0,1,2,3,4\)}{\( x_k = \sqrt[5]{|-32|}( \cos \frac{\pi +k\pi}{5}+ i\sin \frac{\pi +k\pi}{5})\), \(k=0,1,2,3,4\)}{\( x_k = \sqrt[5]{|-32|}( \cos \frac{\pi +2k\pi}{5}+ \sin \frac{\pi +2k\pi}{5})\), \(k=0,1,2,3,4\)}{}{}}
\End

\Data\ID{2010013401}
\Updated{2023-02-21 11:02:52}
\Level{B}
\SubArea{Powers and Roots of Complex Numbers}
\LANGen{
\Question{
Consider the equation \( x^4 + 16 = 0 \). Find the sum of all its roots in the set of complex numbers.}
{\( 0 \)}{\( 2 \)}{\( -2 \)}{\( 16 \)}{\( -16 \)}{}}
\LANGpl{
\Question{
Rozwiąż równanie w zbiorze liczb zespolonych \( x^4 + 16 = 0 \). Znajdź sumę wszystkich jego pierwiastków.}
{\( 0 \)}{\( 2 \)}{\( -2 \)}{\( 16 \)}{\( -16 \)}{}}
\LANGcs{
\Question{
Určete součet komplexních kořenů rovnice \( x^4 + 16 = 0 \).}
{\( 0 \)}{\( 2 \)}{\( -2 \)}{\( 16 \)}{\( -16 \)}{}}
\LANGsk{
\Question{
Súčet koreňov binomickej rovnice \( x^4 + 16 = 0 \) je:}
{\( 0 \)}{\( 2 \)}{\( -2 \)}{\( 16 \)}{\( -16 \)}{}}
\LANGes{
\Question{
Se considera la ecuación \( x^4 + 16 = 0 \). Halla la suma de todas sus raíces en el conjunto de los números complejos.}
{\( 0 \)}{\( 2 \)}{\( -2 \)}{\( 16 \)}{\( -16 \)}{}}
\End

\Data\ID{2010013402}
\Updated{2023-02-21 11:02:04}
\Level{B}
\SubArea{Powers and Roots of Complex Numbers}
\LANGen{
\Question{
Solve the following equation in the set of complex numbers. (Solve the equation by substitution.)
\[ (3x + 2)^4 - 81 = 0 \]}
{\( \left\{-\frac53;\frac13;-\frac23+\mathrm{i} ;-\frac23-\mathrm{i} \right\} \)}{\( \left\{-\frac53;\frac13;\frac23+\mathrm{i} ;\frac23-\mathrm{i} \right\} \)}{\( \left\{\frac53;-\frac13;-\frac23+\mathrm{i} ;-\frac23-\mathrm{i} \right\} \)}{\( \left\{\frac53;-\frac13;\frac23+\mathrm{i} ;\frac23-\mathrm{i} \right\} \)}{}{}}
\LANGpl{
\Question{
Rozwiąż następujące równanie w zbiorze liczb zespolonych. (Rozwiąż równanie przez podstawienie.)
\[ (3x + 2)^4 - 81 = 0 \]}
{\( \left\{-\frac53;\frac13;-\frac23+\mathrm{i} ;-\frac23-\mathrm{i} \right\} \)}{\( \left\{-\frac53;\frac13;\frac23+\mathrm{i} ;\frac23-\mathrm{i} \right\} \)}{\( \left\{\frac53;-\frac13;-\frac23+\mathrm{i} ;-\frac23-\mathrm{i} \right\} \)}{\( \left\{\frac53;-\frac13;\frac23+\mathrm{i} ;\frac23-\mathrm{i} \right\} \)}{}{}}
\LANGcs{
\Question{
Určete množinu komplexních kořenů dané rovnice. (Použijte substituci.)
\[ (3x + 2)^4 - 81 = 0 \]}
{\( \left\{-\frac53;\frac13;-\frac23+\mathrm{i} ;-\frac23-\mathrm{i} \right\} \)}{\( \left\{-\frac53;\frac13;\frac23+\mathrm{i} ;\frac23-\mathrm{i} \right\} \)}{\( \left\{\frac53;-\frac13;-\frac23+\mathrm{i} ;-\frac23-\mathrm{i} \right\} \)}{\( \left\{\frac53;-\frac13;\frac23+\mathrm{i} ;\frac23-\mathrm{i} \right\} \)}{}{}}
\LANGsk{
\Question{
Nájdite množinu komplexných koreňov binomickej rovnice 
\[ (3x + 2)^4 - 81 = 0. \]
(Riešte pomocou substitúcie.)}
{\( \left\{-\frac53;\frac13;-\frac23+\mathrm{i} ;-\frac23-\mathrm{i} \right\} \)}{\( \left\{-\frac53;\frac13;\frac23+\mathrm{i} ;\frac23-\mathrm{i} \right\} \)}{\( \left\{\frac53;-\frac13;-\frac23+\mathrm{i} ;-\frac23-\mathrm{i} \right\} \)}{\( \left\{\frac53;-\frac13;\frac23+\mathrm{i} ;\frac23-\mathrm{i} \right\} \)}{}{}}
\LANGes{
\Question{
Resuelve la siguiente ecuación en el conjunto de los números complejos. (Utiliza un cambio de variable.)
\[ (3x + 2)^4 - 81 = 0 \]}
{\( \left\{-\frac53;\frac13;-\frac23+\mathrm{i} ;-\frac23-\mathrm{i} \right\} \)}{\( \left\{-\frac53;\frac13;\frac23+\mathrm{i} ;\frac23-\mathrm{i} \right\} \)}{\( \left\{\frac53;-\frac13;-\frac23+\mathrm{i} ;-\frac23-\mathrm{i} \right\} \)}{\( \left\{\frac53;-\frac13;\frac23+\mathrm{i} ;\frac23-\mathrm{i} \right\} \)}{}{}}
\End

\Data\ID{2010013404}
\Updated{2023-02-21 11:02:34}
\Level{B}
\SubArea{Powers and Roots of Complex Numbers}
\LANGen{
\Question{
Find the solution set of the following equation in the set of complex numbers.
\[ 
 x^{3} - 8 = 0
\]}
{\(\left\{2;\ -1 - \mathrm{i}\sqrt{3};\ -1 +\mathrm{i}\sqrt{3}\right\}\)}{\(\left\{2;\ 1 - \mathrm{i}\sqrt{3};\ 1 +\mathrm{i}\sqrt{3}\right\}\)}{\(\left\{2;\ 1 - \mathrm{i}\sqrt{3}\right\}\)}{\(\left\{2;\ - 1 + \mathrm{i}\sqrt{3}\right\}\)}{}{}}
\LANGpl{
\Question{
Znajdź zbiór rozwiązań następującego równania w zbiorze liczb zespolonych.
\[ 
 x^{3} - 8 = 0
\]}
{\(\left\{2;\ -1 - \mathrm{i}\sqrt{3};\ -1 +\mathrm{i}\sqrt{3}\right\}\)}{\(\left\{2;\ 1 - \mathrm{i}\sqrt{3};\ 1 +\mathrm{i}\sqrt{3}\right\}\)}{\(\left\{2;\ 1 - \mathrm{i}\sqrt{3}\right\}\)}{\(\left\{2;\ - 1 + \mathrm{i}\sqrt{3}\right\}\)}{}{}}
\LANGcs{
\Question{
Určete množinu všech komplexních kořenů dané rovnice.
\[ 
 x^{3} - 8 = 0
\]}
{\(\left\{2;\ -1 - \mathrm{i}\sqrt{3};\ -1 +\mathrm{i}\sqrt{3}\right\}\)}{\(\left\{2;\ 1 - \mathrm{i}\sqrt{3};\ 1 +\mathrm{i}\sqrt{3}\right\}\)}{\(\left\{2;\ 1 - \mathrm{i}\sqrt{3}\right\}\)}{\(\left\{2;\ - 1 + \mathrm{i}\sqrt{3}\right\}\)}{}{}}
\LANGsk{
\Question{
Nájdite množinu všetkých riešení danej rovnice v množine komplexných čísel.
\[ 
 x^{3} - 8 = 0
\]}
{\(\left\{2;\ -1 - \mathrm{i}\sqrt{3};\ -1 +\mathrm{i}\sqrt{3}\right\}\)}{\(\left\{2;\ 1 - \mathrm{i}\sqrt{3};\ 1 +\mathrm{i}\sqrt{3}\right\}\)}{\(\left\{2;\ 1 - \mathrm{i}\sqrt{3}\right\}\)}{\(\left\{2;\ - 1 + \mathrm{i}\sqrt{3}\right\}\)}{}{}}
\LANGes{
\Question{
Halla todas las soluciones de la siguiente ecuación en el conjunto de los números complejos.
\[ 
 x^{3} - 8 = 0
\]}
{\(\left\{2;\ -1 - \mathrm{i}\sqrt{3};\ -1 +\mathrm{i}\sqrt{3}\right\}\)}{\(\left\{2;\ 1 - \mathrm{i}\sqrt{3};\ 1 +\mathrm{i}\sqrt{3}\right\}\)}{\(\left\{2;\ 1 - \mathrm{i}\sqrt{3}\right\}\)}{\(\left\{2;\ - 1 + \mathrm{i}\sqrt{3}\right\}\)}{}{}}
\End

\Data\ID{2010013405}
\Updated{2023-02-21 11:02:49}
\Level{B}
\SubArea{Powers and Roots of Complex Numbers}
\LANGen{
\Question{
Find the solution set of the following equation in the set of complex numbers.
\[ 
 x^{3} + 27 = 0
\]}
{\(\left\{-3;\ \frac32 - \mathrm{i}\frac{3\sqrt3} {2} ;\ \frac32 +\mathrm{i}\frac{3\sqrt{3}} {2} \right\}\)}{\(\left\{-3;\ -\frac32 + \mathrm{i}\frac{3\sqrt3} {2} ;\ -\frac32 -\mathrm{i}\frac{3\sqrt3} {2} \right\}\)}{\(\left\{-3;\ \frac32 - \mathrm{i}\frac{\sqrt3} {2} ;\ \frac32 +\mathrm{i}\frac{\sqrt3} {2} \right\}\)}{\(\left\{-3;\ -\frac32 + \mathrm{i}\frac{\sqrt3} {2} ;\ -\frac32 -\mathrm{i}\frac{\sqrt3} {2} \right\}\)}{}{}}
\LANGpl{
\Question{
Znajdź zbiór rozwiązań nasępującego równania w zbiorze liczb zespolonych.
\[ 
 x^{3} + 27 = 0
\]}
{\(\left\{-3;\ \frac32 - \mathrm{i}\frac{3\sqrt3} {2} ;\ \frac32 +\mathrm{i}\frac{3\sqrt{3}} {2} \right\}\)}{\(\left\{-3;\ -\frac32 + \mathrm{i}\frac{3\sqrt3} {2} ;\ -\frac32 -\mathrm{i}\frac{3\sqrt3} {2} \right\}\)}{\(\left\{-3;\ \frac32 - \mathrm{i}\frac{\sqrt3} {2} ;\ \frac32 +\mathrm{i}\frac{\sqrt3} {2} \right\}\)}{\(\left\{-3;\ -\frac32 + \mathrm{i}\frac{\sqrt3} {2} ;\ -\frac32 -\mathrm{i}\frac{\sqrt3} {2} \right\}\)}{}{}}
\LANGcs{
\Question{
Určete množinu všech komplexních kořenů dané rovnice.
\[ 
 x^{3} + 27 = 0
\]}
{\(\left\{-3;\ \frac32 - \mathrm{i}\frac{3\sqrt3} {2} ;\ \frac32 +\mathrm{i}\frac{3\sqrt{3}} {2} \right\}\)}{\(\left\{-3;\ -\frac32 + \mathrm{i}\frac{3\sqrt3} {2} ;\ -\frac32 -\mathrm{i}\frac{3\sqrt3} {2} \right\}\)}{\(\left\{-3;\ \frac32 - \mathrm{i}\frac{\sqrt3} {2} ;\ \frac32 +\mathrm{i}\frac{\sqrt3} {2} \right\}\)}{\(\left\{-3;\ -\frac32 + \mathrm{i}\frac{\sqrt3} {2} ;\ -\frac32 -\mathrm{i}\frac{\sqrt3} {2} \right\}\)}{}{}}
\LANGsk{
\Question{
Nájdite množinu všetkých riešení danej rovnice v množine komplexných čísel.
\[ 
 x^{3} + 27 = 0
\]}
{\(\left\{-3;\ \frac32 - \mathrm{i}\frac{3\sqrt3} {2} ;\ \frac32 +\mathrm{i}\frac{3\sqrt{3}} {2} \right\}\)}{\(\left\{-3;\ -\frac32 + \mathrm{i}\frac{3\sqrt3} {2} ;\ -\frac32 -\mathrm{i}\frac{3\sqrt3} {2} \right\}\)}{\(\left\{-3;\ \frac32 - \mathrm{i}\frac{\sqrt3} {2} ;\ \frac32 +\mathrm{i}\frac{\sqrt3} {2} \right\}\)}{\(\left\{-3;\ -\frac32 + \mathrm{i}\frac{\sqrt3} {2} ;\ -\frac32 -\mathrm{i}\frac{\sqrt3} {2} \right\}\)}{}{}}
\LANGes{
\Question{
Calcula todas las soluciones de la siguiente ecuación en el conjunto de los números complejos.
\[ 
 x^{3} + 27 = 0
\]}
{\(\left\{-3;\ \frac32 - \mathrm{i}\frac{3\sqrt3} {2} ;\ \frac32 +\mathrm{i}\frac{3\sqrt{3}} {2} \right\}\)}{\(\left\{-3;\ -\frac32 + \mathrm{i}\frac{3\sqrt3} {2} ;\ -\frac32 -\mathrm{i}\frac{3\sqrt3} {2} \right\}\)}{\(\left\{-3;\ \frac32 - \mathrm{i}\frac{\sqrt3} {2} ;\ \frac32 +\mathrm{i}\frac{\sqrt3} {2} \right\}\)}{\(\left\{-3;\ -\frac32 + \mathrm{i}\frac{\sqrt3} {2} ;\ -\frac32 -\mathrm{i}\frac{\sqrt3} {2} \right\}\)}{}{}}
\End

\Data\ID{9000034301}
\Updated{2023-02-21 11:02:28}
\Level{B}
\SubArea{Powers and Roots of Complex Numbers}
\LANGen{
\Question{
Find the solution set of the following equation in the set of complex numbers.
\[ 
 x^{3} - 1 = 0
\]}
{\(\{1;\ -\frac{1}
{2} + \mathrm{i}\frac{\sqrt{3}} 
 {2} ;\ -\frac{1}
{2} -\mathrm{i}\frac{\sqrt{3}} 
 {2} \}\)}{\(\{1;\ -1 + \mathrm{i}\sqrt{3};\ -1 -\mathrm{i}\sqrt{3}\}\)}{\(\{1;\ -\frac{1}
{2} + \mathrm{i}\frac{\sqrt{3}} 
 {2} \}\)}{\(\{1;\ -\frac{1}
{2} -\mathrm{i}\frac{\sqrt{3}} 
 {2} \}\)}{}{}}
\LANGpl{
\Question{
Wyznacz zbiór rozwiązań równania w zbiorze liczb zespolonych.
\[ 
 x^{3} - 1 = 0
\]}
{\(\{1;\ -\frac{1}
{2} + \mathrm{i}\frac{\sqrt{3}} 
 {2} ;\ -\frac{1}
{2} -\mathrm{i}\frac{\sqrt{3}} 
 {2} \}\)}{\(\{1;\ -1 + \mathrm{i}\sqrt{3};\ -1 -\mathrm{i}\sqrt{3}\}\)}{\(\{1;\ -\frac{1}
{2} + \mathrm{i}\frac{\sqrt{3}} 
 {2} \}\)}{\(\{1;\ -\frac{1}
{2} -\mathrm{i}\frac{\sqrt{3}} 
 {2} \}\)}{}{}}
\LANGcs{
\Question{
Množinou všech komplexních řešení rovnice
\(x^{3} - 1 = 0\) je:}
{\(\{1;\ -\frac{1}
{2} + \mathrm{i}\frac{\sqrt{3}} 
 {2} ;\ -\frac{1}
{2} -\mathrm{i}\frac{\sqrt{3}} 
 {2} \}\)}{\(\{1;\ -1 + \mathrm{i}\sqrt{3};\ -1 -\mathrm{i}\sqrt{3}\}\)}{\(\{1;\ -\frac{1}
{2} + \mathrm{i}\frac{\sqrt{3}} 
 {2} \}\)}{\(\{1;\ -\frac{1}
{2} -\mathrm{i}\frac{\sqrt{3}} 
 {2} \}\)}{}{}}
\LANGsk{
\Question{
Nájdite množinu všetkých riešení danej rovnice v množine komplexných čísel.
\[ 
 x^{3} - 1 = 0
\]}
{\(\{1;\ -\frac{1}
{2} + \mathrm{i}\frac{\sqrt{3}} 
 {2} ;\ -\frac{1}
{2} -\mathrm{i}\frac{\sqrt{3}} 
 {2} \}\)}{\(\{1;\ -1 + \mathrm{i}\sqrt{3};\ -1 -\mathrm{i}\sqrt{3}\}\)}{\(\{1;\ -\frac{1}
{2} + \mathrm{i}\frac{\sqrt{3}} 
 {2} \}\)}{\(\{1;\ -\frac{1}
{2} -\mathrm{i}\frac{\sqrt{3}} 
 {2} \}\)}{}{}}
\LANGes{
\Question{
Determina el conjunto de soluciones de la siguiente ecuación en el conjunto de los números complejos.
\[ 
 x^{3} - 1 = 0
\]}
{\(\{1;\ -\frac{1}
{2} + \mathrm{i}\frac{\sqrt{3}} 
 {2} ;\ -\frac{1}
{2} -\mathrm{i}\frac{\sqrt{3}} 
 {2} \}\)}{\(\{1;\ -1 + \mathrm{i}\sqrt{3};\ -1 -\mathrm{i}\sqrt{3}\}\)}{\(\{1;\ -\frac{1}
{2} + \mathrm{i}\frac{\sqrt{3}} 
 {2} \}\)}{\(\{1;\ -\frac{1}
{2} -\mathrm{i}\frac{\sqrt{3}} 
 {2} \}\)}{}{}}
\End

\Data\ID{9000034302}
\Updated{2023-02-21 11:02:42}
\Level{B}
\SubArea{Powers and Roots of Complex Numbers}
\LANGen{
\Question{
Find the solution set of the following equation in the set of complex numbers.
\[ 
 x^{3} + 8 = 0
\]}
{\(\{ - 2;\ 1 + \mathrm{i}\sqrt{3};\ 1 -\mathrm{i}\sqrt{3}\}\)}{\(\{ - 2;\ -1 + \mathrm{i}\sqrt{3};\ -1 -\mathrm{i}\sqrt{3}\}\)}{\(\{ - 2;\ \frac{1} 
{2} + \mathrm{i}\frac{\sqrt{3}} 
 {2} ;\ \frac{1} 
{2} -\mathrm{i}\frac{\sqrt{3}} 
 {2} \}\)}{\(\{ - 2;\ -\frac{1}
{2} + \mathrm{i}\frac{\sqrt{3}} 
 {2} ;\ -\frac{1}
{2} -\mathrm{i}\frac{\sqrt{3}} 
 {2} \}\)}{}{}}
\LANGpl{
\Question{
Wyznacz zbiór rozwiązań równania w zbiorze liczb zespolonych.
\[ 
 x^{3} + 8 = 0
\]}
{\(\{ - 2;\ 1 + \mathrm{i}\sqrt{3};\ 1 -\mathrm{i}\sqrt{3}\}\)}{\(\{ - 2;\ -1 + \mathrm{i}\sqrt{3};\ -1 -\mathrm{i}\sqrt{3}\}\)}{\(\{ - 2;\ \frac{1} 
{2} + \mathrm{i}\frac{\sqrt{3}} 
 {2} ;\ \frac{1} 
{2} -\mathrm{i}\frac{\sqrt{3}} 
 {2} \}\)}{\(\{ - 2;\ -\frac{1}
{2} + \mathrm{i}\frac{\sqrt{3}} 
 {2} ;\ -\frac{1}
{2} -\mathrm{i}\frac{\sqrt{3}} 
 {2} \}\)}{}{}}
\LANGcs{
\Question{
Množinou všech komplexních řešení rovnice
\(x^{3} + 8 = 0\) je:}
{\(\{ - 2;\ 1 + \mathrm{i}\sqrt{3};\ 1 -\mathrm{i}\sqrt{3}\}\)}{\(\{ - 2;\ -1 + \mathrm{i}\sqrt{3};\ -1 -\mathrm{i}\sqrt{3}\}\)}{\(\{ - 2;\ \frac{1} 
{2} + \mathrm{i}\frac{\sqrt{3}} 
 {2} ;\ \frac{1} 
{2} -\mathrm{i}\frac{\sqrt{3}} 
 {2} \}\)}{\(\{ - 2;\ -\frac{1}
{2} + \mathrm{i}\frac{\sqrt{3}} 
 {2} ;\ -\frac{1}
{2} -\mathrm{i}\frac{\sqrt{3}} 
 {2} \}\)}{}{}}
\LANGsk{
\Question{
Nájdite množinu všetkých riešení danej rovnice v množine komplexných čísel.
\[ 
 x^{3} + 8 = 0
\]}
{\(\{ - 2;\ 1 + \mathrm{i}\sqrt{3};\ 1 -\mathrm{i}\sqrt{3}\}\)}{\(\{ - 2;\ -1 + \mathrm{i}\sqrt{3};\ -1 -\mathrm{i}\sqrt{3}\}\)}{\(\{ - 2;\ \frac{1} 
{2} + \mathrm{i}\frac{\sqrt{3}} 
 {2} ;\ \frac{1} 
{2} -\mathrm{i}\frac{\sqrt{3}} 
 {2} \}\)}{\(\{ - 2;\ -\frac{1}
{2} + \mathrm{i}\frac{\sqrt{3}} 
 {2} ;\ -\frac{1}
{2} -\mathrm{i}\frac{\sqrt{3}} 
 {2} \}\)}{}{}}
\LANGes{
\Question{
Determina el conjunto de soluciones de la siguiente ecuación en el conjunto de los números complejos.
\[ 
 x^{3} + 8 = 0
\]}
{\(\{ - 2;\ 1 + \mathrm{i}\sqrt{3};\ 1 -\mathrm{i}\sqrt{3}\}\)}{\(\{ - 2;\ -1 + \mathrm{i}\sqrt{3};\ -1 -\mathrm{i}\sqrt{3}\}\)}{\(\{ - 2;\ \frac{1} 
{2} + \mathrm{i}\frac{\sqrt{3}} 
 {2} ;\ \frac{1} 
{2} -\mathrm{i}\frac{\sqrt{3}} 
 {2} \}\)}{\(\{ - 2;\ -\frac{1}
{2} + \mathrm{i}\frac{\sqrt{3}} 
 {2} ;\ -\frac{1}
{2} -\mathrm{i}\frac{\sqrt{3}} 
 {2} \}\)}{}{}}
\End

\Data\ID{9000034304}
\Updated{2023-02-21 11:02:13}
\Level{B}
\SubArea{Powers and Roots of Complex Numbers}
\LANGen{
\Question{
Find the solution set of the following equation in the set of complex numbers.
\[ 
 x^{4} - 1 = 0
\]}
{\(\{1;\ -1;\ \mathrm{i};\ -\mathrm{i}\}\)}{\(\{1 -\mathrm{i};\ -1 -\mathrm{i}\}\)}{\(\{1 + \mathrm{i};\ -1 + \mathrm{i}\}\)}{\(\{1 + \mathrm{i};\ 1 -\mathrm{i};\ -1 + \mathrm{i};\ -1 -\mathrm{i}\}\)}{}{}}
\LANGpl{
\Question{
Wyznacz zbiór rozwiązań równania w zbiorze liczb zespolonych.
\[ 
 x^{4} - 1 = 0
\]}
{\(\{1;\ -1;\ \mathrm{i};\ -\mathrm{i}\}\)}{\(\{1 -\mathrm{i};\ -1 -\mathrm{i}\}\)}{\(\{1 + \mathrm{i};\ -1 + \mathrm{i}\}\)}{\(\{1 + \mathrm{i};\ 1 -\mathrm{i};\ -1 + \mathrm{i};\ -1 -\mathrm{i}\}\)}{}{}}
\LANGcs{
\Question{
Množinou všech komplexních řešení rovnice
\(x^{4} - 1 = 0\) je:}
{\(\{1;\ -1;\ \mathrm{i};\ -\mathrm{i}\}\)}{\(\{1 -\mathrm{i};\ -1 -\mathrm{i}\}\)}{\(\{1 + \mathrm{i};\ -1 + \mathrm{i}\}\)}{\(\{1 + \mathrm{i};\ 1 -\mathrm{i};\ -1 + \mathrm{i};\ -1 -\mathrm{i}\}\)}{}{}}
\LANGsk{
\Question{
Nájdite množinu všetkých riešení danej rovnice v množine komplexných čísel. 
\[ 
 x^{4} - 1 = 0
\]}
{\(\{1;\ -1;\ \mathrm{i};\ -\mathrm{i}\}\)}{\(\{1 -\mathrm{i};\ -1 -\mathrm{i}\}\)}{\(\{1 + \mathrm{i};\ -1 + \mathrm{i}\}\)}{\(\{1 + \mathrm{i};\ 1 -\mathrm{i};\ -1 + \mathrm{i};\ -1 -\mathrm{i}\}\)}{}{}}
\LANGes{
\Question{
Determina el conjunto de soluciones de la siguiente ecuación en el conjunto de los números complejos.
\[ 
 x^{4} - 1 = 0
\]}
{\(\{1;\ -1;\ \mathrm{i};\ -\mathrm{i}\}\)}{\(\{1 -\mathrm{i};\ -1 -\mathrm{i}\}\)}{\(\{1 + \mathrm{i};\ -1 + \mathrm{i}\}\)}{\(\{1 + \mathrm{i};\ 1 -\mathrm{i};\ -1 + \mathrm{i};\ -1 -\mathrm{i}\}\)}{}{}}
\End

\Data\ID{9000034305}
\Updated{2023-02-21 11:02:29}
\Level{B}
\SubArea{Powers and Roots of Complex Numbers}
\LANGen{
\Question{
Solve the following equation in the set of complex numbers. 
\[ 
 x^{4} + 16 = 0
\]}
{\(x_{1, 2} = \sqrt{2}(1\pm \mathrm{i}),\ x_{3, 4} = -\sqrt{2}(1\pm \mathrm{i})\)}{\(x_{1, 2} = 1\pm \mathrm{i},\ x_{3, 4} = -1\pm \mathrm{i}\)}{\(x_{1, 2} = 2(1\pm \mathrm{i}),\ x_{3, 4} = -2(1\pm \mathrm{i})\)}{\(x_{1, 2} = \frac{\sqrt{2}} 
 {2} (1\pm \mathrm{i}),\ x_{3, 4} = -\frac{\sqrt{2}}
 {2} (1\pm \mathrm{i})\)}{}{}}
\LANGpl{
\Question{
Rozwiąż równanie w zbiorze liczb zespolonych.
\[ 
 x^{4} + 16 = 0
\]}
{\(x_{1, 2} = \sqrt{2}(1\pm \mathrm{i}),\ x_{3, 4} = -\sqrt{2}(1\pm \mathrm{i})\)}{\(x_{1, 2} = 1\pm \mathrm{i},\ x_{3, 4} = -1\pm \mathrm{i}\)}{\(x_{1, 2} = 2(1\pm \mathrm{i}),\ x_{3, 4} = -2(1\pm \mathrm{i})\)}{\(x_{1, 2} = \frac{\sqrt{2}} 
 {2} (1\pm \mathrm{i}),\ x_{3, 4} = -\frac{\sqrt{2}}
 {2} (1\pm \mathrm{i})\)}{}{}}
\LANGcs{
\Question{
Která z následujících možností vyjadřuje všechna řešení rovnice
\(x^{4} + 16 = 0\) s neznámou
\(x\in \mathbb{C}\)?}
{\(x_{1, 2} = \sqrt{2}(1\pm \mathrm{i}),\ x_{3, 4} = -\sqrt{2}(1\pm \mathrm{i})\)}{\(x_{1, 2} = 1\pm \mathrm{i},\ x_{3, 4} = -1\pm \mathrm{i}\)}{\(x_{1, 2} = 2(1\pm \mathrm{i}),\ x_{3, 4} = -2(1\pm \mathrm{i})\)}{\(x_{1, 2} = \frac{\sqrt{2}} 
 {2} (1\pm \mathrm{i}),\ x_{3, 4} = -\frac{\sqrt{2}}
 {2} (1\pm \mathrm{i})\)}{}{}}
\LANGsk{
\Question{
Vyriešte danú rovnicu v množine komplexných čísel. 
\[ 
 x^{4} + 16 = 0
\]}
{\(x_{1, 2} = \sqrt{2}(1\pm \mathrm{i}),\ x_{3, 4} = -\sqrt{2}(1\pm \mathrm{i})\)}{\(x_{1, 2} = 1\pm \mathrm{i},\ x_{3, 4} = -1\pm \mathrm{i}\)}{\(x_{1, 2} = 2(1\pm \mathrm{i}),\ x_{3, 4} = -2(1\pm \mathrm{i})\)}{\(x_{1, 2} = \frac{\sqrt{2}} 
 {2} (1\pm \mathrm{i}),\ x_{3, 4} = -\frac{\sqrt{2}}
 {2} (1\pm \mathrm{i})\)}{}{}}
\LANGes{
\Question{
Resuelve la siguiente ecuación en el conjunto de los números complejos.
\[ 
 x^{4} + 16 = 0
\]}
{\(x_{1, 2} = \sqrt{2}(1\pm \mathrm{i}),\ x_{3, 4} = -\sqrt{2}(1\pm \mathrm{i})\)}{\(x_{1, 2} = 1\pm \mathrm{i},\ x_{3, 4} = -1\pm \mathrm{i}\)}{\(x_{1, 2} = 2(1\pm \mathrm{i}),\ x_{3, 4} = -2(1\pm \mathrm{i})\)}{\(x_{1, 2} = \frac{\sqrt{2}} 
 {2} (1\pm \mathrm{i}),\ x_{3, 4} = -\frac{\sqrt{2}}
 {2} (1\pm \mathrm{i})\)}{}{}}
\End

\Data\ID{9000034306}
\Updated{2023-02-21 11:02:42}
\Level{B}
\SubArea{Powers and Roots of Complex Numbers}
\LANGen{
\Question{
Solve the following equation in the set of complex numbers. 
\[ 
 x^{6} - 64 = 0
\]}
{\(x_{1, 2} =\pm 2,\ x_{3, 4} = 1\pm \mathrm{i}\sqrt{3},\ x_{5, 6} = -1\pm \mathrm{i}\sqrt{3}\)}{\(x_{1, 2} =\pm 2,\ x_{3, 4} = \frac{1} 
{2}\pm \mathrm{i}\frac{\sqrt{3}} 
 {2} ,\ x_{5, 6} = -\frac{1}
{2}\pm \mathrm{i}\frac{\sqrt{3}} 
 {2} \)}{\(x_{1, 2} =\pm 4,\ x_{3, 4} = 1\pm \mathrm{i}\sqrt{3},\ x_{5, 6} = -1\pm \mathrm{i}\sqrt{3}\)}{\(x_{1, 2} =\pm 8,\ x_{3, 4} = 2\pm 2\mathrm{i}\sqrt{3},\ x_{5, 6} = -2\pm 2\mathrm{i}\sqrt{3}\)}{}{}}
\LANGpl{
\Question{
Rozwiąż równanie w zbiorze liczb zespolonych.
\[ 
 x^{6} - 64 = 0
\]}
{\(x_{1, 2} =\pm 2,\ x_{3, 4} = 1\pm \mathrm{i}\sqrt{3},\ x_{5, 6} = -1\pm \mathrm{i}\sqrt{3}\)}{\(x_{1, 2} =\pm 2,\ x_{3, 4} = \frac{1} 
{2}\pm \mathrm{i}\frac{\sqrt{3}} 
 {2} ,\ x_{5, 6} = -\frac{1}
{2}\pm \mathrm{i}\frac{\sqrt{3}} 
 {2} \)}{\(x_{1, 2} =\pm 4,\ x_{3, 4} = 1\pm \mathrm{i}\sqrt{3},\ x_{5, 6} = -1\pm \mathrm{i}\sqrt{3}\)}{\(x_{1, 2} =\pm 8,\ x_{3, 4} = 2\pm 2\mathrm{i}\sqrt{3},\ x_{5, 6} = -2\pm 2\mathrm{i}\sqrt{3}\)}{}{}}
\LANGcs{
\Question{
Která z následujících možností vyjadřuje všechna řešení rovnice
\(x^{6} - 64 = 0\) s neznámou
\(x\in \mathbb{C}\)?}
{\(x_{1, 2} =\pm 2,\ x_{3, 4} = 1\pm \mathrm{i}\sqrt{3},\ x_{5, 6} = -1\pm \mathrm{i}\sqrt{3}\)}{\(x_{1, 2} =\pm 2,\ x_{3, 4} = \frac{1} 
{2}\pm \mathrm{i}\frac{\sqrt{3}} 
 {2} ,\ x_{5, 6} = -\frac{1}
{2}\pm \mathrm{i}\frac{\sqrt{3}} 
 {2} \)}{\(x_{1, 2} =\pm 4,\ x_{3, 4} = 1\pm \mathrm{i}\sqrt{3},\ x_{5, 6} = -1\pm \mathrm{i}\sqrt{3}\)}{\(x_{1, 2} =\pm 8,\ x_{3, 4} = 2\pm 2\mathrm{i}\sqrt{3},\ x_{5, 6} = -2\pm 2\mathrm{i}\sqrt{3}\)}{}{}}
\LANGsk{
\Question{
Vyriešte nasledujúcu rovnicu v množine komplexných čísel. 
\[ 
 x^{6} - 64 = 0
\]}
{\(x_{1, 2} =\pm 2,\ x_{3, 4} = 1\pm \mathrm{i}\sqrt{3},\ x_{5, 6} = -1\pm \mathrm{i}\sqrt{3}\)}{\(x_{1, 2} =\pm 2,\ x_{3, 4} = \frac{1} 
{2}\pm \mathrm{i}\frac{\sqrt{3}} 
 {2} ,\ x_{5, 6} = -\frac{1}
{2}\pm \mathrm{i}\frac{\sqrt{3}} 
 {2} \)}{\(x_{1, 2} =\pm 4,\ x_{3, 4} = 1\pm \mathrm{i}\sqrt{3},\ x_{5, 6} = -1\pm \mathrm{i}\sqrt{3}\)}{\(x_{1, 2} =\pm 8,\ x_{3, 4} = 2\pm 2\mathrm{i}\sqrt{3},\ x_{5, 6} = -2\pm 2\mathrm{i}\sqrt{3}\)}{}{}}
\LANGes{
\Question{
Resuelve la siguiente ecuación en el conjunto de los números complejos.
\[ 
 x^{6} - 64 = 0
\]}
{\(x_{1, 2} =\pm 2,\ x_{3, 4} = 1\pm \mathrm{i}\sqrt{3},\ x_{5, 6} = -1\pm \mathrm{i}\sqrt{3}\)}{\(x_{1, 2} =\pm 2,\ x_{3, 4} = \frac{1} 
{2}\pm \mathrm{i}\frac{\sqrt{3}} 
 {2} ,\ x_{5, 6} = -\frac{1}
{2}\pm \mathrm{i}\frac{\sqrt{3}} 
 {2} \)}{\(x_{1, 2} =\pm 4,\ x_{3, 4} = 1\pm \mathrm{i}\sqrt{3},\ x_{5, 6} = -1\pm \mathrm{i}\sqrt{3}\)}{\(x_{1, 2} =\pm 8,\ x_{3, 4} = 2\pm 2\mathrm{i}\sqrt{3},\ x_{5, 6} = -2\pm 2\mathrm{i}\sqrt{3}\)}{}{}}
\End

\Data\ID{1003118401}
\Updated{2023-02-21 11:02:32}
\Level{C}
\SubArea{Powers and Roots of Complex Numbers}
\LANGen{
\Question{
Find the solution set of the equation \( x^3 - 8\mathrm{i} = 0 \) in the set of complex numbers.}
{\( \left\{\sqrt3+\mathrm{i}; -\sqrt3+\mathrm{i};-2\mathrm{i} \right\} \)}{\( \left\{ 2\mathrm{i}; -\sqrt3-\mathrm{i}; \sqrt3-\mathrm{i} \right\} \)}{\( \left\{\frac{\sqrt3}2+\frac12\mathrm{i}; -\frac{\sqrt3}2+\frac12\mathrm{i};-\mathrm{i} \right\} \)}{\( \left\{\mathrm{i};-\frac{\sqrt3}2-\frac12\mathrm{i}; \frac{\sqrt3}2-\frac12\mathrm{i} \right\} \)}{}{}}
\LANGpl{
\Question{
Wskaż zbiór rozwiązań równania  \( x^3 - 8\mathrm{i} = 0 \) w zbiorze liczb zespolonych.}
{\( \left\{\sqrt3+\mathrm{i}; -\sqrt3+\mathrm{i};-2\mathrm{i} \right\} \)}{\( \left\{ 2\mathrm{i}; -\sqrt3-\mathrm{i}; \sqrt3-\mathrm{i} \right\} \)}{\( \left\{\frac{\sqrt3}2+\frac12\mathrm{i}; -\frac{\sqrt3}2+\frac12\mathrm{i};-\mathrm{i} \right\} \)}{\( \left\{\mathrm{i};-\frac{\sqrt3}2-\frac12\mathrm{i}; \frac{\sqrt3}2-\frac12\mathrm{i} \right\} \)}{}{}}
\LANGcs{
\Question{
Množina komplexních kořenů binomické rovnice \( x^3 - 8\mathrm{i} = 0 \) je:}
{\( \left\{\sqrt3+\mathrm{i}; -\sqrt3+\mathrm{i};-2\mathrm{i} \right\} \)}{\( \left\{ 2\mathrm{i}; -\sqrt3-\mathrm{i}; \sqrt3-\mathrm{i} \right\} \)}{\( \left\{\frac{\sqrt3}2+\frac12\mathrm{i}; -\frac{\sqrt3}2+\frac12\mathrm{i};-\mathrm{i} \right\} \)}{\( \left\{\mathrm{i};-\frac{\sqrt3}2-\frac12\mathrm{i}; \frac{\sqrt3}2-\frac12\mathrm{i} \right\} \)}{}{}}
\LANGsk{
\Question{
Množina komplexných koreňov binomickej rovnice \( x^3 - 8\mathrm{i} = 0 \) je:}
{\( \left\{\sqrt3+\mathrm{i}; -\sqrt3+\mathrm{i};-2\mathrm{i} \right\} \)}{\( \left\{ 2\mathrm{i}; -\sqrt3-\mathrm{i}; \sqrt3-\mathrm{i} \right\} \)}{\( \left\{\frac{\sqrt3}2+\frac12\mathrm{i}; -\frac{\sqrt3}2+\frac12\mathrm{i};-\mathrm{i} \right\} \)}{\( \left\{\mathrm{i};-\frac{\sqrt3}2-\frac12\mathrm{i}; \frac{\sqrt3}2-\frac12\mathrm{i} \right\} \)}{}{}}
\LANGes{
\Question{
Determina el conjunto de soluciones de la ecuación \( x^3 - 8\mathrm{i} = 0 \) en el conjunto de los números complejos.}
{\( \left\{\sqrt3+\mathrm{i}; -\sqrt3+\mathrm{i};-2\mathrm{i} \right\} \)}{\( \left\{ 2\mathrm{i}; -\sqrt3-\mathrm{i}; \sqrt3-\mathrm{i} \right\} \)}{\( \left\{\frac{\sqrt3}2+\frac12\mathrm{i}; -\frac{\sqrt3}2+\frac12\mathrm{i};-\mathrm{i} \right\} \)}{\( \left\{\mathrm{i};-\frac{\sqrt3}2-\frac12\mathrm{i}; \frac{\sqrt3}2-\frac12\mathrm{i} \right\} \)}{}{}}
\End

\Data\ID{1003118402}
\Updated{2023-02-21 11:02:52}
\Level{C}
\SubArea{Powers and Roots of Complex Numbers}
\LANGen{
\Question{
Which of the given complex numbers is not a root of the equation \( x^6  +  8\mathrm{i}  =  0 \)?}
{\( 1-\mathrm{i} \)}{\( 1+\mathrm{i} \)}{\( -1-\mathrm{i} \)}{\( \sqrt2\left(\cos\frac{7\pi}{12}+\mathrm{i}\cdot\sin\frac{7\pi}{12}\right) \)}{\( \sqrt2\left(\cos\frac{23\pi}{12}+\mathrm{i}\cdot\sin\frac{23\pi}{12}\right) \)}{}}
\LANGpl{
\Question{
Która z podanych liczb zespolonych nie jest pierwiastkiem równania \( x^6  +  8\mathrm{i}  =  0 \)?}
{\( 1-\mathrm{i} \)}{\( 1+\mathrm{i} \)}{\( -1-\mathrm{i} \)}{\( \sqrt2\left(\cos\frac{7\pi}{12}+\mathrm{i}\cdot\sin\frac{7\pi}{12}\right) \)}{\( \sqrt2\left(\cos\frac{23\pi}{12}+\mathrm{i}\cdot\sin\frac{23\pi}{12}\right) \)}{}}
\LANGcs{
\Question{
Kořenem binomické rovnice \( x^6  +  8\mathrm{i}  =  0 \) není:}
{\( 1-\mathrm{i} \)}{\( 1+\mathrm{i} \)}{\( -1-\mathrm{i} \)}{\( \sqrt2\left(\cos\frac{7\pi}{12}+\mathrm{i}\cdot\sin\frac{7\pi}{12}\right) \)}{\( \sqrt2\left(\cos\frac{23\pi}{12}+\mathrm{i}\cdot\sin\frac{23\pi}{12}\right) \)}{}}
\LANGsk{
\Question{
Koreňom binomickej rovnice \( x^6  +  8\mathrm{i}  =  0 \) nie je:}
{\( 1-\mathrm{i} \)}{\( 1+\mathrm{i} \)}{\( -1-\mathrm{i} \)}{\( \sqrt2\left(\cos\frac{7\pi}{12}+\mathrm{i}\cdot\sin\frac{7\pi}{12}\right) \)}{\( \sqrt2\left(\cos\frac{23\pi}{12}+\mathrm{i}\cdot\sin\frac{23\pi}{12}\right) \)}{}}
\LANGes{
\Question{
¿Cuál de los números complejos dados no es una raíz de la ecuación \( x^6  +  8\mathrm{i}  =  0 \)?}
{\( 1-\mathrm{i} \)}{\( 1+\mathrm{i} \)}{\( -1-\mathrm{i} \)}{\( \sqrt2\left(\cos\frac{7\pi}{12}+\mathrm{i}\cdot\sin\frac{7\pi}{12}\right) \)}{\( \sqrt2\left(\cos\frac{23\pi}{12}+\mathrm{i}\cdot\sin\frac{23\pi}{12}\right) \)}{}}
\End

\Data\ID{1003118405}
\Updated{2023-02-21 11:02:18}
\Level{C}
\SubArea{Powers and Roots of Complex Numbers}
\LANGen{
\Question{
All the solutions of the equation \( x^6-4\sqrt3+4\mathrm{i} = 0 \) can be  displayed as points in the  rectangular coordinate system. What is the distance of the two most distant points?}
{\( 2\sqrt2 \)}{\( \sqrt2 \)}{\( 2\sqrt[3]4 \)}{\( \sqrt[3]4 \)}{\( 2\sqrt3 \)}{\( \sqrt3 \)}}
\LANGpl{
\Question{
Wszystkie rozwiązania równania  \( x^6-4\sqrt3+4\mathrm{i} = 0 \) mogą być zaznaczone  jako punkty w  układzie współrzędnych. Jaka jest odległość dwóch najbardziej odległych punktów?}
{\( 2\sqrt2 \)}{\( \sqrt2 \)}{\( 2\sqrt[3]4 \)}{\( \sqrt[3]4 \)}{\( 2\sqrt3 \)}{\( \sqrt3 \)}}
\LANGcs{
\Question{
Vzdálenost dvou nejvzdálenějších obrazů kořenů binomické rovnice  \( x^6-4\sqrt3+4\mathrm{i} = 0 \), zobrazených v Gaussově rovině, je:}
{\( 2\sqrt2 \)}{\( \sqrt2 \)}{\( 2\sqrt[3]4 \)}{\( \sqrt[3]4 \)}{\( 2\sqrt3 \)}{\( \sqrt3 \)}}
\LANGsk{
\Question{
Vzdialenosť dvoch najvzdialenejších obrazov koreňov binomickej rovnice  \( x^6-4\sqrt3+4\mathrm{i} = 0 \), zobrazených v Gaussovej rovine, je:}
{\( 2\sqrt2 \)}{\( \sqrt2 \)}{\( 2\sqrt[3]4 \)}{\( \sqrt[3]4 \)}{\( 2\sqrt3 \)}{\( \sqrt3 \)}}
\LANGes{
\Question{
Todas las soluciones de la ecuación \( x^6-4\sqrt3+4\mathrm{i} = 0 \) se pueden mostrar como puntos en el sistema de coordenadas rectangular. ¿Cuál es la distancia de los dos puntos más distantes?}
{\( 2\sqrt2 \)}{\( \sqrt2 \)}{\( 2\sqrt[3]4 \)}{\( \sqrt[3]4 \)}{\( 2\sqrt3 \)}{\( \sqrt3 \)}}
\End

\Data\ID{1003118406}
\Updated{2023-02-21 11:02:37}
\Level{C}
\SubArea{Powers and Roots of Complex Numbers}
\LANGen{
\Question{
All the solutions of the equation \( x^4+1+\sqrt3\mathrm{i} = 0 \) are complex numbers, whose arguments are from the interval \( [0; 2\pi) \). Find the sum of the arguments of all the solutions of the equation.}
{\( \frac{13}3\pi \)}{\( 4\pi \)}{\( \frac{25}6\pi \)}{\( \frac92\pi \)}{}{}}
\LANGpl{
\Question{
Wszystkie rozwiązania równania \( x^4+1+\sqrt3\mathrm{i} = 0 \) są liczbami zespolonymi, których argumenty należą  do  przedziału \( \langle0; 2\pi) \). Znajdź sumę argumentów wszystkich rozwiązań równania.}
{\( \frac{13}3\pi \)}{\( 4\pi \)}{\( \frac{25}6\pi \)}{\( \frac92\pi \)}{}{}}
\LANGcs{
\Question{
Všechny kořeny rovnice  \( x^4+1+\sqrt3\mathrm{i} = 0 \) jsou komplexní čísla s argumenty z intervalu \( \langle0; 2\pi) \). Určete součet argumentů všech kořenů rovnice.}
{\( \frac{13}3\pi \)}{\( 4\pi \)}{\( \frac{25}6\pi \)}{\( \frac92\pi \)}{}{}}
\LANGsk{
\Question{
Všetky korene rovnice  \( x^4+1+\sqrt3\mathrm{i} = 0 \) sú komplexné čísla s argumentami z intervalu \( \langle0; 2\pi) \). Určte súčet argumentov všetkých koreňov rovnice.}
{\( \frac{13}3\pi \)}{\( 4\pi \)}{\( \frac{25}6\pi \)}{\( \frac92\pi \)}{}{}}
\LANGes{
\Question{
Todas las soluciones de la ecuación \( x^4+1+\sqrt3\mathrm{i} = 0 \) son números complejos, cuyos argumentos pertenecen al intervalo \( [0; 2\pi) \). Calcula la suma de los argumentos de todas las soluciones de la ecuación.}
{\( \frac{13}3\pi \)}{\( 4\pi \)}{\( \frac{25}6\pi \)}{\( \frac92\pi \)}{}{}}
\End

\Data\ID{1103118403}
\Updated{2023-02-21 11:02:59}
\Level{C}
\SubArea{Powers and Roots of Complex Numbers}
\LANGen{
\Question{
Consider the equation  \( x^4+2-2\sqrt3\mathrm{i} = 0 \).  In one of the following pictures the points depicted in black correspond to the roots of the equation. Find the figure.}
{}{}{}{}{}{}}
\LANGpl{
\Question{
Dane jest równanie  \( x^4+2-2\sqrt3\mathrm{i} = 0 \).  Na jednym z poniższych rysunków punkty oznaczone kolorem czarnym odpowiadają pierwiastkom danego równania. Wskaż figurę.}
{}{}{}{}{}{}}
\LANGcs{
\Question{
Vyberte obrázek, na kterém jsou černými body zobrazeny kořeny binomické rovnice:  \( x^4+2-2\sqrt3\mathrm{i} = 0 \).}
{}{}{}{}{}{}}
\LANGsk{
\Question{
Vyberte obrázok, na ktorom sú čiernymi bodmi zobrazené korene binomickej rovnice:  \( x^4+2-2\sqrt3\mathrm{i} = 0 \).}
{}{}{}{}{}{}}
\LANGes{
\Question{
Considera la ecuación \( x^4+2-2\sqrt3\mathrm{i} = 0 \).  En una de las siguientes figuras, los puntos representados en negro corresponden a las raíces de la ecuación. Determina la figura.}
{}{}{}{}{}{}}

\IMGanswer{1103118403a_0.pdf}
\IMGanswer{1103118403b_0.pdf}
\IMGanswer{1103118403c_0.pdf}
\IMGanswer{1103118403d_0.pdf}\End

\Data\ID{1103118404}
\Updated{2023-02-21 11:02:14}
\Level{C}
\SubArea{Powers and Roots of Complex Numbers}
\IMG{1103118404_0.pdf}
\LANGen{
\Question{
Consider an equation \( x^n+b=0 \), where \( n \) is a positive integer and \( b \) is a complex number. In the picture the points that correspond to the roots of the equation are depicted in black. Find the equation.}
{\( x^3 + 4\sqrt2 - 4\sqrt2\mathrm{i} = 0 \)}{\( x^3 + 4\sqrt2 +4\sqrt2\mathrm{i} = 0 \)}{\( x^3 - 4\sqrt2 - 4\sqrt2\mathrm{i} = 0 \)}{\( x^3 - 4\sqrt2 +4\sqrt2\mathrm{i} = 0 \)}{}{}}
\LANGpl{
\Question{
Dane jest równanie \( x^n+b=0 \), gdzie \( n \) jest dodatnią liczbą naturalna, a  \( b \) jest liczbą zespoloną. Na rysunku punkty oznaczone kolorem czarnym odpowiadają pierwiastkom równania. Wskaż równanie.}
{\( x^3 + 4\sqrt2 - 4\sqrt2\mathrm{i} = 0 \)}{\( x^3 + 4\sqrt2 +4\sqrt2\mathrm{i} = 0 \)}{\( x^3 - 4\sqrt2 - 4\sqrt2\mathrm{i} = 0 \)}{\( x^3 - 4\sqrt2 +4\sqrt2\mathrm{i} = 0 \)}{}{}}
\LANGcs{
\Question{
Uvažujte rovnici $x^n+b=0$, kde $n\in\mathbb{N}^+$ a $b$ je komplexní číslo. Na obrázku jsou černými body zobrazeny kořeny binomické rovnice:}
{\( x^3 + 4\sqrt2 - 4\sqrt2\mathrm{i} = 0 \)}{\( x^3 + 4\sqrt2 +4\sqrt2\mathrm{i} = 0 \)}{\( x^3 - 4\sqrt2 - 4\sqrt2\mathrm{i} = 0 \)}{\( x^3 - 4\sqrt2 +4\sqrt2\mathrm{i} = 0 \)}{}{}}
\LANGsk{
\Question{
Uvažujte rovnicu $x^n+b=0$, kde $n\in\mathbb{N}^+$ a $b$ je komplexné číslo. Na obrázku sú čiernymi bodmi zobrazené korene binomickej rovnice:}
{\( x^3 + 4\sqrt2 - 4\sqrt2\mathrm{i} = 0 \)}{\( x^3 + 4\sqrt2 +4\sqrt2\mathrm{i} = 0 \)}{\( x^3 - 4\sqrt2 - 4\sqrt2\mathrm{i} = 0 \)}{\( x^3 - 4\sqrt2 +4\sqrt2\mathrm{i} = 0 \)}{}{}}
\LANGes{
\Question{
Considera la ecuación \( x^n+b=0 \), donde \( n \) es un número natural y \( b \) es un número complejo.  En la imagen, los puntos que corresponden a las raíces de la ecuación se representan en negro. Determina la ecuación.}
{\( x^3 + 4\sqrt2 - 4\sqrt2\mathrm{i} = 0 \)}{\( x^3 + 4\sqrt2 +4\sqrt2\mathrm{i} = 0 \)}{\( x^3 - 4\sqrt2 - 4\sqrt2\mathrm{i} = 0 \)}{\( x^3 - 4\sqrt2 +4\sqrt2\mathrm{i} = 0 \)}{}{}}
\End

\Data\ID{2000002607}
\Updated{2023-02-21 11:02:50}
\Level{C}
\SubArea{Powers and Roots of Complex Numbers}
\LANGen{
\Question{
Which of the following equations has not the solution \(x=i\)?}
{\( x^6 -1 =0\)}{\( x^6 +1 =0\)}{\( x^3 +i =0\)}{\( x^5 -i=0\)}{}{}}
\LANGpl{
\Question{
Które z podanych równań nie ma rozwiązania \(x=i\)?}
{\( x^6 -1 =0\)}{\( x^6 +1 =0\)}{\( x^3 +i =0\)}{\( x^5 -i=0\)}{}{}}
\LANGcs{
\Question{
Která z následujících rovnic nemá řešení \(x=i\)?}
{\( x^6 -1 =0\)}{\( x^6 +1 =0\)}{\( x^3 +i =0\)}{\( x^5 -i=0\)}{}{}}
\LANGsk{
\Question{
Ktorá z následujúcich rovníc nemá riešenie \(x=i\)?}
{\( x^6 -1 =0\)}{\( x^6 +1 =0\)}{\( x^3 +i =0\)}{\( x^5 -i=0\)}{}{}}
\LANGes{
\Question{
¿Cuál de las siguientes ecuaciones no tiene la solución \(x=i\)?}
{\( x^6 -1 =0\)}{\( x^6 +1 =0\)}{\( x^3 +i =0\)}{\( x^5 -i=0\)}{}{}}
\End

\Data\ID{2000002610}
\Updated{2023-02-21 11:02:27}
\Level{C}
\SubArea{Powers and Roots of Complex Numbers}
\LANGen{
\Question{
One solution of the equation \( x^2-72i=0\) is \(x_1 = 6(1+i)\). Find the missing solution.}
{\( x_2 = -6(1+i) \)}{\( x_2 = 6(1-i) \)}{\( x_2 = 6(-1+i) \)}{\( x_2 = -6(1-i) \)}{}{}}
\LANGpl{
\Question{
Jednym z rozwiązań równania  \( x^2-72i=0\) jest \(x_1 = 6(1+i)\). Znajdź brakujące rozwiązanie.}
{\( x_2 = -6(1+i) \)}{\( x_2 = 6(1-i) \)}{\( x_2 = 6(-1+i) \)}{\( x_2 = -6(1-i) \)}{}{}}
\LANGcs{
\Question{
Jedním z řešení \( x^2-72i=0\) je \(x_1 = 6(1+i)\). Určete chybějící řešení.}
{\( x_2 = -6(1+i) \)}{\( x_2 = 6(1-i) \)}{\( x_2 = 6(-1+i) \)}{\( x_2 = -6(1-i) \)}{}{}}
\LANGsk{
\Question{
Jeden z koreňov rovnice \( x^2-72i=0\) je \(x_1 = 6(1+i)\). Určte chýbajúce riešenie.}
{\( x_2 = -6(1+i) \)}{\( x_2 = 6(1-i) \)}{\( x_2 = 6(-1+i) \)}{\( x_2 = -6(1-i) \)}{}{}}
\LANGes{
\Question{
Una solución de la ecuación \( x^2-72i=0\) es \(x_1 = 6(1+i)\). Determina la solución que falta.}
{\( x_2 = -6(1+i) \)}{\( x_2 = 6(1-i) \)}{\( x_2 = 6(-1+i) \)}{\( x_2 = -6(1-i) \)}{}{}}
\End

\Data\ID{2010013403}
\Updated{2023-02-21 11:02:20}
\Level{C}
\SubArea{Powers and Roots of Complex Numbers}
\LANGen{
\Question{
All the solutions of the equation \( x^6+3\sqrt5-6\mathrm{i} = 0 \) can be  displayed as points in the  rectangular coordinate system. What is the distance of the two most distant points?}
{\( 2\sqrt[3]3 \)}{\( 2\sqrt3 \)}{\( \sqrt3 \)}{\( \sqrt[3]9\)}{\( 2\sqrt[3]9\)}{}}
\LANGpl{
\Question{
Wszystkie rozwiązania równania \( x^6+3\sqrt5-6\mathrm{i} = 0 \) mogą być przedstawione, jako punkty w prostokątnym układzie współrzędnych. Jaka jest odległość dwóch najbardziej odległych punktów?}
{\( 2\sqrt[3]3 \)}{\( 2\sqrt3 \)}{\( \sqrt3 \)}{\( \sqrt[3]9\)}{\( 2\sqrt[3]9\)}{}}
\LANGcs{
\Question{
Všechna řešení rovnice \( x^6+3\sqrt5-6\mathrm{i} = 0 \) mohou být zobrazena v pravúhlém souřadnicovém systému. Jaká je vzdálenost nejvzdálenějších bodů?}
{\( 2\sqrt[3]3 \)}{\( 2\sqrt3 \)}{\( \sqrt3 \)}{\( \sqrt[3]9\)}{\( 2\sqrt[3]9\)}{}}
\LANGsk{
\Question{
Vzdialenosť dvoch najvzdialenejších obrazov koreňov binomickej rovnice \( x^6+3\sqrt5-6\mathrm{i} = 0 \) zobrazených v Gaussovej rovine, je:}
{\( 2\sqrt[3]3 \)}{\( 2\sqrt3 \)}{\( \sqrt3 \)}{\( \sqrt[3]9\)}{\( 2\sqrt[3]9\)}{}}
\LANGes{
\Question{
Todas las soluciones de la ecuación \( x^6+3\sqrt5-6\mathrm{i} = 0 \) se pueden representar en el plano cartesiano. ¿Cuál es la distancia entre los dos puntos más lejanos?}
{\( 2\sqrt[3]3 \)}{\( 2\sqrt3 \)}{\( \sqrt3 \)}{\( \sqrt[3]9\)}{\( 2\sqrt[3]9\)}{}}
\End

\Data\ID{2010013406}
\Updated{2023-02-21 11:02:17}
\Level{C}
\SubArea{Powers and Roots of Complex Numbers}
\LANGen{
\Question{
Find the solution set of the following equation in the set of complex numbers.
\[ 
 x^{3} + 8\mathrm{i} = 0
\]}
{\(\left\{2\mathrm{i};\ \sqrt{3} -\mathrm{i};\ -\sqrt{3}-\mathrm{i}\right\}\)}{\(\left\{ -2\mathrm{i};\ \sqrt{3} -\mathrm{i};\ -\sqrt{3}-\mathrm{i}\right\}\)}{\(\left\{ -2;\ -\sqrt{3} +\mathrm{i};\ \sqrt{3}+\mathrm{i}\right\}\)}{\(\left\{ 2;\ -\sqrt{3} +\mathrm{i};\ \sqrt{3}+\mathrm{i}\right\}\)}{}{}}
\LANGpl{
\Question{
Znajdź zbiór rozwiązań następującego równania w zbiorze liczb zespolonych. 
\[ 
 x^{3} + 8\mathrm{i} = 0
\]}
{\(\left\{2\mathrm{i};\ \sqrt{3} -\mathrm{i};\ -\sqrt{3}-\mathrm{i}\right\}\)}{\(\left\{ -2\mathrm{i};\ \sqrt{3} -\mathrm{i};\ -\sqrt{3}-\mathrm{i}\right\}\)}{\(\left\{ -2;\ -\sqrt{3} +\mathrm{i};\ \sqrt{3}+\mathrm{i}\right\}\)}{\(\left\{ 2;\ -\sqrt{3} +\mathrm{i};\ \sqrt{3}+\mathrm{i}\right\}\)}{}{}}
\LANGcs{
\Question{
Určete množinu všech komplexních kořenů dané rovnice.
\[ 
 x^{3} + 8\mathrm{i} = 0
\]}
{\(\left\{2\mathrm{i};\ \sqrt{3} -\mathrm{i};\ -\sqrt{3}-\mathrm{i}\right\}\)}{\(\left\{ -2\mathrm{i};\ \sqrt{3} -\mathrm{i};\ -\sqrt{3}-\mathrm{i}\right\}\)}{\(\left\{ -2;\ -\sqrt{3} +\mathrm{i};\ \sqrt{3}+\mathrm{i}\right\}\)}{\(\left\{ 2;\ -\sqrt{3} +\mathrm{i};\ \sqrt{3}+\mathrm{i}\right\}\)}{}{}}
\LANGsk{
\Question{
Nájdite množinu všetkých riešení danej rovnice v množine komplexných čísel.
\[ 
 x^{3} + 8\mathrm{i} = 0
\]}
{\(\left\{2\mathrm{i};\ \sqrt{3} -\mathrm{i};\ -\sqrt{3}-\mathrm{i}\right\}\)}{\(\left\{ -2\mathrm{i};\ \sqrt{3} -\mathrm{i};\ -\sqrt{3}-\mathrm{i}\right\}\)}{\(\left\{ -2;\ -\sqrt{3} +\mathrm{i};\ \sqrt{3}+\mathrm{i}\right\}\)}{\(\left\{ 2;\ -\sqrt{3} +\mathrm{i};\ \sqrt{3}+\mathrm{i}\right\}\)}{}{}}
\LANGes{
\Question{
Halla todas las soluciones de la siguiente ecuación en el conjunto de los números complejos.
\[ 
 x^{3} + 8\mathrm{i} = 0
\]}
{\(\left\{2\mathrm{i};\ \sqrt{3} -\mathrm{i};\ -\sqrt{3}-\mathrm{i}\right\}\)}{\(\left\{ -2\mathrm{i};\ \sqrt{3} -\mathrm{i};\ -\sqrt{3}-\mathrm{i}\right\}\)}{\(\left\{ -2;\ -\sqrt{3} +\mathrm{i};\ \sqrt{3}+\mathrm{i}\right\}\)}{\(\left\{ 2;\ -\sqrt{3} +\mathrm{i};\ \sqrt{3}+\mathrm{i}\right\}\)}{}{}}
\End

\Data\ID{2010013407}
\Updated{2023-02-21 11:02:56}
\Level{C}
\SubArea{Powers and Roots of Complex Numbers}
\LANGen{
\Question{
Two solutions of the equation 
\[ 
 x^{3} + 1 - \mathrm{i} = 0
\]
are 
\[ 
\begin{aligned}x_{1}& = \root{6}\of{2}\left (\cos \frac{\pi}
{4} + \mathrm{i}\sin \frac{\pi}
{4} \right ),&
\\x_{2}& = \root{6}\of{2}\left (\cos \frac{11}
{12}\pi + \mathrm{i}\sin \frac{11}
{12}\pi \right ).
\\ \end{aligned}
\]
Find the third solution.}
{\(x_{3} = \root{6}\of{2}\left (\cos \frac{19}
{12}\pi + \mathrm{i}\sin \frac{19}
{12}\pi \right )\)}{\(x_{3} = \root{6}\of{2}\left (\cos \frac{7}
{12}\pi + \mathrm{i}\sin \frac{7}
{12}\pi \right )\)}{\(x_{3} = \root{6}\of{2}\left (\cos \frac{5}
{12}\pi + \mathrm{i}\sin \frac{5}
{12}\pi \right )\)}{\(x_{3} = \root{6}\of{2}\left (\cos \frac{13}
{12}\pi + \mathrm{i}\sin \frac{13}
{12}\pi \right )\)}{}{}}
\LANGpl{
\Question{
Dwa rozwiązania równania  
\[ 
 x^{3} + 1 - \mathrm{i} = 0
\]
są równe:
\[ 
\begin{aligned}x_{1}& = \root{6}\of{2}\left (\cos \frac{\pi}
{4} + \mathrm{i}\sin \frac{\pi}
{4} \right ),&
\\x_{2}& = \root{6}\of{2}\left (\cos \frac{11}
{12}\pi + \mathrm{i}\sin \frac{11}
{12}\pi \right ).
\\ \end{aligned}
\]
Znajdź trzecie rozwiązanie.}
{\(x_{3} = \root{6}\of{2}\left (\cos \frac{19}
{12}\pi + \mathrm{i}\sin \frac{19}
{12}\pi \right )\)}{\(x_{3} = \root{6}\of{2}\left (\cos \frac{7}
{12}\pi + \mathrm{i}\sin \frac{7}
{12}\pi \right )\)}{\(x_{3} = \root{6}\of{2}\left (\cos \frac{5}
{12}\pi + \mathrm{i}\sin \frac{5}
{12}\pi \right )\)}{\(x_{3} = \root{6}\of{2}\left (\cos \frac{13}
{12}\pi + \mathrm{i}\sin \frac{13}
{12}\pi \right )\)}{}{}}
\LANGcs{
\Question{
Dva kořeny rovnice
\[ 
 x^{3} + 1 - \mathrm{i} = 0
\]
jsou
\[ 
\begin{aligned}x_{1}& = \root{6}\of{2}\left (\cos \frac{\pi}
{4} + \mathrm{i}\sin \frac{\pi}
{4} \right ),&
\\x_{2}& = \root{6}\of{2}\left (\cos \frac{11}
{12}\pi + \mathrm{i}\sin \frac{11}
{12}\pi \right ).
\\ \end{aligned}
\]
Určete třetí kořen.}
{\(x_{3} = \root{6}\of{2}\left (\cos \frac{19}
{12}\pi + \mathrm{i}\sin \frac{19}
{12}\pi \right )\)}{\(x_{3} = \root{6}\of{2}\left (\cos \frac{7}
{12}\pi + \mathrm{i}\sin \frac{7}
{12}\pi \right )\)}{\(x_{3} = \root{6}\of{2}\left (\cos \frac{5}
{12}\pi + \mathrm{i}\sin \frac{5}
{12}\pi \right )\)}{\(x_{3} = \root{6}\of{2}\left (\cos \frac{13}
{12}\pi + \mathrm{i}\sin \frac{13}
{12}\pi \right )\)}{}{}}
\LANGsk{
\Question{
Dve riešenia rovnice
\[ 
 x^{3} + 1 - \mathrm{i} = 0
\]
sú 
\[ 
\begin{aligned}x_{1}& = \root{6}\of{2}\left (\cos \frac{\pi}
{4} + \mathrm{i}\sin \frac{\pi}
{4} \right ),&
\\x_{2}& = \root{6}\of{2}\left (\cos \frac{11}
{12}\pi + \mathrm{i}\sin \frac{11}
{12}\pi \right ).
\\ \end{aligned}
\]
Nájdite tretie riešenie.}
{\(x_{3} = \root{6}\of{2}\left (\cos \frac{19}
{12}\pi + \mathrm{i}\sin \frac{19}
{12}\pi \right )\)}{\(x_{3} = \root{6}\of{2}\left (\cos \frac{7}
{12}\pi + \mathrm{i}\sin \frac{7}
{12}\pi \right )\)}{\(x_{3} = \root{6}\of{2}\left (\cos \frac{5}
{12}\pi + \mathrm{i}\sin \frac{5}
{12}\pi \right )\)}{\(x_{3} = \root{6}\of{2}\left (\cos \frac{13}
{12}\pi + \mathrm{i}\sin \frac{13}
{12}\pi \right )\)}{}{}}
\LANGes{
\Question{
Dos soluciones de la ecuación 
\[ 
 x^{3} + 1 - \mathrm{i} = 0
\]
son 
\[ 
\begin{aligned}x_{1}& = \root{6}\of{2}\left (\cos \frac{\pi}
{4} + \mathrm{i}\sin \frac{\pi}
{4} \right ),&
\\x_{2}& = \root{6}\of{2}\left (\cos \frac{11}
{12}\pi + \mathrm{i}\sin \frac{11}
{12}\pi \right ).
\\ \end{aligned}
\]
Calcula la tercera solución.}
{\(x_{3} = \root{6}\of{2}\left (\cos \frac{19}
{12}\pi + \mathrm{i}\sin \frac{19}
{12}\pi \right )\)}{\(x_{3} = \root{6}\of{2}\left (\cos \frac{7}
{12}\pi + \mathrm{i}\sin \frac{7}
{12}\pi \right )\)}{\(x_{3} = \root{6}\of{2}\left (\cos \frac{5}
{12}\pi + \mathrm{i}\sin \frac{5}
{12}\pi \right )\)}{\(x_{3} = \root{6}\of{2}\left (\cos \frac{13}
{12}\pi + \mathrm{i}\sin \frac{13}
{12}\pi \right )\)}{}{}}
\End

\Data\ID{2010013408}
\Updated{2023-02-21 11:02:49}
\Level{C}
\SubArea{Powers and Roots of Complex Numbers}
\LANGen{
\Question{
Three solutions of the equation 
\[ 
 x^{4} - 2\mathrm{i} = 0
\]
are 
 \[\begin{aligned}x_{1} = \root{4}\of{2}\left (\cos \frac{1}{8}\pi + \mathrm{i}\sin \frac{1}{8}\pi   \right ),\\ x_{2} = \root{4}\of{2}\left (\cos \frac{5}{8}\pi + \mathrm{i}\sin \frac{5}{8}\pi   \right ),\\ x_{3} = \root{4}\of{2}\left (\cos \frac{9}{8}\pi + \mathrm{i}\sin \frac{9}{8}\pi   \right ).\\ \end{aligned}\]
Find the fourth solution.}
{\(x_{4} = \root{4}\of{2}\left (\cos \frac{13}{8}\pi + \mathrm{i}\sin \frac{13}{8}\pi   \right )\)}{\(x_{4} = \root{4}\of{2}\left (\cos \frac{11}{8}\pi + \mathrm{i}\sin \frac{11}{8}\pi   \right )\)}{\(x_{4} = \root{4}\of{2}\left (\cos \frac{15}{8}\pi + \mathrm{i}\sin \frac{15}{8}\pi   \right )\)}{\(x_{4} = \root{4}\of{2}\left (\cos \frac{3}{8}\pi + \mathrm{i}\sin \frac{3}{8}\pi   \right )\)}{}{}}
\LANGpl{
\Question{
Trzy rozwiązania równania 
\[ 
 x^{4} - 2\mathrm{i} = 0
\]
są równe:
 \[\begin{aligned}x_{1} = \root{4}\of{2}\left (\cos \frac{1}{8}\pi + \mathrm{i}\sin \frac{1}{8}\pi   \right ),\\ x_{2} = \root{4}\of{2}\left (\cos \frac{5}{8}\pi + \mathrm{i}\sin \frac{5}{8}\pi   \right ),\\ x_{3} = \root{4}\of{2}\left (\cos \frac{9}{8}\pi + \mathrm{i}\sin \frac{9}{8}\pi   \right ).\\ \end{aligned}\]
Znajdź czwarte rozwiązanie.}
{\(x_{4} = \root{4}\of{2}\left (\cos \frac{13}{8}\pi + \mathrm{i}\sin \frac{13}{8}\pi   \right )\)}{\(x_{4} = \root{4}\of{2}\left (\cos \frac{11}{8}\pi + \mathrm{i}\sin \frac{11}{8}\pi   \right )\)}{\(x_{4} = \root{4}\of{2}\left (\cos \frac{15}{8}\pi + \mathrm{i}\sin \frac{15}{8}\pi   \right )\)}{\(x_{4} = \root{4}\of{2}\left (\cos \frac{3}{8}\pi + \mathrm{i}\sin \frac{3}{8}\pi   \right )\)}{}{}}
\LANGcs{
\Question{
Tři kořeny rovnice
\[ 
 x^{4} - 2\mathrm{i} = 0
\]
jsou
  \[\begin{aligned}x_{1} = \root{4}\of{2}\left (\cos \frac{1}{8}\pi + \mathrm{i}\sin \frac{1}{8}\pi   \right ),\\ x_{2} = \root{4}\of{2}\left (\cos \frac{5}{8}\pi + \mathrm{i}\sin \frac{5}{8}\pi   \right ),\\ x_{3} = \root{4}\of{2}\left (\cos \frac{9}{8}\pi + \mathrm{i}\sin \frac{9}{8}\pi   \right ).\\ \end{aligned}\]
Určete čtvrtý kořen.}
{\(x_{4} = \root{4}\of{2}\left (\cos \frac{13}{8}\pi + \mathrm{i}\sin \frac{13}{8}\pi   \right )\)}{\(x_{4} = \root{4}\of{2}\left (\cos \frac{11}{8}\pi + \mathrm{i}\sin \frac{11}{8}\pi   \right )\)}{\(x_{4} = \root{4}\of{2}\left (\cos \frac{15}{8}\pi + \mathrm{i}\sin \frac{15}{8}\pi   \right )\)}{\(x_{4} = \root{4}\of{2}\left (\cos \frac{3}{8}\pi + \mathrm{i}\sin \frac{3}{8}\pi   \right )\)}{}{}}
\LANGsk{
\Question{
Tri riešenia rovnice
\[ 
 x^{4} - 2\mathrm{i} = 0
\]
sú 
 \[\begin{aligned}x_{1} = \root{4}\of{2}\left (\cos \frac{1}{8}\pi + \mathrm{i}\sin \frac{1}{8}\pi   \right ),\\ x_{2} = \root{4}\of{2}\left (\cos \frac{5}{8}\pi + \mathrm{i}\sin \frac{5}{8}\pi   \right ),\\ x_{3} = \root{4}\of{2}\left (\cos \frac{9}{8}\pi + \mathrm{i}\sin \frac{9}{8}\pi   \right ).\\ \end{aligned}\]
Nájdite štvrté riešenie.}
{\(x_{4} = \root{4}\of{2}\left (\cos \frac{13}{8}\pi + \mathrm{i}\sin \frac{13}{8}\pi   \right )\)}{\(x_{4} = \root{4}\of{2}\left (\cos \frac{11}{8}\pi + \mathrm{i}\sin \frac{11}{8}\pi   \right )\)}{\(x_{4} = \root{4}\of{2}\left (\cos \frac{15}{8}\pi + \mathrm{i}\sin \frac{15}{8}\pi   \right )\)}{\(x_{4} = \root{4}\of{2}\left (\cos \frac{3}{8}\pi + \mathrm{i}\sin \frac{3}{8}\pi   \right )\)}{}{}}
\LANGes{
\Question{
Tres soluciones de la ecuación 
\[ 
 x^{4} - 2\mathrm{i} = 0
\]
son 
 \[\begin{aligned}x_{1} = \root{4}\of{2}\left (\cos \frac{1}{8}\pi + \mathrm{i}\sin \frac{1}{8}\pi   \right ),\\ x_{2} = \root{4}\of{2}\left (\cos \frac{5}{8}\pi + \mathrm{i}\sin \frac{5}{8}\pi   \right ),\\ x_{3} = \root{4}\of{2}\left (\cos \frac{9}{8}\pi + \mathrm{i}\sin \frac{9}{8}\pi   \right ).\\ \end{aligned}\]
Calcula la cuarta solución.}
{\(x_{4} = \root{4}\of{2}\left (\cos \frac{13}{8}\pi + \mathrm{i}\sin \frac{13}{8}\pi   \right )\)}{\(x_{4} = \root{4}\of{2}\left (\cos \frac{11}{8}\pi + \mathrm{i}\sin \frac{11}{8}\pi   \right )\)}{\(x_{4} = \root{4}\of{2}\left (\cos \frac{15}{8}\pi + \mathrm{i}\sin \frac{15}{8}\pi   \right )\)}{\(x_{4} = \root{4}\of{2}\left (\cos \frac{3}{8}\pi + \mathrm{i}\sin \frac{3}{8}\pi   \right )\)}{}{}}
\End

\Data\ID{2010013409}
\Updated{2023-02-21 11:02:00}
\Level{C}
\SubArea{Powers and Roots of Complex Numbers}
\LANGen{
\Question{
Three solutions of the equation 
\[ 
 x^{4} + 8\mathrm{i} = 0
\]
are 
 \[\begin{aligned}x_{1} = \root{4}\of{8}\left (\cos \frac{3}{8}\pi + \mathrm{i}\sin \frac{3}{8}\pi   \right ), \\ x_{2} = \root{4}\of{8}\left (\cos \frac{7}{8}\pi + \mathrm{i}\sin \frac{7}{8}\pi   \right ),\\ x_{3} = \root{4}\of{8}\left (\cos \frac{15}{8}\pi + \mathrm{i}\sin \frac{15}{8}\pi   \right ).\\ \end{aligned}\]
Find the fourth solution.}
{\(x_{4} = \root{4}\of{8}\left (\cos \frac{11}{8}\pi + \mathrm{i}\sin \frac{11}{8}\pi   \right )\)}{\(x_{4} = \root{4}\of{8}\left (\cos \frac{9}{8}\pi + \mathrm{i}\sin \frac{9}{8}\pi   \right )\)}{\(x_{4} = \root{4}\of{8}\left (\cos \frac{5}{8}\pi + \mathrm{i}\sin \frac{5}{8}\pi   \right )\)}{\(x_{4} = \root{4}\of{8}\left (\cos \frac{1}{8}\pi + \mathrm{i}\sin \frac{1}{8}\pi   \right )\)}{}{}}
\LANGpl{
\Question{
Trzy rozwiązania równania 
\[ 
 x^{4} + 8\mathrm{i} = 0
\]
są równe:
 \[\begin{aligned}x_{1} = \root{4}\of{8}\left (\cos \frac{3}{8}\pi + \mathrm{i}\sin \frac{3}{8}\pi   \right ), \\ x_{2} = \root{4}\of{8}\left (\cos \frac{7}{8}\pi + \mathrm{i}\sin \frac{7}{8}\pi   \right ),\\ x_{3} = \root{4}\of{8}\left (\cos \frac{15}{8}\pi + \mathrm{i}\sin \frac{15}{8}\pi   \right ).\\ \end{aligned}\]
Znajdź czwarte rozwiązanie.}
{\(x_{4} = \root{4}\of{8}\left (\cos \frac{11}{8}\pi + \mathrm{i}\sin \frac{11}{8}\pi   \right )\)}{\(x_{4} = \root{4}\of{8}\left (\cos \frac{9}{8}\pi + \mathrm{i}\sin \frac{9}{8}\pi   \right )\)}{\(x_{4} = \root{4}\of{8}\left (\cos \frac{5}{8}\pi + \mathrm{i}\sin \frac{5}{8}\pi   \right )\)}{\(x_{4} = \root{4}\of{8}\left (\cos \frac{1}{8}\pi + \mathrm{i}\sin \frac{1}{8}\pi   \right )\)}{}{}}
\LANGcs{
\Question{
Tři kořeny rovnice
\[ 
 x^{4} + 8\mathrm{i} = 0
\]
jsou
 \[\begin{aligned}x_{1} = \root{4}\of{8}\left (\cos \frac{3}{8}\pi + \mathrm{i}\sin \frac{3}{8}\pi   \right ), \\ x_{2} = \root{4}\of{8}\left (\cos \frac{7}{8}\pi + \mathrm{i}\sin \frac{7}{8}\pi   \right ),\\ x_{3} = \root{4}\of{8}\left (\cos \frac{15}{8}\pi + \mathrm{i}\sin \frac{15}{8}\pi   \right ).\\ \end{aligned}\]
Určete čtvrtý kořen.}
{\(x_{4} = \root{4}\of{8}\left (\cos \frac{11}{8}\pi + \mathrm{i}\sin \frac{11}{8}\pi   \right )\)}{\(x_{4} = \root{4}\of{8}\left (\cos \frac{9}{8}\pi + \mathrm{i}\sin \frac{9}{8}\pi   \right )\)}{\(x_{4} = \root{4}\of{8}\left (\cos \frac{5}{8}\pi + \mathrm{i}\sin \frac{5}{8}\pi   \right )\)}{\(x_{4} = \root{4}\of{8}\left (\cos \frac{1}{8}\pi + \mathrm{i}\sin \frac{1}{8}\pi   \right )\)}{}{}}
\LANGsk{
\Question{
Tri riešenia rovnice
\[ 
 x^{4} + 8\mathrm{i} = 0
\]
sú 
 \[\begin{aligned}x_{1} = \root{4}\of{8}\left (\cos \frac{3}{8}\pi + \mathrm{i}\sin \frac{3}{8}\pi   \right ), \\ x_{2} = \root{4}\of{8}\left (\cos \frac{7}{8}\pi + \mathrm{i}\sin \frac{7}{8}\pi   \right ),\\ x_{3} = \root{4}\of{8}\left (\cos \frac{15}{8}\pi + \mathrm{i}\sin \frac{15}{8}\pi   \right ).\\ \end{aligned}\]
Nájdite štvrté riešenie.}
{\(x_{4} = \root{4}\of{8}\left (\cos \frac{11}{8}\pi + \mathrm{i}\sin \frac{11}{8}\pi   \right )\)}{\(x_{4} = \root{4}\of{8}\left (\cos \frac{9}{8}\pi + \mathrm{i}\sin \frac{9}{8}\pi   \right )\)}{\(x_{4} = \root{4}\of{8}\left (\cos \frac{5}{8}\pi + \mathrm{i}\sin \frac{5}{8}\pi   \right )\)}{\(x_{4} = \root{4}\of{8}\left (\cos \frac{1}{8}\pi + \mathrm{i}\sin \frac{1}{8}\pi   \right )\)}{}{}}
\LANGes{
\Question{
Tres soluciones de la ecuación 
\[ 
 x^{4} + 8\mathrm{i} = 0
\]
son 
 \[\begin{aligned}x_{1} = \root{4}\of{8}\left (\cos \frac{3}{8}\pi + \mathrm{i}\sin \frac{3}{8}\pi   \right ), \\ x_{2} = \root{4}\of{8}\left (\cos \frac{7}{8}\pi + \mathrm{i}\sin \frac{7}{8}\pi   \right ),\\ x_{3} = \root{4}\of{8}\left (\cos \frac{15}{8}\pi + \mathrm{i}\sin \frac{15}{8}\pi   \right ).\\ \end{aligned}\]
Calcula la cuarta solución.}
{\(x_{4} = \root{4}\of{8}\left (\cos \frac{11}{8}\pi + \mathrm{i}\sin \frac{11}{8}\pi   \right )\)}{\(x_{4} = \root{4}\of{8}\left (\cos \frac{9}{8}\pi + \mathrm{i}\sin \frac{9}{8}\pi   \right )\)}{\(x_{4} = \root{4}\of{8}\left (\cos \frac{5}{8}\pi + \mathrm{i}\sin \frac{5}{8}\pi   \right )\)}{\(x_{4} = \root{4}\of{8}\left (\cos \frac{1}{8}\pi + \mathrm{i}\sin \frac{1}{8}\pi   \right )\)}{}{}}
\End

\Data\ID{2010013410}
\Updated{2023-02-21 11:02:18}
\Level{C}
\SubArea{Powers and Roots of Complex Numbers}
\LANGen{
\Question{
Consider the equation
\[ 
 x^{4} = ( 2 - \mathrm{i} )^{4}.
\]
Find the product of all its roots in the set of complex numbers.}
{\(7+24\mathrm{i}\)}{\(-7-24\mathrm{i}\)}{\(16\)}{\(-16\)}{}{}}
\LANGpl{
\Question{
Rozwiąż równanie w zbiorze liczb zespolonych
\[ 
 x^{4} = ( 2 - \mathrm{i} )^{4}.
\]
Znajdź iloczyn wszystkich jego pierwiastków.}
{\(7+24\mathrm{i}\)}{\(-7-24\mathrm{i}\)}{\(16\)}{\(-16\)}{}{}}
\LANGcs{
\Question{
Určete součin všech komplexních kořenů rovnice 
\[ 
 x^{4} = ( 2 - \mathrm{i} )^{4}.
\]}
{\(7+24\mathrm{i}\)}{\(-7-24\mathrm{i}\)}{\(16\)}{\(-16\)}{}{}}
\LANGsk{
\Question{
Určte súčin všetkých komplexných koreňov rovnice
\[ 
 x^{4} = ( 2 - \mathrm{i} )^{4}.
\]}
{\(7+24\mathrm{i}\)}{\(-7-24\mathrm{i}\)}{\(16\)}{\(-16\)}{}{}}
\LANGes{
\Question{
Dada al ecuación
\[ 
 x^{4} = ( 2 - \mathrm{i} )^{4}.
\]
Halla el producto de todas sus raíces en el conjunto de los números complejos.}
{\(7+24\mathrm{i}\)}{\(-7-24\mathrm{i}\)}{\(16\)}{\(-16\)}{}{}}
\End

\Data\ID{2010013411}
\Updated{2023-02-21 11:02:31}
\Level{C}
\SubArea{Powers and Roots of Complex Numbers}
\LANGen{
\Question{
Consider the equation
\[ 
 x^{4} = ( 2 + \mathrm{i} )^{4}.
\]
Find the product of all its roots in the set of complex numbers.}
{\(7-24\mathrm{i}\)}{\(-7+24\mathrm{i}\)}{\(16\)}{\(-16\)}{}{}}
\LANGpl{
\Question{
Rozwiąż równanie w zbiorze liczb zespolonych
\[ 
 x^{4} = ( 2 + \mathrm{i} )^{4}.
\]
Znajdź iloczyn wszystkich jego pierwiastków.}
{\(7-24\mathrm{i}\)}{\(-7+24\mathrm{i}\)}{\(16\)}{\(-16\)}{}{}}
\LANGcs{
\Question{
Určete součin všech komplexních kořenů rovnice
\[ 
 x^{4} = ( 2 + \mathrm{i} )^{4}.
\]}
{\(7-24\mathrm{i}\)}{\(-7+24\mathrm{i}\)}{\(16\)}{\(-16\)}{}{}}
\LANGsk{
\Question{
Určte súčin všetkých komplexných koreňov rovnice
\[ 
 x^{4} = ( 2 + \mathrm{i} )^{4}.
\]}
{\(7-24\mathrm{i}\)}{\(-7+24\mathrm{i}\)}{\(16\)}{\(-16\)}{}{}}
\LANGes{
\Question{
Dada la ecuación
\[ 
 x^{4} = ( 2 + \mathrm{i} )^{4}.
\]
Halla el producto de todas sus raíces en el conjunto de los números complejos.}
{\(7-24\mathrm{i}\)}{\(-7+24\mathrm{i}\)}{\(16\)}{\(-16\)}{}{}}
\End

\Data\ID{2010013412}
\Updated{2023-02-21 11:02:43}
\Level{C}
\SubArea{Powers and Roots of Complex Numbers}
\LANGen{
\Question{
Which of the given numbers does not belong to the solution set of the following equation?
\[x^{4}-1+\mathrm{i}=0\]}
{\(\root{8}\of{2}\left (\cos \frac{3\pi}{16} + \mathrm{i}\sin \frac{3\pi}{16}\right )\)}{\(-\root{4}\of{1-\mathrm{i}}\)}{\(-\mathrm{i}\root{4}\of{1-\mathrm{i}}\)}{\(\root{8}\of{2}\left (\cos \left(-\frac{\pi}{16}\right) + \mathrm{i}\sin \left(-\frac{\pi}{16}\right)\right )\)}{}{}}
\LANGpl{
\Question{
Która z liczb nie należy do zbioru rozwiązań poniższego równania?
\[x^{4}-1+\mathrm{i}=0\]}
{\(\root{8}\of{2}\left (\cos \frac{3\pi}{16} + \mathrm{i}\sin \frac{3\pi}{16}\right )\)}{\(-\root{4}\of{1-\mathrm{i}}\)}{\(-\mathrm{i}\root{4}\of{1-\mathrm{i}}\)}{\(\root{8}\of{2}\left (\cos \left(-\frac{\pi}{16}\right) + \mathrm{i}\sin \left(-\frac{\pi}{16}\right)\right )\)}{}{}}
\LANGcs{
\Question{
Které z čísel nepatří do množiny řešení dané rovnice?
\[x^{4}-1+\mathrm{i}=0\]}
{\(\root{8}\of{2}\left (\cos \frac{3\pi}{16} + \mathrm{i}\sin \frac{3\pi}{16}\right )\)}{\(-\root{4}\of{1-\mathrm{i}}\)}{\(-\mathrm{i}\root{4}\of{1-\mathrm{i}}\)}{\(\root{8}\of{2}\left (\cos \left(-\frac{\pi}{16}\right) + \mathrm{i}\sin \left(-\frac{\pi}{16}\right)\right )\)}{}{}}
\LANGsk{
\Question{
Ktoré z daných čísel nepatrí do množiny riešení danej rovnice?
\[x^{4}-1+\mathrm{i}=0\]}
{\(\root{8}\of{2}\left (\cos \frac{3\pi}{16} + \mathrm{i}\sin \frac{3\pi}{16}\right )\)}{\(-\root{4}\of{1-\mathrm{i}}\)}{\(-\mathrm{i}\root{4}\of{1-\mathrm{i}}\)}{\(\root{8}\of{2}\left (\cos \left(-\frac{\pi}{16}\right) + \mathrm{i}\sin \left(-\frac{\pi}{16}\right)\right )\)}{}{}}
\LANGes{
\Question{
¿Cuál de estos números no pertenece al conjunto de soluciones de la siguiente ecuación?
\[x^{4}-1+\mathrm{i}=0\]}
{\(\root{8}\of{2}\left (\cos \frac{3\pi}{16} + \mathrm{i}\sin \frac{3\pi}{16}\right )\)}{\(-\root{4}\of{1-\mathrm{i}}\)}{\(-\mathrm{i}\root{4}\of{1-\mathrm{i}}\)}{\(\root{8}\of{2}\left (\cos \left(-\frac{\pi}{16}\right) + \mathrm{i}\sin \left(-\frac{\pi}{16}\right)\right )\)}{}{}}
\End

\Data\ID{2010013413}
\Updated{2023-02-21 11:02:58}
\Level{C}
\SubArea{Powers and Roots of Complex Numbers}
\LANGen{
\Question{
Which of the given numbers does not belong to the solution set of the following equation?
\[x^{4}+1+\mathrm{i}=0\]}
{\(\root{8}\of{2}\left (\cos \frac{3\pi}{16} + \mathrm{i}\sin \frac{3\pi}{16}\right )\)}{\(\mathrm{i}\root{4}\of{-1-\mathrm{i}}\)}{\(\root{4}\of{-1-\mathrm{i}}\)}{\(\root{8}\of{2}\left (\cos \frac{5\pi}{16} + \mathrm{i}\sin \frac{5\pi}{16}\right )\)}{}{}}
\LANGpl{
\Question{
Która z podanych liczb nie należy do zbioru rozwiązań poniższego równania?
\[x^{4}+1+\mathrm{i}=0\]}
{\(\root{8}\of{2}\left (\cos \frac{3\pi}{16} + \mathrm{i}\sin \frac{3\pi}{16}\right )\)}{\(\mathrm{i}\root{4}\of{-1-\mathrm{i}}\)}{\(\root{4}\of{-1-\mathrm{i}}\)}{\(\root{8}\of{2}\left (\cos \frac{5\pi}{16} + \mathrm{i}\sin \frac{5\pi}{16}\right )\)}{}{}}
\LANGcs{
\Question{
Které z čísel nepatří do množiny řešení dané rovnice?
\[x^{4}+1+\mathrm{i}=0\]}
{\(\root{8}\of{2}\left (\cos \frac{3\pi}{16} + \mathrm{i}\sin \frac{3\pi}{16}\right )\)}{\(\mathrm{i}\root{4}\of{-1-\mathrm{i}}\)}{\(\root{4}\of{-1-\mathrm{i}}\)}{\(\root{8}\of{2}\left (\cos \frac{5\pi}{16} + \mathrm{i}\sin \frac{5\pi}{16}\right )\)}{}{}}
\LANGsk{
\Question{
Ktoré z daných čísel nepatrí do množiny riešení danej rovnice?
\[x^{4}+1+\mathrm{i}=0\]}
{\(\root{8}\of{2}\left (\cos \frac{3\pi}{16} + \mathrm{i}\sin \frac{3\pi}{16}\right )\)}{\(\mathrm{i}\root{4}\of{-1-\mathrm{i}}\)}{\(\root{4}\of{-1-\mathrm{i}}\)}{\(\root{8}\of{2}\left (\cos \frac{5\pi}{16} + \mathrm{i}\sin \frac{5\pi}{16}\right )\)}{}{}}
\LANGes{
\Question{
¿Cuál de estos números no pertenece al conjunto de soluciones de la siguiente ecuación?
\[x^{4}+1+\mathrm{i}=0\]}
{\(\root{8}\of{2}\left (\cos \frac{3\pi}{16} + \mathrm{i}\sin \frac{3\pi}{16}\right )\)}{\(\mathrm{i}\root{4}\of{-1-\mathrm{i}}\)}{\(\root{4}\of{-1-\mathrm{i}}\)}{\(\root{8}\of{2}\left (\cos \frac{5\pi}{16} + \mathrm{i}\sin \frac{5\pi}{16}\right )\)}{}{}}
\End

\Data\ID{9000034303}
\Updated{2023-02-21 11:02:58}
\Level{C}
\SubArea{Powers and Roots of Complex Numbers}
\LANGen{
\Question{
Find the solution set of the following equation in the set of complex numbers.
\[ 
 x^{3} + \mathrm{i} = 0
\]}
{\(\{\mathrm{i};\ \frac{\sqrt{3}} 
 {2} -\frac{1} 
{2}\mathrm{i};\ -\frac{\sqrt{3}}
 {2} -\frac{1} 
{2}\mathrm{i}\}\)}{\(\{ - 1;\ -\frac{\sqrt{3}}
 {2} + \frac{1} 
{2}\mathrm{i};\ -\frac{\sqrt{3}}
 {2} -\frac{1} 
{2}\mathrm{i}\}\)}{\(\{ - 1;\ \frac{\sqrt{3}} 
 {2} -\frac{1} 
{2}\mathrm{i};\ -\frac{\sqrt{3}}
 {2} -\frac{1} 
{2}\mathrm{i}\}\)}{\(\{\mathrm{i};\ -\frac{\sqrt{3}}
 {2} + \frac{1} 
{2}\mathrm{i};\ -\frac{\sqrt{3}}
 {2} -\frac{1} 
{2}\mathrm{i}\}\)}{}{}}
\LANGpl{
\Question{
Wyznacz zbiór rozwiązań równania w zbiorze liczb zespolonych.
\[ 
 x^{3} + \mathrm{i} = 0
\]}
{\(\{\mathrm{i};\ \frac{\sqrt{3}} 
 {2} -\frac{1} 
{2}\mathrm{i};\ -\frac{\sqrt{3}}
 {2} -\frac{1} 
{2}\mathrm{i}\}\)}{\(\{ - 1;\ -\frac{\sqrt{3}}
 {2} + \frac{1} 
{2}\mathrm{i};\ -\frac{\sqrt{3}}
 {2} -\frac{1} 
{2}\mathrm{i}\}\)}{\(\{ - 1;\ \frac{\sqrt{3}} 
 {2} -\frac{1} 
{2}\mathrm{i};\ -\frac{\sqrt{3}}
 {2} -\frac{1} 
{2}\mathrm{i}\}\)}{\(\{\mathrm{i};\ -\frac{\sqrt{3}}
 {2} + \frac{1} 
{2}\mathrm{i};\ -\frac{\sqrt{3}}
 {2} -\frac{1} 
{2}\mathrm{i}\}\)}{}{}}
\LANGcs{
\Question{
Množinou všech komplexních řešení rovnice
\(x^{3} + \mathrm{i} = 0\) je:}
{\(\{\mathrm{i};\ \frac{\sqrt{3}} 
 {2} -\frac{1} 
{2}\mathrm{i};\ -\frac{\sqrt{3}}
 {2} -\frac{1} 
{2}\mathrm{i}\}\)}{\(\{ - 1;\ -\frac{\sqrt{3}}
 {2} + \frac{1} 
{2}\mathrm{i};\ -\frac{\sqrt{3}}
 {2} -\frac{1} 
{2}\mathrm{i}\}\)}{\(\{ - 1;\ \frac{\sqrt{3}} 
 {2} -\frac{1} 
{2}\mathrm{i};\ -\frac{\sqrt{3}}
 {2} -\frac{1} 
{2}\mathrm{i}\}\)}{\(\{\mathrm{i};\ -\frac{\sqrt{3}}
 {2} + \frac{1} 
{2}\mathrm{i};\ -\frac{\sqrt{3}}
 {2} -\frac{1} 
{2}\mathrm{i}\}\)}{}{}}
\LANGsk{
\Question{
Nájdite množinu všetkých riešení danej rovnice v množine komplexných čísel. 
\[ 
 x^{3} + \mathrm{i} = 0
\]}
{\(\{\mathrm{i};\ \frac{\sqrt{3}} 
 {2} -\frac{1} 
{2}\mathrm{i};\ -\frac{\sqrt{3}}
 {2} -\frac{1} 
{2}\mathrm{i}\}\)}{\(\{ - 1;\ -\frac{\sqrt{3}}
 {2} + \frac{1} 
{2}\mathrm{i};\ -\frac{\sqrt{3}}
 {2} -\frac{1} 
{2}\mathrm{i}\}\)}{\(\{ - 1;\ \frac{\sqrt{3}} 
 {2} -\frac{1} 
{2}\mathrm{i};\ -\frac{\sqrt{3}}
 {2} -\frac{1} 
{2}\mathrm{i}\}\)}{\(\{\mathrm{i};\ -\frac{\sqrt{3}}
 {2} + \frac{1} 
{2}\mathrm{i};\ -\frac{\sqrt{3}}
 {2} -\frac{1} 
{2}\mathrm{i}\}\)}{}{}}
\LANGes{
\Question{
Determina el conjunto de soluciones de la siguiente ecuación en el conjunto de los números complejos.
\[ 
 x^{3} + \mathrm{i} = 0
\]}
{\(\{\mathrm{i};\ \frac{\sqrt{3}} 
 {2} -\frac{1} 
{2}\mathrm{i};\ -\frac{\sqrt{3}}
 {2} -\frac{1} 
{2}\mathrm{i}\}\)}{\(\{ - 1;\ -\frac{\sqrt{3}}
 {2} + \frac{1} 
{2}\mathrm{i};\ -\frac{\sqrt{3}}
 {2} -\frac{1} 
{2}\mathrm{i}\}\)}{\(\{ - 1;\ \frac{\sqrt{3}} 
 {2} -\frac{1} 
{2}\mathrm{i};\ -\frac{\sqrt{3}}
 {2} -\frac{1} 
{2}\mathrm{i}\}\)}{\(\{\mathrm{i};\ -\frac{\sqrt{3}}
 {2} + \frac{1} 
{2}\mathrm{i};\ -\frac{\sqrt{3}}
 {2} -\frac{1} 
{2}\mathrm{i}\}\)}{}{}}
\End

\Data\ID{9000034307}
\Updated{2023-02-21 11:02:55}
\Level{C}
\SubArea{Powers and Roots of Complex Numbers}
\LANGen{
\Question{
Let \(x\) 
be any of the complex solutions of the following equation. 
\[ 
 x^{5} + \sqrt{3} -\mathrm{i} = 0
\]
Find the absolute value of \(x\).}
{\(\root{5}\of{2}\)}{\(2\)}{\(\root{5}\of{4}\)}{\(\root{10}\of{2}\)}{}{}}
\LANGpl{
\Question{
Wartość bezwzględna każdego rozwiązania równania 
\[ 
 x^{5} + \sqrt{3} -\mathrm{i} = 0
\]
wynosi:}
{\(\root{5}\of{2}\)}{\(2\)}{\(\root{5}\of{4}\)}{\(\root{10}\of{2}\)}{}{}}
\LANGcs{
\Question{
Absolutní hodnota každého řešení rovnice
\(x^{5} + \sqrt{3} -\mathrm{i} = 0\) je:}
{\(\root{5}\of{2}\)}{\(2\)}{\(\root{5}\of{4}\)}{\(\root{10}\of{2}\)}{}{}}
\LANGsk{
\Question{
Nech \(x\) 
je jedno z komplexných riešení danej rovnice.  
\[ 
 x^{5} + \sqrt{3} -\mathrm{i} = 0
\]
Nájdite absolútnu hodnotu \(x\).}
{\(\root{5}\of{2}\)}{\(2\)}{\(\root{5}\of{4}\)}{\(\root{10}\of{2}\)}{}{}}
\LANGes{
\Question{
Sea \(x\) 
cualquiera de las soluciones complejas de la siguiente ecuación.
\[ 
 x^{5} + \sqrt{3} -\mathrm{i} = 0
\]
Calcula el valor absoluto de \(x\).}
{\(\root{5}\of{2}\)}{\(2\)}{\(\root{5}\of{4}\)}{\(\root{10}\of{2}\)}{}{}}
\End

\Data\ID{9000034308}
\Updated{2023-02-21 11:02:08}
\Level{C}
\SubArea{Powers and Roots of Complex Numbers}
\LANGen{
\Question{
Two solutions of the equation 
\[ 
 x^{3} + 1 + \mathrm{i} = 0
\]
are 
\[ 
\begin{aligned}x_{1}& = \root{6}\of{2}\left (\cos \frac{5}
{12}\pi + \mathrm{i}\sin \frac{5}
{12}\pi \right ),&
\\x_{2}& = \root{6}\of{2}\left (\cos \frac{13}
{12}\pi + \mathrm{i}\sin \frac{13}
{12}\pi \right ).
\\ \end{aligned}
\]
Find the third solution.}
{\(x_{3} = \root{6}\of{2}\left (\cos \frac{21}
{12}\pi + \mathrm{i}\sin \frac{21}
{12}\pi \right )\)}{\(x_{3} = \root{6}\of{2}\left (\cos \frac{9}
{12}\pi + \mathrm{i}\sin \frac{9}
{12}\pi \right )\)}{\(x_{3} = \root{6}\of{2}\left (\cos \frac{17}
{12}\pi + \mathrm{i}\sin \frac{17}
{12}\pi \right )\)}{\(x_{3} = \root{6}\of{2}\left (\cos \frac{19}
{12}\pi + \mathrm{i}\sin \frac{19}
{12}\pi \right )\)}{}{}}
\LANGpl{
\Question{
Równanie
\[ 
 x^{3} + 1 + \mathrm{i} = 0
\]
ma dwa rozwiązania
\[ 
\begin{aligned}x_{1}& = \root{6}\of{2}\left (\cos \frac{5}
{12}\pi + \mathrm{i}\sin \frac{5}
{12}\pi \right ),&
\\x_{2}& = \root{6}\of{2}\left (\cos \frac{13}
{12}\pi + \mathrm{i}\sin \frac{13}
{12}\pi \right ).
\\ \end{aligned}
\]
Wskaż trzecie.}
{\(x_{3} = \root{6}\of{2}\left (\cos \frac{21}
{12}\pi + \mathrm{i}\sin \frac{21}
{12}\pi \right )\)}{\(x_{3} = \root{6}\of{2}\left (\cos \frac{9}
{12}\pi + \mathrm{i}\sin \frac{9}
{12}\pi \right )\)}{\(x_{3} = \root{6}\of{2}\left (\cos \frac{17}
{12}\pi + \mathrm{i}\sin \frac{17}
{12}\pi \right )\)}{\(x_{3} = \root{6}\of{2}\left (\cos \frac{19}
{12}\pi + \mathrm{i}\sin \frac{19}
{12}\pi \right )\)}{}{}}
\LANGcs{
\Question{
Dvě z řešení rovnice
 \[x^{3} + 1 + \mathrm{i} = 0\]
jsou
\[ 
 x_{1} = \root{6}\of{2}\left (\cos \frac{5}
{12}\pi + \mathrm{i}\sin \frac{5}
{12}\pi \right ),
\]
\[ 
 x_{2} = \root{6}\of{2}\left (\cos \frac{13}
{12}\pi + \mathrm{i}\sin \frac{13}
{12}\pi \right ).
\]
 Třetím řešení rovnice je:}
{\(x_{3} = \root{6}\of{2}\left (\cos \frac{21}
{12}\pi + \mathrm{i}\sin \frac{21}
{12}\pi \right )\)}{\(x_{3} = \root{6}\of{2}\left (\cos \frac{9}
{12}\pi + \mathrm{i}\sin \frac{9}
{12}\pi \right )\)}{\(x_{3} = \root{6}\of{2}\left (\cos \frac{17}
{12}\pi + \mathrm{i}\sin \frac{17}
{12}\pi \right )\)}{\(x_{3} = \root{6}\of{2}\left (\cos \frac{19}
{12}\pi + \mathrm{i}\sin \frac{19}
{12}\pi \right )\)}{}{}}
\LANGsk{
\Question{
Dve riešenia rovnice
\[ 
 x^{3} + 1 + \mathrm{i} = 0
\]
sú 
\[ 
\begin{aligned}x_{1}& = \root{6}\of{2}\left (\cos \frac{5}
{12}\pi + \mathrm{i}\sin \frac{5}
{12}\pi \right ),&
\\x_{2}& = \root{6}\of{2}\left (\cos \frac{13}
{12}\pi + \mathrm{i}\sin \frac{13}
{12}\pi \right ).
\\ \end{aligned}
\]
Nájdite tretie riešenie.}
{\(x_{3} = \root{6}\of{2}\left (\cos \frac{21}
{12}\pi + \mathrm{i}\sin \frac{21}
{12}\pi \right )\)}{\(x_{3} = \root{6}\of{2}\left (\cos \frac{9}
{12}\pi + \mathrm{i}\sin \frac{9}
{12}\pi \right )\)}{\(x_{3} = \root{6}\of{2}\left (\cos \frac{17}
{12}\pi + \mathrm{i}\sin \frac{17}
{12}\pi \right )\)}{\(x_{3} = \root{6}\of{2}\left (\cos \frac{19}
{12}\pi + \mathrm{i}\sin \frac{19}
{12}\pi \right )\)}{}{}}
\LANGes{
\Question{
Dos soluciones de la ecuación
\[ 
 x^{3} + 1 + \mathrm{i} = 0
\]
son
\[ 
\begin{aligned}x_{1}& = \root{6}\of{2}\left (\cos \frac{5}
{12}\pi + \mathrm{i}\sin \frac{5}
{12}\pi \right ),&
\\x_{2}& = \root{6}\of{2}\left (\cos \frac{13}
{12}\pi + \mathrm{i}\sin \frac{13}
{12}\pi \right ).
\\ \end{aligned}
\]
Calcula la tercera solución.}
{\(x_{3} = \root{6}\of{2}\left (\cos \frac{21}
{12}\pi + \mathrm{i}\sin \frac{21}
{12}\pi \right )\)}{\(x_{3} = \root{6}\of{2}\left (\cos \frac{9}
{12}\pi + \mathrm{i}\sin \frac{9}
{12}\pi \right )\)}{\(x_{3} = \root{6}\of{2}\left (\cos \frac{17}
{12}\pi + \mathrm{i}\sin \frac{17}
{12}\pi \right )\)}{\(x_{3} = \root{6}\of{2}\left (\cos \frac{19}
{12}\pi + \mathrm{i}\sin \frac{19}
{12}\pi \right )\)}{}{}}
\End

\Data\ID{9000034309}
\Updated{2023-02-21 11:02:19}
\Level{C}
\SubArea{Powers and Roots of Complex Numbers}
\LANGen{
\Question{
Find the angle \(\varphi \) 
such that the angles in the polar form of any two solutions of the equation
\[ 
 x^{5} - 1 + \mathrm{i}\sqrt{3} = 0
\]
differ by an integer multiple of \(\varphi \).}
{\(\varphi = \frac{2} 
{5}\pi \)}{\(\varphi = \frac{3} 
{5}\pi \)}{\(\varphi = \frac{4} 
{5}\pi \)}{\(\varphi =\pi \)}{}{}}
\LANGpl{
\Question{
Wskaż kąt  \(\varphi \) 
tak, aby kąty  w postaci biegunowej jakichkolwiek dwóch rozwiązań równania
\[ 
 x^{5} - 1 + \mathrm{i}\sqrt{3} = 0
\]
różniły się liczbą całkowitą wielokrotności  \(\varphi \).}
{\(\varphi = \frac{2} 
{5}\pi \)}{\(\varphi = \frac{3} 
{5}\pi \)}{\(\varphi = \frac{4} 
{5}\pi \)}{\(\varphi =\pi \)}{}{}}
\LANGcs{
\Question{
Argumenty libovolných dvou řešení rovnice
\(x^{5} - 1 + \mathrm{i}\sqrt{3} = 0\) se
liší o celočíselný násobek čísla:}
{\(\varphi=\frac{2}
{5}\pi \)}{\(\varphi=\frac{3}
{5}\pi \)}{\(\varphi=\frac{4}
{5}\pi \)}{\(\varphi=\pi \)}{}{}}
\LANGsk{
\Question{
Argumenty ľubovoľných dvoch riešení rovnice 
\[ 
 x^{5} - 1 + \mathrm{i}\sqrt{3} = 0
\]
sa líšia o celočíselný násobok čísla:}
{\(\varphi = \frac{2} 
{5}\pi \)}{\(\varphi = \frac{3} 
{5}\pi \)}{\(\varphi = \frac{4} 
{5}\pi \)}{\(\varphi =\pi \)}{}{}}
\LANGes{
\Question{
Calcula el ángulo \(\varphi \) 
suponiendo que los ángulos en la forma polar de cualquiera de las dos soluciones de la ecuación
\[ 
 x^{5} - 1 + \mathrm{i}\sqrt{3} = 0
\]
difieren por un múltiplo entero de \(\varphi \).}
{\(\varphi = \frac{2} 
{5}\pi \)}{\(\varphi = \frac{3} 
{5}\pi \)}{\(\varphi = \frac{4} 
{5}\pi \)}{\(\varphi =\pi \)}{}{}}
\End

\Data\ID{9000035810}
\Updated{2023-02-21 11:02:56}
\Level{C}
\SubArea{Powers and Roots of Complex Numbers}
\LANGen{
\Question{
Given the complex number \(z = -2 + 2\mathrm{i}\),
find all the roots of \(\root{3}\of{z}\).}
{\(\begin{aligned}[t] &w_{0} = \root{6}\of{8}\left (\cos \frac{\pi }
{4} + \mathrm{i}\sin \frac{\pi }
{4}\right ) &
\\&w_{1} = \root{6}\of{8}\left (\cos \frac{11\pi }
{12} + \mathrm{i}\sin \frac{11\pi }
{12}\right )
\\&w_{2} = \root{6}\of{8}\left (\cos \frac{19\pi }
{12} + \mathrm{i}\sin \frac{19\pi }
{12}\right )
\\ \end{aligned}\)}{\(\begin{aligned}[t] &w_{0} = 2\left (\cos \frac{\pi }
{4} + \mathrm{i}\sin \frac{\pi }
{4}\right ) &
\\&w_{1} = 2\left (\cos \frac{11\pi }
{12} + \mathrm{i}\sin \frac{11\pi }
{12}\right )
\\&w_{2} = 2\left (\cos \frac{19\pi }
{12} + \mathrm{i}\sin \frac{19\pi }
{12}\right )
\\ \end{aligned}\)}{\(\root{3}\of{-2} + \root{3}\of{2}\)}{\(\begin{aligned}[t] &w_{0} = 2\left (\cos \frac{\pi }
{3} + \mathrm{i}\sin \frac{\pi }
{3}\right )&
\\&w_{1} = 2\left (\cos \pi +\mathrm{i}\sin \pi \right )
\\&w_{2} = 2\left (\cos \frac{5\pi }
{3} + \mathrm{i}\sin \frac{5\pi }
{3}\right )
\\ \end{aligned}\)}{}{}}
\LANGpl{
\Question{
Podano liczbę zespoloną \(z = -2 + 2\mathrm{i}\),
wskaż wszystkie pierwiastki \(\root{3}\of{z}\).}
{\(\begin{aligned}[t] &w_{0} = \root{6}\of{8}\left (\cos \frac{\pi }
{4} + \mathrm{i}\sin \frac{\pi }
{4}\right ) &
\\&w_{1} = \root{6}\of{8}\left (\cos \frac{11\pi }
{12} + \mathrm{i}\sin \frac{11\pi }
{12}\right )
\\&w_{2} = \root{6}\of{8}\left (\cos \frac{19\pi }
{12} + \mathrm{i}\sin \frac{19\pi }
{12}\right )
\\ \end{aligned}\)}{\(\begin{aligned}[t] &w_{0} = 2\left (\cos \frac{\pi }
{4} + \mathrm{i}\sin \frac{\pi }
{4}\right ) &
\\&w_{1} = 2\left (\cos \frac{11\pi }
{12} + \mathrm{i}\sin \frac{11\pi }
{12}\right )
\\&w_{2} = 2\left (\cos \frac{19\pi }
{12} + \mathrm{i}\sin \frac{19\pi }
{12}\right )
\\ \end{aligned}\)}{\(\root{3}\of{-2} + \root{3}\of{2}\)}{\(\begin{aligned}[t] &w_{0} = 2\left (\cos \frac{\pi }
{3} + \mathrm{i}\sin \frac{\pi }
{3}\right )&
\\&w_{1} = 2\left (\cos \pi +\mathrm{i}\sin \pi \right )
\\&w_{2} = 2\left (\cos \frac{5\pi }
{3} + \mathrm{i}\sin \frac{5\pi }
{3}\right )
\\ \end{aligned}\)}{}{}}
\LANGcs{
\Question{
Je dáno komplexní číslo \(z = -2 + 2\mathrm{i}\).
Všechny navzájem různé hodnoty \(\root{3}\of{z}\) 
jsou:}
{\(\begin{aligned}[t] &w_{0} = \root{6}\of{8}\left (\cos \frac{\pi }
{4} + \mathrm{i}\sin \frac{\pi }
{4}\right ) &
\\&w_{1} = \root{6}\of{8}\left (\cos \frac{11\pi }
{12} + \mathrm{i}\sin \frac{11\pi }
{12}\right )
\\&w_{2} = \root{6}\of{8}\left (\cos \frac{19\pi }
{12} + \mathrm{i}\sin \frac{19\pi }
{12}\right )
\\ \end{aligned}\)}{\(\begin{aligned}[t] &w_{0} = 2\left (\cos \frac{\pi }
{4} + \mathrm{i}\sin \frac{\pi }
{4}\right ) &
\\&w_{1} = 2\left (\cos \frac{11\pi }
{12} + \mathrm{i}\sin \frac{11\pi }
{12}\right )
\\&w_{2} = 2\left (\cos \frac{19\pi }
{12} + \mathrm{i}\sin \frac{19\pi }
{12}\right )
\\ \end{aligned}\)}{\(\root{3}\of{-2} + \root{3}\of{2}\)}{\(\begin{aligned}[t] &w_{0} = 2\left (\cos \frac{\pi }
{3} + \mathrm{i}\sin \frac{\pi }
{3}\right )&
\\&w_{1} = 2\left (\cos \pi +\mathrm{i}\sin \pi \right )
\\&w_{2} = 2\left (\cos \frac{5\pi }
{3} + \mathrm{i}\sin \frac{5\pi }
{3}\right )
\\ \end{aligned}\)}{}{}}
\LANGsk{
\Question{
Je dané komplexné číslo \(z = -2 + 2\mathrm{i}\).
Všetky navzájom rôzne hodnoty \(\root{3}\of{z}\) sú:}
{\(\begin{aligned}[t] &w_{0} = \root{6}\of{8}\left (\cos \frac{\pi }
{4} + \mathrm{i}\sin \frac{\pi }
{4}\right ) &
\\&w_{1} = \root{6}\of{8}\left (\cos \frac{11\pi }
{12} + \mathrm{i}\sin \frac{11\pi }
{12}\right )
\\&w_{2} = \root{6}\of{8}\left (\cos \frac{19\pi }
{12} + \mathrm{i}\sin \frac{19\pi }
{12}\right )
\\ \end{aligned}\)}{\(\begin{aligned}[t] &w_{0} = 2\left (\cos \frac{\pi }
{4} + \mathrm{i}\sin \frac{\pi }
{4}\right ) &
\\&w_{1} = 2\left (\cos \frac{11\pi }
{12} + \mathrm{i}\sin \frac{11\pi }
{12}\right )
\\&w_{2} = 2\left (\cos \frac{19\pi }
{12} + \mathrm{i}\sin \frac{19\pi }
{12}\right )
\\ \end{aligned}\)}{\(\root{3}\of{-2} + \root{3}\of{2}\)}{\(\begin{aligned}[t] &w_{0} = 2\left (\cos \frac{\pi }
{3} + \mathrm{i}\sin \frac{\pi }
{3}\right )&
\\&w_{1} = 2\left (\cos \pi +\mathrm{i}\sin \pi \right )
\\&w_{2} = 2\left (\cos \frac{5\pi }
{3} + \mathrm{i}\sin \frac{5\pi }
{3}\right )
\\ \end{aligned}\)}{}{}}
\LANGes{
\Question{
Dados los números complejos \(z = -2 + 2\mathrm{i}\),
determina todas las raíces de \(\root{3}\of{z}\).}
{\(\begin{aligned}[t] &w_{0} = \root{6}\of{8}\left (\cos \frac{\pi }
{4} + \mathrm{i}\sin \frac{\pi }
{4}\right ) &
\\&w_{1} = \root{6}\of{8}\left (\cos \frac{11\pi }
{12} + \mathrm{i}\sin \frac{11\pi }
{12}\right )
\\&w_{2} = \root{6}\of{8}\left (\cos \frac{19\pi }
{12} + \mathrm{i}\sin \frac{19\pi }
{12}\right )
\\ \end{aligned}\)}{\(\begin{aligned}[t] &w_{0} = 2\left (\cos \frac{\pi }
{4} + \mathrm{i}\sin \frac{\pi }
{4}\right ) &
\\&w_{1} = 2\left (\cos \frac{11\pi }
{12} + \mathrm{i}\sin \frac{11\pi }
{12}\right )
\\&w_{2} = 2\left (\cos \frac{19\pi }
{12} + \mathrm{i}\sin \frac{19\pi }
{12}\right )
\\ \end{aligned}\)}{\(\root{3}\of{-2} + \root{3}\of{2}\)}{\(\begin{aligned}[t] &w_{0} = 2\left (\cos \frac{\pi }
{3} + \mathrm{i}\sin \frac{\pi }
{3}\right )&
\\&w_{1} = 2\left (\cos \pi +\mathrm{i}\sin \pi \right )
\\&w_{2} = 2\left (\cos \frac{5\pi }
{3} + \mathrm{i}\sin \frac{5\pi }
{3}\right )
\\ \end{aligned}\)}{}{}}
\End
